\documentclass[12pt]{report}

% ---------------------------------------------------------------
% 字型與中文排版
% ---------------------------------------------------------------
% 作者在 macOS 上用 Times New Roman／Songti TC／PingFang TC；
% 其他系統（例如 Linux）找不到這些字型時，改用 TeX Gyre Termes 與
% Noto Serif／Sans CJK TC，同一份原始碼兩邊都能編。
\usepackage{fontspec}
\usepackage{xeCJK}
\IfFontExistsTF{Times New Roman}
    {\setmainfont{Times New Roman}}
    {\setmainfont{TeX Gyre Termes}}
\IfFontExistsTF{Songti TC}
    {\setCJKmainfont{Songti TC}}
    {\setCJKmainfont{Noto Serif CJK TC}}
\IfFontExistsTF{PingFang TC}
    {\setCJKsansfont{PingFang TC}}
    {\setCJKsansfont{Noto Sans CJK TC}}
% 等寬字用 Latin Modern Mono。它不是系統字型，只在 TeX Live 樹裡；
% 用「檔名＋副檔名」讓 XeTeX 透過 kpathsea 找檔案，不寫死 TeX Live
% 的年份路徑。
\setmonofont{lmmono10-regular}[
    Extension = .otf,
    AutoFakeBold = 2.5,
]

% ---------------------------------------------------------------
% 本書排版套件（packages/ 目錄）
% ---------------------------------------------------------------
\usepackage{sleek}
\usepackage{sleek-code}
\usepackage{sleek-diagrams}
\usepackage{sleek-recap}
\usepackage{sleek-aki}
% 注意：sleek-theorems.sty 不載入——本書沒有用到定理類環境，而它依賴的 thmtools
% 和 2026-11 以後的 LaTeX 核心衝突（\newtheorem{lem}[thm] 會報「\c@lem already
% defined」），載入的話整本書編不過。
% 注意：sleek-listings.sty 不載入——它是舊版 lstlistings/inconsolata 排版套件，
% 本書實際用的是 sleek-code.sty 的 minted 版 code/example 環境，兩者同時載入會讓
% minted 抓不到 \setmonofont 設定的等寬字（回退成 Latin Modern Roman），
% 且會讓程式碼裡的直引號被轉成全形彎引號。

% 章節標題中文化：「第 N 章」
\usepackage{titlesec}
\titleformat{\chapter}{\normalfont\huge\bfseries}{第~\thechapter~章}{1em}{}
\titlespacing*{\chapter}{0pt}{-20pt}{20pt}
% 頁首與目錄標題中文化（預設是 CHAPTER 1. 和 CONTENTS）
\renewcommand{\contentsname}{目錄}
\renewcommand{\chaptermark}[1]{\markboth{第~\thechapter~章\quad #1}{}}

% 行內程式碼（\texttt）不能斷行，段落裡出現較長的識別字時容易凸出右邊界；
% 允許 TeX 多拉開一點字距來排版。
\setlength{\emergencystretch}{3em}

% PDF 的書名資訊另外寫，不從 \title 抓（\title 裡有換行和系列名稱）。
\hypersetup{
    pdftitle={C++ 演算法解題 1：APCS 實作題核心},
    pdfauthor={林子祐},
    pdfsubject={註解 Annotated 演算法系列 第 1 本},
}

\title{C++ 演算法解題 1\\[0.3em]{\LARGE APCS 實作題核心}\\[1.2em]{\large 註解 Annotated 演算法系列}}
\author{林子祐}
\date{\today}

\begin{document}

\maketitle
\romantableofcontents

% ---------------------------------------------------------------
% 各章內容放在 "chapters/<單元>/chapter.tex"。
% 尚未動筆的章先註解掉，整本書才編得過。
% ---------------------------------------------------------------

\chapter{解題前的準備}
\label{ch:prep}

一間小診所早上收到一疊掛號單，順序是亂的。 護理師有兩種叫號的做法。 第一種：
每次要叫下一位，就把還沒叫的單子從頭翻一遍，找出號碼最小的那張。 第二種：開診
前先把每張單子插進一個照號碼分好格的架子，之後每次叫號只要從架子上依序拿下一
張。 插一張單子比翻一張慢，大約要花翻十張的時間。 病人只有五位時，哪一種比較省
事？ 如果有一百位、一千位呢？ 往下讀之前，先自己估估看，人數到多少時第二種開始
划算。

這個問題沒有用到任何演算法，只用到一件事：數清楚一個做法要做幾次動作，再看這個
次數隨人數怎麼長大。 寫解題程式也是一樣。 動手之前先算程式大約要做多少次運算，
就知道它在題目的時間限制內跑不跑得完；這比寫完、送出、等結果，再回頭改快得多。

這一章不教特定的演算法，教的是後面每一章都會用到的四件事：

\begin{itemize}
\item 數清楚一段程式執行幾次，以及這個次數隨資料量怎麼長大（第~\ref{sec:count-loops}、\ref{sec:growth}~節）。
\item 從題目給的資料範圍，判斷需要多快的做法（第~\ref{sec:growth}~節）。
\item 避開兩種和演算法無關、卻很常讓答案錯掉的問題：整數溢位和輸入輸出（第~\ref{sec:overflow}、\ref{sec:io}~節）。
\item 自己產生測試資料，用一個慢但簡單的程式來檢查快的程式（第~\ref{sec:testing}~節）。
\end{itemize}

這一章會用到《C++ 程式設計：給學過 C 的人》（以下簡稱銜接書）裡的幾個地方：
解題程式的骨架、從資料範圍估計時間的表、用檔案測試、\texttt{long long} 的溢位在
第~10~章 10.1 節；輸入加速的兩行在第~1~章；\texttt{getline} 在第~3~章；
\texttt{vector}、\texttt{pair}、\texttt{sort}、\texttt{set}／\texttt{map}、
\texttt{lower\_bound} 分別在第~3、8、9~章。 本書從第~2~章開始會一直用到這些容器
和函式，章末的實作練習也放了一組複習題。 沒讀過銜接書的讀者，只要會寫一支從標準
輸入讀資料、把答案印到標準輸出的 C++ 程式，就能讀這一章。

\section{數一數程式做了幾次}
\label{sec:count-loops}

先把章首的第一種叫號方法寫成程式。 \texttt{num} 存著 $n$ 張單子的號碼，每次叫號，
都在還沒叫過的單子裡找最小的，找到後換到前面，表示已經叫過：

\begin{code}[label=code:call-loop]{cpp}{第一種叫號方法：每次翻遍還沒叫的單子}
long long flips = 0;                      // 比較了幾次號碼
for (int call = 0; call < n; call++) {    // 第 call 次叫號
    int best = call;
    for (int i = call + 1; i < n; i++) {  // 只翻還沒叫過的
        flips++;
        if (num[i] < num[best]) {
            best = i;
        }
    }
    swap(num[call], num[best]);           // 叫過的換到前面
}
\end{code}

內層迴圈的範圍從 \texttt{call + 1} 開始，跟著外層變。 這種迴圈不能用「外層 $n$ 次
乘內層 $n$ 次」來算，要一輪一輪列出來。 $n = 5$ 時：

\begin{center}
\begin{tabular}{cccc}
\toprule
\texttt{call} & 內層的 \texttt{i} & 這一輪的 \texttt{flips++} & 累計 \\
\midrule
0 & 1, 2, 3, 4 & 4 & 4 \\
1 & 2, 3, 4 & 3 & 7 \\
2 & 3, 4 & 2 & 9 \\
3 & 4 & 1 & 10 \\
4 & （不執行） & 0 & 10 \\
\bottomrule
\end{tabular}
\end{center}

每一輪少一次，加起來是 $4 + 3 + 2 + 1 + 0 = 10$。 一般的 $n$ 是
\[
(n - 1) + (n - 2) + \cdots + 1 + 0 = \frac{n(n - 1)}{2}.
\]
這裡數的是 \texttt{flips++} 這一行，也就是比較號碼的次數。 迴圈本身的 \texttt{i++}、
\texttt{i < n} 也要花時間，但它們的次數跟 \texttt{flips++} 差不多，數最內層那一行
就足夠代表整段程式的工作量。

\subsection{三種常見的迴圈形狀}

大部分迴圈的次數，可以拆成下面三種形狀來算。

\begin{itemize}
\item \textbf{固定範圍。} \texttt{for (i = 0; i < n; i++)} 跑 $n$ 次；兩層都是固定範
圍時相乘，跑 $n \times m$ 次。
\item \textbf{內層範圍跟著外層變。} 像上面的叫號程式，要一輪一輪加起來。 常見的結果
是 $1 + 2 + \cdots + n = n(n + 1)/2$，大約是 $n^2/2$。
\item \textbf{每次乘或除一個數。} \texttt{for (i = 1; i <= n; i *= 2)} 的 \texttt{i}
依序是 $1, 2, 4, 8, \ldots$，到超過 $n$ 為止。 $n = 20$ 時是 $1, 2, 4, 8, 16$，跑
$5$ 次；$n = 1000$ 時只跑 $10$ 次。 一般來說是 $\lfloor \log_2 n \rfloor + 1$
次，也就是 $n$ 寫成二進位的位數。 每次除以 $3$ 的迴圈也一樣，次數大約是
$\log_3 n$。
\end{itemize}

第三種形狀的次數長得非常慢：$n$ 從一千變成一百萬，次數只從 $10$ 變成 $20$。 所以
看到「兩層迴圈」先別急著說它很慢，要看清楚每一層的變數是怎麼變的。 下面的例題就
是三段看起來都是兩層、實際次數卻差很多的程式。

\subsection{例題：迴圈到底跑幾次}
\label{sec:ex-loop-count}

\textbf{題意}（\akiref{how-many-times-does-the-loop-run}{迴圈到底跑幾次}）

\begin{itemize}
\item 有三段程式片段，變數都是 \texttt{long long}，\texttt{cnt} 從 $0$ 開始（見下面的
片語）。
\item 有 $q$ 個問題，每個問題給片段編號 $t$ 和正整數 $n$，問那一段跑完後
\texttt{cnt} 是多少。
\item $q \le 10^5$，$n \le 10^{18}$。 題目保證片段裡的變數不會溢位。
\end{itemize}

\begin{code}[label=code:three-pieces]{cpp}{例題的三段片段}
// 片段 1
for (i = 1; i <= n; i++) {
    j = i;
    while (j <= n) {
        cnt++;
        j *= 2;
    }
}
// 片段 2
for (i = 1; i <= n; i *= 2)
    for (j = 1; j <= i; j++)
        cnt++;
// 片段 3（i /= 3 是整數除法）
for (i = n; i > 0; i /= 3)
    for (j = 0; j <= i; j++)
        cnt++;
\end{code}

例如 $n = 9$ 時，片段 3 的外層 \texttt{i} 依序是 $9, 3, 1$，下一個是 $0$ 就停了。
內層的 \texttt{j} 從 $0$ 跑到 \texttt{i}，三輪分別跑 $10$、$4$、$2$ 次，答案是
$16$。

\textbf{想法}

最直接的做法，是把三段片段原封不動放進程式裡跑：寫一個函式
\texttt{run(t, n)}，裡面用 \texttt{if} 分成三段，每一段就是上面的片語照抄，最後回傳
\texttt{cnt}；\texttt{main} 讀每個問題、印出 \texttt{run(t, n)}。 這就是這題的
\textbf{暴力解}（brute force）：照題目的字面意思做，不動任何腦筋，所以幾乎不會錯。

這個程式要做的事，跟答案一樣多：答案是多少，\texttt{cnt++} 就執行幾次。 片段 1
在 $n = 10^9$ 時的答案大約是 $2 \times 10^9$，在本書測試的環境下這支程式跑一筆要
$1.1$ 秒左右。 $n$ 到 $10^{18}$ 時是它的十億倍，而且最多有 $10^5$ 筆，完全不可能。
所以要找一個不必真的跑完、就能算出次數的方法。

\textbf{片段 2 和片段 3：只跑外層。} 片段 2 的外層 \texttt{i} 每次乘 $2$，片段 3 的
外層每次除以 $3$，外層最多跑六十幾輪。 內層的次數直接寫得出來：片段 2 是 \texttt{i}
次，片段 3 是 \texttt{i + 1} 次（\texttt{j} 從 $0$ 開始）。 所以只要跑外層，每一輪
把內層的次數加上去。

\textbf{片段 1：換一個方向數。} 片段 1 的外層是 \texttt{i++}，要跑 $n$ 輪，只跑外層
也來不及。 先看 $n = 6$ 時每一輪的 \texttt{j}：

\begin{center}
\begin{tikzpicture}[x=1cm, y=1cm]
\foreach \i in {1,...,6} {\node[idxlbl] at (\i,0.75) {\texttt{i} = \i};}
\node[left, font=\small] at (0.4,0) {乘 $0$ 次};
\node[left, font=\small] at (0.4,-1) {乘 $1$ 次};
\node[left, font=\small] at (0.4,-2) {乘 $2$ 次};
\foreach \i in {1,...,6} {\node[hlcell] at (\i,0) {\i};}
\foreach \i/\j in {1/2,2/4,3/6} {\node[hlcell] at (\i,-1) {\j};}
\foreach \i in {4,5,6} {\node[cell, draw=black!20] at (\i,-1) {};}
\node[hlcell] at (1,-2) {4};
\foreach \i in {2,...,6} {\node[cell, draw=black!20] at (\i,-2) {};}
\node[font=\small] at (7.4,0) {$6$ 格};
\node[font=\small] at (7.4,-1) {$3$ 格};
\node[font=\small] at (7.4,-2) {$1$ 格};
\end{tikzpicture}
\diagcaption{片段 1 在 $n = 6$ 時，每一格是 \texttt{cnt++} 執行時 \texttt{j} 的值。
直著看是一輪外層，橫著看是「乘了幾次 $2$」。}
\label{fig:piece1-grid}
\end{center}

題目問的是格子總數。 直著一行一行數（照程式的順序）是 $3 + 2 + 2 + 1 + 1 + 1 = 10$，
要數 $n$ 行。 改成橫著一列一列數：第 $k$ 列是 \texttt{j} 已經乘了 $k$ 次 $2$ 的那
些格子，\texttt{j} $= i \times 2^k$。 這一格存在，剛好是 $i \times 2^k \le n$，也就
是 $i \le n / 2^k$，所以第 $k$ 列有 $\lfloor n / 2^k \rfloor$ 格。 總數是
\[
n + \left\lfloor \frac{n}{2} \right\rfloor + \left\lfloor \frac{n}{4} \right\rfloor + \cdots,
\]
加到 $\lfloor n / 2^k \rfloor = 0$ 為止。 $n = 6$ 時是 $6 + 3 + 1 = 10$，和直著數一
樣。 這個和最多六十幾項，用一個 \texttt{m /= 2} 的迴圈就能算出來。 而且
$\lfloor n / 2^k \rfloor$ 也等於 $\lfloor \lfloor n / 2^{k-1} \rfloor / 2 \rfloor$，
所以 \texttt{m} 從 $n$ 開始一直除以 $2$ 就好，不用算 $2^k$。

順便看出一件事：這個和小於 $n + n/2 + n/4 + \cdots = 2n$。 片段 1 雖然有兩層迴圈，
\texttt{cnt++} 只執行不到 $2n$ 次，跟 $n^2$ 差得很遠。

\textbf{程式}

\begin{example}[label=ex:loop-count]{cpp}{(loop\_count.cpp) 迴圈到底跑幾次}
#include <iostream>
using namespace std;

// 片段 1：第 k 列的 j 是 i 乘上 2 的 k 次方，
// 這一列有 n / 2^k 格（整數除法）
long long piece1(long long n) {
    long long total = 0;
    for (long long m = n; m > 0; m /= 2) {
        total += m;
    }
    return total;
}

// 片段 2：外層 i = 1, 2, 4, ...，內層各跑 i 次
long long piece2(long long n) {
    long long total = 0;
    for (long long i = 1; i <= n; i *= 2) {
        total += i;
    }
    return total;
}

// 片段 3：外層 i = n, n/3, n/9, ...，內層 j 從 0 到 i，跑 i + 1 次
long long piece3(long long n) {
    long long total = 0;
    for (long long i = n; i > 0; i /= 3) {
        total += i + 1;
    }
    return total;
}

int main() {
    ios::sync_with_stdio(false);
    cin.tie(nullptr);
    int q;
    cin >> q;
    while (q--) {
        int t;
        long long n;
        cin >> t >> n;
        long long ans;
        if (t == 1) {
            ans = piece1(n);
        } else if (t == 2) {
            ans = piece2(n);
        } else {
            ans = piece3(n);
        }
        cout << ans << '\n';
    }
    return 0;
}
\end{example}

\iosample{
5 \\
3 1 \\
2 7 \\
2 8 \\
3 9 \\
1 10 \\
}{
2 \\
7 \\
15 \\
16 \\
18 \\
}

\begin{itemize}
\item 把同一組輸入拿給暴力解跑，輸出也相同。 兩支程式做法完全不同，答案卻一樣，這本
身就是一種檢查；第~\ref{sec:testing}~節會把這個想法變成固定的流程。
\item 三個函式的迴圈都和原本的片段很像，差別只在內層被換成一個加法。 這樣寫比較容
易和原本的片段逐行對照。
\item 答案最大大約是 $2 \times 10^{18}$，在 \texttt{long long} 的範圍內（上限大約
$9.2 \times 10^{18}$）；\texttt{int} 只到大約 $2.1 \times 10^9$，一定不夠。
\end{itemize}

\textbf{複雜度}

每個問題的三個函式各最多跑 $\log_2 10^{18} \approx 60$ 輪，$10^5$ 個問題總共幾百
萬次加法。 最大的測試資料在本書測試的環境下約 $0.01$ 秒。

\subsection{常見錯誤}

\begin{itemize}
\item \textbf{看到兩層就套公式}：把片段 1 當成「外層 $n$ 次、內層 $\log_2 n$ 次」，寫成
$n \times (\lfloor \log_2 n \rfloor + 1)$。 範例 \texttt{1 10} 會印出 \texttt{40}，正確
是 \texttt{18}。 內層的 \texttt{j} 從 \texttt{i} 開始而不是從 $1$ 開始，\texttt{i}
越大，內層越短。
\item \textbf{漏掉從 $0$ 開始的那一次}：片段 3 寫成 \texttt{total += i}，範例
\texttt{3 1} 會印出 \texttt{1}，正確是 \texttt{2}。 \texttt{j} 從 $0$ 跑到
\texttt{i} 是 $i + 1$ 個數。
\item \textbf{用浮點數算次數}：片段 2 的外層跑 $\lfloor \log_2 n \rfloor + 1$ 輪，有人
用 \texttt{log2((double)n)} 算出輪數，再套 $2^{m+1} - 1$。 輸入 \texttt{2
576460752303423487}（$n = 2^{59} - 1$）時，正確答案是
\texttt{576460752303423487}，這樣寫在本書測試的環境下印出
\texttt{1152921504606846975}，多算了一輪。 原因在第~\ref{sec:overflow}~節：
\texttt{double} 存不下這麼大的整數，$2^{59} - 1$ 被捨入成 $2^{59}$。
\item \textbf{直接跑片段 1}：只把片段 2、3 改成只跑外層，片段 1 還是照原樣跑。 範例都
會對，但 $n = 10^{18}$ 的資料永遠跑不完。
\end{itemize}

\section{從次數到成長速度}
\label{sec:growth}

回到章首的診所。 第一種方法翻單子 $n(n - 1)/2$ 次。 第二種方法每張單子插一次、拿
一次，插一張算翻十張的時間，總共相當於翻 $10n + n = 11n$ 張。 兩種方法在幾個人數下
的比較：

\begin{center}
\begin{tabular}{rrrl}
\toprule
人數 $n$ & 第一種 $n(n-1)/2$ & 第二種 $11n$ & 比較省事的 \\
\midrule
$5$ & $10$ & $55$ & 第一種 \\
$20$ & $190$ & $220$ & 第一種 \\
$23$ & $253$ & $253$ & 一樣 \\
$24$ & $276$ & $264$ & 第二種 \\
$100$ & $4950$ & $1100$ & 第二種 \\
$1000$ & $499500$ & $11000$ & 第二種 \\
$10^5$ & 約 $5 \times 10^9$ & $1.1 \times 10^6$ & 第二種 \\
\bottomrule
\end{tabular}
\end{center}

人少的時候，第一種方法比較省事，因為它不用花插單子的那段固定成本。 兩邊在 $23$
人時剛好打平，從 $24$ 人起第二種方法一直比較快，而且差距越拉越大：一千人時差了
四十多倍，十萬人時差了四千多倍。

\begin{center}
\begin{tikzpicture}[x=0.25cm, y=0.0075cm]
\draw[->, black!60] (0,0) -- (42,0) node[right, font=\small] {$n$};
\draw[->, black!60] (0,0) -- (0,800) node[above, font=\small] {翻單子的次數};
\foreach \t in {10,20,30,40} {\draw[black!60] (\t,0) -- (\t,-15) node[below, font=\scriptsize] {\t};}
\foreach \t in {200,400,600} {\draw[black!60] (0,\t) -- (-0.4,\t) node[left, font=\scriptsize] {\t};}
\draw[thick, optStroke, domain=0:40, samples=60] plot (\x, {\x*(\x-1)/2});
\draw[thick, coreStroke, domain=0:40] plot (\x, {11*\x});
\draw[dashed, black!50] (23,0) -- (23,253);
\fill[black] (23,253) circle (2.5pt);
\node[font=\small, optStroke] at (29,620) {第一種};
\node[font=\small, coreStroke] at (39,350) {第二種};
\node[font=\scriptsize, above left] at (23,253) {$n = 23$};
\end{tikzpicture}
\diagcaption{兩種叫號方法翻單子的次數。 交叉點左邊第一種較少，右邊第二種較少。}
\label{fig:crossover}
\end{center}

如果插單子更慢，例如一張算翻一百張，交叉點會往右移到兩百人左右，但右邊最後還是第
二種贏。 理由是兩個次數長大的方式不一樣：第二種和 $n$ 成正比，人數變十倍，次數也
變十倍；第一種大約是 $n^2/2$，人數變十倍，次數變一百倍。 只要 $n$ 夠大，前面的係數
$10$ 或 $100$ 都追不上這個差距。

\subsection{只看最大的那一項}

把上面的觀察寫成一般的規則：$n$ 很大時，決定次數的是長得最快的那一項，係數和比較
小的項都不重要。 $n(n - 1)/2 = n^2/2 - n/2$，最大的項是 $n^2/2$，去掉係數就是
$n^2$，記作 $O(n^2)$；$11n$ 記作 $O(n)$。 這種寫法叫\textbf{時間複雜度}（time
complexity），讀作「big-O」。 白話的意思是：$n$ 夠大以後，次數不超過括號裡那個式
子的某個固定倍數。

\begin{itemize}
\item 兩段程式先後執行，複雜度取比較大的那段。 先排序 $O(n \log n)$、再掃一遍
$O(n)$，合起來是 $O(n \log n)$。
\item 一段 $O(f)$ 的程式被執行 $g$ 次，合起來是 $O(f \times g)$。
\item 次數會隨輸入的內容變（例如找到就 \texttt{break}）時，看最壞的情況。
\item $\log$ 的底數只差一個固定倍數（$\log_3 n$ 大約是 $\log_2 n$ 的 $0.63$ 倍），所以
$O(\log n)$ 不寫底數。
\end{itemize}

要分清楚兩種問題。 「這段程式的 \texttt{cnt++} 執行幾次」要的是精確的次數，例如
$n(n - 1)/2$，上一節的例題和本章練習的追蹤題都是這種。 「這個做法來不來得及」只要
量級，例如 $O(n^2)$，接下來估時間用的都是這種。 兩者都要會：精確的次數算錯了，量
級常常也跟著錯。

\subsection{實際量一量}

量級只說次數怎麼長大，沒說一秒能做幾次。 拿一個 $O(n^2)$ 的雙層迴圈來量：
對每一對 $i < j$ 檢查 \texttt{a[i] + a[j]} 是否等於某個數，共比較 $n(n - 1)/2$ 次。
在本書測試的環境下，三種大小各跑一次：

\begin{center}
\begin{tabular}{rrr}
\toprule
$n$ & 比較次數 $n(n-1)/2$ & 時間 \\
\midrule
$10^4$ & 約 $5 \times 10^7$ & $0.02$ 秒 \\
$3 \times 10^4$ & 約 $4.5 \times 10^8$ & $0.16$ 秒 \\
$10^5$ & 約 $5 \times 10^9$ & $1.8$ 秒 \\
\bottomrule
\end{tabular}
\end{center}

$n$ 變三倍，時間大約變九倍；$n$ 變十倍，時間大約變一百倍，和 $O(n^2)$ 的預測一致。
照這個比例，$n = 2 \times 10^5$ 會要七秒左右，大部分題目的時間限制是一兩秒，來不及。

一秒能做幾次，要看迴圈裡做的是什麼。 同一台機器上，另外量了兩種比較慢的操作：

\begin{center}
\begin{tabular}{lr}
\toprule
迴圈裡做的事 & 每秒大約幾次 \\
\midrule
加法和比較，依序讀一個陣列（上面的迴圈） & $3 \times 10^9$ \\
在一個 64 MB 的陣列裡，照亂數跳著讀 & $10^8$ \\
在約 60 萬個鍵的 \texttt{map} 裡插入或查詢 & $2 \times 10^6$ \\
\bottomrule
\end{tabular}
\end{center}

最快和最慢差了上千倍。 依序讀陣列時，電腦可以一次搬一大塊資料、同時做好幾個比較；
跳著讀時，每次都要等資料從記憶體送過來；\texttt{map} 的每次操作本身就要跳十幾二十
次。 銜接書 10.1 節的表用「一秒大約 $10^8$ 次簡單運算」來估，這個數字對單純的迴圈
很保守，但加上這些慢的操作之後剛好合用。 本書也用它：算出來在 $10^8$ 以下，一定
來得及；到 $10^9$ 附近，要看迴圈裡做的事有多單純；超過 $10^{10}$，換做法。

\subsection{從資料範圍反推需要的複雜度}
\label{sec:from-limits}

題目的限制欄寫著 $n \le 2 \times 10^5$，這個數字除了告訴你陣列要開多大，也在告訴你
需要多快的做法。 銜接書 10.1 節有一張「$n$ 的上限 $\to$ 可以接受的複雜度」的表，
本書不重印，只把用法整理成步驟：

\begin{enumerate}
\item 找出所有會讓工作量變大的數：資料筆數 $n$、詢問數 $q$、測試資料的組數、數值的
範圍。
\item 寫出你的做法的複雜度，把\textbf{每一個}數的上限都代進去。 詢問 $q$ 次、每次
$O(n)$，就是 $n \times q$，不是 $n$。
\item 和 $10^8$ 比。
\end{enumerate}

第二步最容易漏。 下面是幾個例子：

\begin{center}
\begin{tabular}{p{0.32\linewidth}p{0.22\linewidth}p{0.36\linewidth}}
\toprule
限制 & 想用的做法 & 估計 \\
\midrule
$n \le 2 \times 10^5$ & 排序後掃一遍 & $n \log_2 n \approx 3.6 \times 10^6$，可以 \\
$n, q \le 2 \times 10^5$ & 每個詢問掃一遍陣列 & $n \times q = 4 \times 10^{10}$，不行 \\
$n \le 5000$ & 檢查每一對 & $n^2/2 \approx 1.3 \times 10^7$，可以 \\
組數 $\le 10^5$，所有組的 $n$ 總和 $\le 10^6$ & 每組 $O(n)$ & 加起來 $O(10^6)$，可以 \\
同上，但每組都清空一個大小 $10^6$ 的陣列 & & $10^5 \times 10^6 = 10^{11}$，不行 \\
$n \le 20$ & 枚舉所有子集合 & $2^{20} \approx 10^6$，可以 \\
\bottomrule
\end{tabular}
\end{center}

倒數第二列是第~\ref{sec:ex-vote}~小節例題的重點：題目給「總和」的限制時，每一組
的工作量也要跟那一組的大小成正比，不能跟陣列宣告的大小成正比。

反過來，限制也會透露出題者想要的做法。 $n \le 20$ 幾乎是在說「可以試遍所有子集
合」；$n \le 5000$ 通常是 $O(n^2)$；$n$ 到 $10^5$ 以上就要 $O(n \log n)$ 或 $O(n)$；
$n$ 到 $10^{18}$ 時連 $O(n)$ 都不行，像上一節的例題，只能找公式或 $O(\log n)$ 的做法。

\begin{apcsnote}
實作題的題目會寫出資料範圍。 如果題目把測試資料分成幾個子題，範圍小的子題通常是給
比較直接的做法。 想不到夠快的做法時，先寫一個暴力解把小範圍的子題拿下來；這個暴力
解等一下也能拿來檢查你的快速版本（第~\ref{sec:testing}~節）。
\end{apcsnote}

\section{整數溢位與浮點數}
\label{sec:overflow}

\texttt{int} 能存的範圍大約是正負 $2.1 \times 10^9$，\texttt{long long} 大約是正負
$9.2 \times 10^{18}$。 有號整數算出超過範圍的值，叫\textbf{溢位}（overflow）。 在
C++ 裡這是未定義行為：程式不會停下來，也不會有錯誤訊息，只是印出一個錯的數字，而
且同一支程式換個編譯選項，錯的數字可能還不一樣。 銜接書 10.1 節的常見錯誤示範過
總和用 \texttt{int} 存的情形，這一節看幾種比較隱蔽的。

\subsection{先估最大值}

動手之前先算答案和中間值最大可能多大。 常見的幾種：

\begin{center}
\begin{tabular}{lll}
\toprule
算的東西 & 最大大約是 & 例子（$n = 2 \times 10^5$，值 $\le 10^9$） \\
\midrule
$n$ 個數的總和 & $n \times$ 最大值 & $2 \times 10^{14}$ \\
兩個數的乘積 & 最大值的平方 & $10^{18}$ \\
有幾對 & $n^2/2$ & $2 \times 10^{10}$ \\
$1 + 2 + \cdots + n$ & $n^2/2$ & $2 \times 10^{10}$ \\
\bottomrule
\end{tabular}
\end{center}

超過 $2 \times 10^9$ 就用 \texttt{long long}。 超過 $9 \times 10^{18}$ 時連
\texttt{long long} 都不夠，要換個算法，例如下面的「用除法比大小」。

\subsection{型別的上下限與型別別名}

範圍不用背。 \texttt{<limits>} 裡的 \texttt{numeric\_limits<T>::max()} 和
\texttt{min()} 給出型別 \texttt{T} 能存的最大值和最小值。 另外，\texttt{long long}
要一直寫很長，解題程式常在開頭寫 \texttt{using ll = long long;}，之後 \texttt{ll}
就是 \texttt{long long} 的另一個名字。 這叫\textbf{型別別名}（type alias），和 C 的
\texttt{typedef long long ll;} 意思完全一樣；\texttt{using} 的寫法是「新名字 =
舊型別」，從左讀到右，型別比較長時比較好讀。

\begin{example}[label=ex:limits]{cpp}{(limits.cpp) 型別的上下限與「無限大」}
#include <iostream>
#include <limits>
using namespace std;

using ll = long long;                   // 之後寫 ll 就是 long long

int main() {
    cout << numeric_limits<int>::max() << '\n';
    cout << numeric_limits<int>::min() << '\n';
    cout << numeric_limits<ll>::max() << '\n';
    const ll INF = 1e18;                // 當成「無限大」用
    cout << INF + INF << '\n';          // 兩個相加還在範圍內
    return 0;
}
\end{example}

\iosample{}{
2147483647 \\
-2147483648 \\
9223372036854775807 \\
2000000000000000000 \\
}

求最小值時，常先把答案設成「無限大」，再一個一個比。 直覺上會想用型別的最大值，可是之後只要拿它加任何正數就溢位了。 例如最短
路徑常寫「目前距離加上一條邊的長度」：

\begin{code}{cpp}{拿型別最大值當無限大（錯誤示範）}
ll dist = numeric_limits<ll>::max();   // 還沒走到
ll cand = dist + w;                    // w = 5
if (cand < dist) { ... }               // 以為找到更短的路
\end{code}

在本書測試的環境下，\texttt{cand} 變成 \texttt{-9223372036854775804}，比
\texttt{dist} 小，程式以為找到了一條更短的路。 所以「無限大」通常改用
\texttt{1e18} 這種常數：它比任何真正的答案都大，兩個相加也還在 \texttt{long long} 的
範圍內。 \texttt{1e18} 本身是 \texttt{double}，但 $10^{18}$ 能被 \texttt{double} 精確
表示，存進 \texttt{ll} 不會差。 \texttt{int} 的版本常用 \texttt{1e9}。

\subsection{結果存進 \texttt{long long} 也沒用}

下面這段程式的變數都宣告成 \texttt{long long} 了，可是三行裡有兩行是錯的：

\begin{code}[label=code:overflow-demo]{cpp}{三種在算式裡發生的溢位}
int w, h, n, b;
cin >> w >> h >> n >> b;               // 輸入 100000 100000 80000 36
long long area = w * h;                // 錯
long long area2 = (long long)w * h;    // 對
long long pairs = n * (n - 1) / 2;     // 錯
long long pairs2 = (long long)n * (n - 1) / 2;
long long mask = (1 << b) - 1;         // 錯
long long mask2 = (1LL << b) - 1;      // 對
\end{code}

在本書測試的環境下，用 \texttt{-O2} 編譯，印出這六個變數是：

\begin{code}{text}{上面六個變數的值}
area  = 1410065408     area2  = 10000000000
pairs = 1052476352     pairs2 = 3199960000
mask  = 15             mask2  = 68719476735
\end{code}

C++ 算 \texttt{w * h} 時只看兩個運算元的型別：兩個都是 \texttt{int}，就用
\texttt{int} 乘，乘完已經溢位，才轉成 \texttt{long long} 存起來。 等號左邊是什麼型
別，影響不到右邊怎麼算。 修正的方法是在相乘之前，先把其中一個運算元轉成
\texttt{long long}，另一個會自動跟著轉。 下面兩種寫法要分清楚：

\begin{itemize}
\item \texttt{(long long)w * h}：先轉 \texttt{w}，再乘，對。
\item \texttt{(long long)(w * h)}：先用 \texttt{int} 乘完，再轉，跟沒轉一樣。
\end{itemize}

\texttt{1 << b} 也是同一回事：\texttt{1} 是 \texttt{int}，左移 $36$ 位超過了
\texttt{int} 的 $32$ 位。 要寫 \texttt{1LL << b}，\texttt{LL} 讓這個 \texttt{1} 一開
始就是 \texttt{long long}。 這一行在加上 \texttt{-fsanitize=undefined} 編譯時印出
的是 \texttt{-1}，不是 \texttt{15}，這就是「未定義行為」的意思：結果沒有任何保證。

\subsection{用除法比大小}

有時候乘出來的值連 \texttt{long long} 都存不下，但我們只想知道它有沒有超過某個數。
例如想檢查 $n^2 \le T$，$T$ 到 $10^{18}$，$n$ 也可能到 $10^{18}$。 這時把乘法換成除
法：
\[
n \times n \le T \quad\text{等於}\quad n \le \left\lfloor \frac{T}{n} \right\rfloor
\qquad (n \ge 1).
\]
因為 $n$ 是整數，「$n \le T/n$」和「$n$ 不超過 $T/n$ 的整數部分」是同一件事。 右邊
是整數除法，結果不會比 $T$ 大，一定不會溢位。 三個數相乘也一樣，$n^3 \le T$ 可以
寫成 \texttt{n <= T / n / n}。

另一個常見的地方是兩個數的中點。 \texttt{lo} 和 \texttt{hi} 都是 \texttt{int}、可能
到 $2 \times 10^9$ 附近時，\texttt{(lo + hi) / 2} 的加法就溢位了。 改寫成
\texttt{lo + (hi - lo) / 2}，算出來的值一樣，但中間值不會超過 \texttt{hi}。 第~3~章
的二分搜尋會一直用到這個寫法。

\subsection{浮點數存不下大整數}

\texttt{double} 有 $53$ 個二進位位數存數字，能精確表示的整數只到
$2^{53} \approx 9 \times 10^{15}$。 再大的整數會被捨入成最接近的可表示值，例如
$2^{53} + 1 = 9007199254740993$ 轉成 \texttt{double} 再轉回來，會變成
$9007199254740992$。

這會讓「用 \texttt{sqrt}、\texttt{log2} 算整數答案」的程式在大數上差一：

\begingroup\sloppy
\begin{itemize}
\item $T = 10^{18} - 1$，想求最大的 $n$ 使 $n^2 \le T$，答案是 $999999999$。
\texttt{(long long)sqrt((double)T)} 卻得到 $1000000000$，因為 $10^{18} - 1$ 轉成
\texttt{double} 時被捨入成 $10^{18}$。
\item $n = 2^{59} - 1$，$\lfloor \log_2 n \rfloor = 58$。
\texttt{(long long)log2((double)n)} 得到 $59$，理由一樣。
\end{itemize}
\endgroup

浮點數還是有用處：拿它估一個大概的值，再用整數修正。 下一小節的例題會示範：
先用 \texttt{sqrt} 估，再用上一小節的除法檢查，大了就減一、小了就加一。

\subsection{例題：預算內的最大規模}
\label{sec:ex-max-size}

\textbf{題意}（\akiref{largest-size-within-the-budget}{預算內的最大規模}）

\begin{itemize}
\item 有六個函式：$f_1(n) = n$；$f_2(n) = n \times b(n)$，$b(n)$ 是 $n$ 的二進位位數，
$f_2(0) = 0$；$f_3(n) = n^2$；$f_4(n) = n^3$；$f_5(n) = 2^n$；$f_6(n) = n!$。
\item 有 $q$ 個詢問，每個給代號 $t$ 和預算 $T$，求最大的非負整數 $n$ 使
$f_t(n) \le T$。
\item $q \le 10^5$，$1 \le T \le 10^{18}$。
\end{itemize}

$T = 10^8$ 時六個答案依序是 $10^8$、$4347826$、$10^4$、$464$、$26$、$11$。 這一列
就是銜接書那張表的由來：一秒做 $10^8$ 次的話，$O(n^2)$ 的做法能處理一萬筆，$O(2^n)$
的做法只能處理二十幾筆。 $f_2$ 是 $n \log n$ 的一種精確版本。

\textbf{想法}

暴力解是 $n$ 從 $0$ 開始一個一個試，直到 $f(n) > T$。 $f_5$、$f_6$ 長得很快，試幾十
個就停了；可是 $f_1$ 在 $T = 10^{18}$ 時要試 $10^{18}$ 次，$f_3$ 要試 $10^9$ 次，再乘上
$10^5$ 個詢問，不可能。 六個函式分開處理：

\begin{itemize}
\item \textbf{$f_1$}：答案就是 $T$。
\item \textbf{$f_5$、$f_6$}：照暴力解一個一個試，最多六十步。 判斷「下一個還在預算內」
時用除法：$2^{n+1} \le T$ 寫成 \texttt{value <= T / 2}，$(n+1)! \le T$ 寫成
\texttt{value <= T / (n + 1)}。 這兩題其實直接乘也不會溢位，但要確定這件事得自己算
一遍；用除法就不用算。
\item \textbf{$f_3$、$f_4$}：用 \texttt{sqrt}、\texttt{cbrt} 估一個 $n$，再用除法修正：
$n^2 > T$ 就往下減，$(n+1)^2 \le T$ 就往上加。 浮點數只差一兩個，修正的迴圈只跑幾
次。
\item \textbf{$f_2$}：依位數分組。 有 $b$ 位的數是 $2^{b-1}$ 到 $2^b - 1$，這一組裡
$f_2(n) = n \times b$，條件 $n \times b \le T$ 等於 $n \le \lfloor T / b \rfloor$，
所以這一組最大的答案是 $\min(2^b - 1, \lfloor T / b \rfloor)$，前提是它不小於
$2^{b-1}$（否則這一組沒有任何數在預算內）。 $f_2$ 隨 $n$ 遞增，所以位數越多、還能
在預算內的組，答案越大；$b$ 從 $1$ 試到 $62$，留下最後一個合格的就好。
\end{itemize}

$f_2$ 的上限為什麼試到 $62$？ $T \le 10^{18} < 2^{60}$，答案不超過 $T$，最多 $60$
位；$62$ 只是留一點餘裕，而且 \texttt{1LL << 62} 還在 \texttt{long long} 的範圍內。

\textbf{程式}

\begin{example}[label=ex:max-size]{cpp}{(max\_size.cpp) 預算內的最大規模}
#include <cmath>
#include <iostream>
using namespace std;

// 每個函式回答：最大的 n，使 f(n) <= T。比大小一律用除法，不做會溢位的乘法。
// f(n) = n 的答案就是 T，直接寫在 main 裡。

long long nLogN(long long T) {              // f(n) = n * (n 的二進位位數)
    long long best = 0;
    for (int b = 1; b <= 62; b++) {         // 有 b 位的 n：2^(b-1) 到 2^b - 1
        long long lo = 1LL << (b - 1), hi = (1LL << b) - 1;
        long long n = min(hi, T / b);       // n * b <= T  等於  n <= T / b
        if (n >= lo) {
            best = n;
        }
    }
    return best;
}

// n^k <= T 嗎？（n >= 1）把 T 連除 k 次 n，還剩至少 1 就是
bool fits(long long n, int k, long long T) {
    for (int i = 0; i < k; i++) {
        T /= n;
    }
    return T >= 1;
}

long long root(long long T, int k) {        // f(n) = n^k，k = 2 或 3
    double guess = (k == 2) ? sqrt((double)T) : cbrt((double)T);
    long long n = (long long)guess;         // 先用浮點數估一個大概
    while (n > 0 && !fits(n, k, T)) {       // 太大，往下修
        n--;
    }
    while (fits(n + 1, k, T)) {             // 還能再大，往上修
        n++;
    }
    return n;
}

long long power2(long long T) {             // f(n) = 2^n
    long long n = 0, value = 1;             // value = 2^n
    while (value <= T / 2) {                // 2^(n+1) <= T
        value *= 2;
        n++;
    }
    return n;
}

long long factorial(long long T) {          // f(n) = n!
    long long n = 1, value = 1;             // value = n!，0! = 1! = 1
    while (value <= T / (n + 1)) {          // (n+1)! <= T
        n++;
        value *= n;
    }
    return n;
}

int main() {
    ios::sync_with_stdio(false);
    cin.tie(nullptr);
    int q;
    cin >> q;
    while (q--) {
        int t;
        long long T;
        cin >> t >> T;
        long long ans;
        if (t == 1) {
            ans = T;
        } else if (t == 2) {
            ans = nLogN(T);
        } else if (t <= 4) {
            ans = root(T, t - 1);
        } else if (t == 5) {
            ans = power2(T);
        } else {
            ans = factorial(T);
        }
        cout << ans << '\n';
    }
    return 0;
}
\end{example}

\iosample{
6 \\
1 100000000 \\
2 100000000 \\
3 100000000 \\
4 100000000 \\
5 100000000 \\
6 100000000 \\
}{
100000000 \\
4347826 \\
10000 \\
464 \\
26 \\
11 \\
}

\begin{itemize}
\item \texttt{fits} 把 $n^k \le T$ 換成「$T$ 連除 $k$ 次 $n$ 之後至少是 $1$」，和上一
小節的除法是同一個想法。 \texttt{root} 的第一個 \texttt{while} 要先檢查
\texttt{n > 0}，因為除數不能是 $0$。
\item \texttt{factorial} 從 $n = 1$ 開始，因為 $0! = 1! = 1 \le T$，答案至少是 $1$。
\texttt{power2} 從 $n = 0$ 開始：$T = 1$ 時 $2^1 = 2 > 1$，答案是 $0$。
\item \texttt{nLogN} 裡的 \texttt{1LL << b} 不能寫成 \texttt{1 << b}，$b \ge 31$ 時就
是上一小節的位移溢位。
\end{itemize}

\textbf{複雜度}

每個詢問最多六十幾步（$f_2$ 的 $62$ 組、$f_5$ 的六十次倍增），修正迴圈只跑幾次。
$10^5$ 個詢問大約 $6 \times 10^6$ 步，最大的測試資料在本書測試的環境下約 $0.01$ 秒。

\subsection{常見錯誤}

\begingroup\sloppy
\begin{itemize}
\item \textbf{只用浮點數}：$f_3$ 寫成 \texttt{(long long)sqrt((double)T)} 就回傳，輸入
\texttt{3 999999999999999999} 會印出 \texttt{1000000000}，正確是 \texttt{999999999}。
$f_4$ 用 \texttt{cbrt} 也一樣，同一個 $T$ 印出 \texttt{1000000}，正確是
\texttt{999999}。 $f_5$ 用 \texttt{log2}，輸入 \texttt{5 576460752303423487} 會印出
\texttt{59}，正確是 \texttt{58}。 範例的 $T$ 都很小，這三種寫法都會通過範例。
\item \textbf{比大小時直接乘}：\texttt{fits} 寫成把 $n$ 連乘 $k$ 次再和 $T$ 比。 這裡
$n$ 只到 $10^9$ 左右，乘起來還在範圍內；但同樣的寫法搬到 $f_2$ 的 \texttt{hi * b} 上，
$b = 62$ 時 \texttt{hi} 大約 $4.6 \times 10^{18}$，乘上 $62$ 就溢位了。 與其每次都去確認會不會溢位，不如一律用
除法。
\item \textbf{位數算少一位}：把 $b(n)$ 當成 $\lfloor \log_2 n \rfloor$，也就是從 $0$
開始數位數。 範例的第二列會印出 \texttt{4545454}，比正確的 \texttt{4347826} 大。
$8 = 1000_2$ 有 $4$ 位，不是 $3$ 位。
\item \textbf{小的邊界}：忘了 $0! = 1$，或以為 $T = 1$ 時 $f_5$ 的答案是 $1$。 這些
在範例 2 裡都有，交出去前把範例全部跑過一次。
\end{itemize}
\endgroup

\section{輸入與輸出}
\label{sec:io}

演算法再快，資料讀得太慢一樣會超過時間限制；答案算對了，多印一個空白也可能被判錯。
這一節講三件事：讀寫的速度、讀到檔案結束、多組資料。

\subsection{讀寫的速度}

\begin{bridgerecap}[輸入加速的兩行]
\texttt{main} 一開頭加上
\texttt{ios::sync\_with\_stdio(false);} 和 \texttt{cin.tie(nullptr);}（銜接書第~1~章）。
第一行關掉 \texttt{cin}／\texttt{cout} 和 C 的 \texttt{scanf}／\texttt{printf} 之間的同
步；第二行讓 \texttt{cin} 讀資料前不再先把 \texttt{cout} 的緩衝區送出去。 加了之後
不要再混用 \texttt{scanf}／\texttt{printf}，否則輸出的順序可能會亂掉。 換行用
\texttt{'\textbackslash n'}，不要用 \texttt{endl}：\texttt{endl} 除了換行，還會把緩
衝區立刻送出去，每一行都送一次就很慢。
\end{bridgerecap}

這三件事的效果，在本書測試的環境下量了一次。 程式讀 $10^6$ 個整數，每讀一個就印一
行，輸入和輸出都是檔案：

\begin{center}
\begin{tabular}{lr}
\toprule
寫法 & 時間 \\
\midrule
什麼都不加，換行用 \texttt{'\textbackslash n'} & $0.70$ 秒 \\
只加 \texttt{sync\_with\_stdio(false)} & $0.44$ 秒 \\
兩行都加，但換行用 \texttt{endl} & $0.47$ 秒 \\
兩行都加，換行用 \texttt{'\textbackslash n'} & $0.10$ 秒 \\
\bottomrule
\end{tabular}
\end{center}

少做任何一件，都會慢四倍以上。 只讀不印的時候，讀 $2 \times 10^6$ 個整數，不加速要
$0.33$ 秒，加了兩行是 $0.06$ 秒，和 \texttt{scanf} 的 $0.07$ 秒差不多。 一題的時間
限制常常只有一秒，這幾倍的差距可能就是過與不過的差別。 所以本書每一支讀大量資料的程
式，開頭都有這兩行。

\subsection{讀到檔案結束}

有些題目不告訴你有幾筆資料，資料一直到檔案結束為止。 \texttt{cin >> x} 這個運算式
本身有值：讀到東西時當成 \texttt{true}，讀不到（檔案結束，或內容不是數字）時當成
\texttt{false}。 所以 \texttt{while (cin >> x)} 會一直讀，讀到檔案結束就停。

\begin{code}[label=code:delivery]{cpp}{讀到檔案結束：不知道有幾筆的外送距離}
long long d, total = 0;
int count = 0;
while (cin >> d) {              // 讀不到數字（檔案結束）時停下
    count++;
    total += d;
}
cout << count << ' ' << total << '\n';
\end{code}

每個數字是一筆外送的距離（公尺），輸入 \texttt{1200}、\texttt{3500 800}、
\texttt{2750}（三行）時印出 \texttt{4 8250}。 \texttt{>>} 會跳過所有空白和換行，所
以一行有幾個數字都沒關係，檔案最後有沒有換行也沒關係。 在鍵盤上手動輸入時，要按
Ctrl+D（Windows 是 Ctrl+Z 再按 Enter）表示輸入結束；用 \texttt{./a.out < in.txt}
從檔案讀時，讀完檔案自然就結束了。

如果資料是「一行一筆」，而一行裡的個數不固定，就要知道哪些數字在同一行。 這時用
\texttt{getline} 讀一整行，再用 \texttt{istringstream} 把這一行當成一個小的輸入來拆。
\texttt{istringstream} 在 \texttt{<sstream>} 裡，用法和 \texttt{cin} 一樣：

\begin{example}[label=ex:orders]{cpp}{(orders.cpp) 每位客人點了幾樣}
#include <iostream>
#include <sstream>
#include <string>
using namespace std;

int main() {
    int n;
    cin >> n;
    string line;
    getline(cin, line);             // 讀掉 n 後面剩下的換行
    for (int i = 1; i <= n; i++) {
        getline(cin, line);         // 一整行：這位客人點的所有品項
        istringstream in(line);
        string item;
        int count = 0;
        while (in >> item) {        // 從這一行裡一個一個拆出品項
            count++;
        }
        cout << "#" << i << ' ' << count << '\n';
    }
    return 0;
}
\end{example}

\iosample{
3 \\
latte bagel \\
tea \\
mocha muffin scone \\
}{
\#1 2 \\
\#2 1 \\
\#3 3 \\
}

第 10 行的 \texttt{getline} 是必要的。 \texttt{cin >> n} 讀完 \texttt{3} 就停下來，後
面的換行還留在輸入裡；少了這一行，第一次 \texttt{getline} 會讀到那個空的行尾，程式
印出 \texttt{\#1 0}、\texttt{\#2 2}、\texttt{\#3 1}，整個錯開一行。 這是銜接書第~3~章
講過的陷阱，在混用 \texttt{>>} 和 \texttt{getline} 的題目裡一定會遇到。

\subsection{控制小數位數與寬度}

題目常要求「四捨五入到小數點後第二位」或「時間寫成 \texttt{09:05}」。 這些用
\texttt{<iomanip>} 裡的格式設定：\texttt{fixed} 加 \texttt{setprecision(k)} 印出小數
點後 $k$ 位；\texttt{setw(w)} 讓下一個輸出至少佔 $w$ 格，\texttt{setfill(c)} 決定空出
來的格子填什麼字元。

\begin{example}[label=ex:price]{cpp}{(price.cpp) 平均金額與補零的時間}
#include <iomanip>
#include <iostream>
using namespace std;

int main() {
    int n;
    cin >> n;
    long long total = 0;
    for (int i = 0; i < n; i++) {
        int x;
        cin >> x;
        total += x;
    }
    double avg = (double)total / n;
    cout << fixed << setprecision(2) << avg << '\n';   // 小數點後兩位
    int h, m;
    cin >> h >> m;
    cout << setfill('0') << setw(2) << h << ':' << setw(2) << m << '\n';
    return 0;
}
\end{example}

\iosample{
3 \\
1000 2050 1235 \\
9 5 \\
}{
1428.33 \\
09:05 \\
}

\begin{itemize}
\item \textbf{少了 \texttt{fixed}}：\texttt{setprecision(2)} 的意思變成「全部只留兩位
有效數字」，同一個平均值會印成 \texttt{1.4e+03}。
\item \textbf{設定會一直留著}：\texttt{fixed}、\texttt{setprecision}、\texttt{setfill}
設了之後對後面所有輸出都有效，之後再印 \texttt{3.0} 會變成 \texttt{3.00}。 只有
\texttt{setw} 是只管下一個輸出，所以上面的程式寫了兩次。
\item \textbf{四捨五入不一定照十進位}：\texttt{2.675} 在 \texttt{double} 裡其實比
$2.675$ 小一點，印到兩位是 \texttt{2.67}；剛好在中間的 \texttt{0.125} 也印成
\texttt{0.12}。 題目要求精確的四捨五入時，金額先用整數算（以分為單位），最後再轉。
\end{itemize}

同一套設定也能用在 \texttt{ostringstream} 上：把它當成 \texttt{cout} 來輸出，最後用
\texttt{.str()} 取出字串，就能先把格式化好的結果存起來，例如拼出一行再一次印出。

\subsection{例題：校慶攤位主題票選}
\label{sec:ex-vote}

\textbf{題意}（\akiref{voting-on-school-fair-booth-themes}{校慶攤位主題票選}）

\begin{itemize}
\item 輸入有很多組，每組是一個班級：第一行是選票張數 $n$ 和主題數 $k$；$n \ge 1$ 時下
一行有 $n$ 個主題編號（$1$ 到 $k$），$n = 0$ 時沒有這一行。
\item 最後一行是 \texttt{0 0}，表示輸入結束。
\item 每組輸出得票最多的主題編號，同票時取編號最小的；沒有人投票時輸出 $0$。
\item 組數最多 $5 \times 10^5$，每組的 $n$、$k$ 最多 $5 \times 10^5$，所有組的 $n + k$
加起來不超過 $10^6$。
\end{itemize}

例如輸入是 \texttt{3 5}、\texttt{5 5 2}、\texttt{0 4}、\texttt{1 2}、\texttt{2}、
\texttt{0 0}（每項一行）。 第一組主題 $5$ 得兩票，輸出 $5$；第二組有四個主題、沒有
人投票，輸出 $0$；第三組輸出 $2$。

\textbf{想法}

每一組的計票很直接：開一個計數陣列，\texttt{cnt[c]} 是主題 $c$ 的票數，讀一張票就加
一，最後從 $1$ 到 $k$ 找最大的。 這題要想清楚的是另外三件事。

\textbf{結束的條件。} 第二組的 \texttt{0 4} 是一組合法的資料：有四個主題，只是沒有人
投票。 只有 $n$ 和 $k$ 都是 $0$ 才是結束標記。 讀到 $n = 0$ 時也不用特別處理「沒有
選票那一行」，因為 \texttt{>>} 本來就會跳過換行，直接讀下一組的 $n$、$k$。

\textbf{每一組的工作量。} 用第~\ref{sec:from-limits}~小節的方法估。 如果宣告一個大小
$5 \times 10^5 + 1$ 的陣列，每組開始前整個清成 $0$，最壞的情況是 $5 \times 10^5$ 組
都要清一次，共 $2.5 \times 10^{11}$ 次，來不及。 題目限制的是所有組的 $n + k$ 總和，
所以每一組的工作量要和這一組的 $n + k$ 成正比。 最簡單的做法，是每一組各自開一個大小
$k + 1$ 的 \texttt{vector}：開的時候順便清成 $0$，花 $O(k)$；讀票 $O(n)$；找最大值
$O(k)$。 所有組加起來是 $O(\sum (n + k)) = O(10^6)$。

\textbf{同票和沒有票。} 從 \texttt{best = 0} 開始，主題 $c$ 的票數嚴格大於
\texttt{cnt[best]} 時才換成 $c$。 \texttt{cnt[0]} 永遠是 $0$，所以沒有人投票時
\texttt{best} 一直是 $0$，剛好是題目要的輸出；有人投票時，第一個票數大於 $0$ 的主題
會換上來。 用嚴格大於，同票時就會留下編號小的那個。

同樣的模型也出現在別的地方：客服系統每天統計哪一類問題被反映得最多，外送平台統計每
家店當天最多人點的品項。 資料一天一組、類別數每天不同，就是這題的形狀。

\textbf{程式}

\begin{example}[label=ex:vote]{cpp}{(vote.cpp) 校慶攤位主題票選}
#include <iostream>
#include <vector>
using namespace std;

int main() {
    ios::sync_with_stdio(false);
    cin.tie(nullptr);
    int n, k;
    while (cin >> n >> k) {
        if (n == 0 && k == 0) {           // 兩個都是 0 才是結束標記
            break;
        }
        vector<int> cnt(k + 1, 0);        // 這一組自己的計數，大小跟著 k
        for (int i = 0; i < n; i++) {
            int a;
            cin >> a;
            cnt[a]++;
        }
        int best = 0;                     // cnt[0] 一定是 0
        for (int c = 1; c <= k; c++) {
            if (cnt[c] > cnt[best]) {     // 嚴格大於才換，同票留小的編號
                best = c;
            }
        }
        cout << best << '\n';
    }
    return 0;
}
\end{example}

\iosample{
3 5 \\
5 5 2 \\
0 4 \\
1 2 \\
2 \\
0 0 \\
}{
5 \\
0 \\
2 \\
}

\begin{itemize}
\item \texttt{while (cin >> n >> k)} 同時處理兩種結束：讀到 \texttt{0 0}，或是檔案在那
之前就結束了。 題目保證有 \texttt{0 0}，多這一層保護不花任何成本。
\item \texttt{vector<int> cnt(k + 1, 0)} 每一組都重新建立，上一組的票數不會留下來，也
不用自己清。
\end{itemize}

\textbf{複雜度}

所有組加起來 $O(\sum (n + k))$，最多大約 $10^6$ 次操作，加上讀 $10^6$ 個整數。 最大
的測試資料在本書測試的環境下約 $0.03$ 秒。

\subsection{常見錯誤}

\begin{itemize}
\item \textbf{看到 $n = 0$ 就結束}：結束條件寫成 \texttt{if (n == 0) break;}，上面的範
例只會印出 \texttt{5}，第二組和第三組都漏掉了。
\item \textbf{每組清空整個大陣列}：宣告一個大小 $5 \times 10^5 + 1$ 的計數陣列，每組都
整個清成 $0$。 答案是對的，但組數很多的測試資料在本書測試的環境下要跑 $9$ 秒左右，
正確的程式只要 $0.03$ 秒。
\item \textbf{忘了清}：計數陣列宣告在迴圈外面，每一組都沒有清。 上面的範例會印出
\texttt{5}、\texttt{2}、\texttt{2}：第二組沒有人投票，卻讀到第一組留下的票數。
\item \textbf{同票用 \texttt{>=}}：比較寫成 \texttt{cnt[c] >= cnt[best]}，同票時會換成
編號大的。 輸入 \texttt{4 3}、\texttt{1 3 3 1}、\texttt{0 0} 會印出 \texttt{3}，正確
是 \texttt{1}。
\end{itemize}

\subsection{輸出的格式}

評分系統拿你的輸出和正確答案比對。 有些系統會忽略行尾的空白和最後多出來的空行，有
些逐字比對；照最嚴格的那種準備就不用猜。 數字之間的空白、換行、\texttt{Case \#1: 5}
這種格式的大小寫和冒號，都照題目寫；一行印好幾個數字時，把空白印在數字「前面」（第
一個除外），最後就不會多一個空白；題目沒要求的提示文字一個字都不要印；除錯用的輸出改
印到 \texttt{cerr}，它不會被比對，但交出去前還是刪掉，印太多也花時間。

\section{自己測試：範例之外}
\label{sec:testing}

題目給的範例是用來說明題意的，通常很小，也不會故意去碰每一個邊界。 本章前面的常見
錯誤裡，有好幾個錯法都能通過全部範例：用 \texttt{sqrt} 求整數平方根、片段 1 直接跑、
清空整個大陣列。 交出去之前，至少要做兩件事：自己補邊界的輸入，以及在大量的隨機小輸
入上，拿暴力解來比對。

\subsection{手造邊界輸入}

想邊界輸入時，從題目的每一個限制下手：

\begin{itemize}
\item \textbf{最小}：$n = 1$、$n = 0$（題目允許的話）、只有一組資料。
\item \textbf{最大}：每個數都取上限，檢查溢位和時間。 迴圈次數的例題要測
$n = 10^{18}$，主題票選要測「很多組、每組很小」和「一組、很大」兩種。
\item \textbf{剛好在分界上}：題目的條件用到 $\le$ 或 $<$ 的地方，取剛好相等的值。 例如
$T$ 剛好是某個完全平方數、比它少一，或 $n$ 剛好是 $2$ 的次方、比它少一。
\item \textbf{全部相同、已經排好、完全反過來}：很多錯誤只在這些情況出現。
\item \textbf{格式的邊界}：結束標記前一組是空的（\texttt{0 4}）、檔案最後沒有換行。
\end{itemize}

每一筆都要先手算出正確答案，再和程式比。 手算不出來的大輸入，至少看程式有沒有在時
間內跑完，答案的大小合不合理。

\subsection{用暴力解比對}

手造的輸入受限於你想得到的情況。 想不到的那些，交給電腦：寫一個慢但簡單的暴力解，
產生幾千組隨機的小輸入，看兩支程式的輸出是否完全相同。 這種做法俗稱
\textbf{對拍}。 需要三個程式和一個腳本：

\begin{center}
\begin{tikzpicture}[node distance=0.8cm and 0.8cm, font=\small]
\node[process] (gen) {產生器};
\node[ioNode, right=of gen] (in) {\texttt{test.in}};
\node[process, above right=0.2cm and 0.8cm of in] (mine) {要測的程式};
\node[process, below right=0.2cm and 0.8cm of in] (brute) {暴力解};
\node[decision, right=3.4cm of in, aspect=1.8, inner sep=1pt] (cmp) {相同？};
\draw[flowarrow] (gen) -- (in);
\draw[flowarrow] (in) -- (mine.west);
\draw[flowarrow] (in) -- (brute.west);
\draw[flowarrow] (mine.east) -- (cmp);
\draw[flowarrow] (brute.east) -- (cmp);
\draw[flowarrow] (cmp.south) -- ++(0,-1.3) -| node[pos=0.25, below] {相同：換下一個種子} (gen.south);
\node[right=0.4cm of cmp, align=left] (stop) {不同：印出\\\texttt{test.in}，停下};
\draw[flowarrow] (cmp) -- (stop);
\end{tikzpicture}
\diagcaption{用暴力解比對的流程。 每一輪用不同的亂數種子產生一組輸入。}
\label{fig:check-loop}
\end{center}

以主題票選為例。 暴力解要寫得和正式的程式越不一樣越好，這樣兩邊才不會犯同一個錯。
正式的程式用計數陣列，暴力解就對每個主題把整組票從頭數一次，$O(nk)$，很慢，但幾乎
不可能寫錯：

\begin{example}[label=ex:vote-brute]{cpp}{(vote\_brute.cpp) 主題票選的暴力解}
#include <iostream>
#include <vector>
using namespace std;

// 暴力解：每個主題都把整組票從頭數一次
int main() {
    int n, k;
    while (cin >> n >> k) {
        if (n == 0 && k == 0) {
            break;
        }
        vector<int> a(n);
        for (int &x : a) {
            cin >> x;
        }
        int best = 0, bestVotes = 0;
        for (int c = 1; c <= k; c++) {
            int votes = 0;
            for (int x : a) {
                if (x == c) {
                    votes++;
                }
            }
            if (votes > bestVotes) {
                best = c;
                bestVotes = votes;
            }
        }
        cout << best << '\n';
    }
    return 0;
}
\end{example}

\iosample{
5 3 \\
2 3 2 1 3 \\
4 4 \\
4 1 4 1 \\
0 0 \\
}{
2 \\
1 \\
}

產生器需要亂數。 C++ 用 \texttt{<random>} 裡的 \texttt{mt19937} 產生亂數：建立時給
一個\textbf{種子}（seed），之後每次呼叫 \texttt{rng()} 得到一個 $0$ 到 $2^{32} - 1$ 的
整數；再交給 \texttt{uniform\_int\_distribution<int>(lo, hi)}，就變成 \texttt{lo} 到
\texttt{hi} 之間平均分布的整數。 同一個種子一定得到同一串亂數：

\begin{example}[label=ex:rng]{cpp}{(rng\_demo.cpp) 固定種子的亂數}
#include <iostream>
#include <random>
using namespace std;

int main() {
    mt19937 rng(2026);                      // 種子固定，每次執行都一樣
    cout << rng() << ' ' << rng() << ' ' << rng() << '\n';  // 0 到 2^32 - 1
    uniform_int_distribution<int> die(1, 6);    // 1 到 6 平均分布
    for (int i = 0; i < 10; i++) {
        cout << (i > 0 ? " " : "") << die(rng);
    }
    cout << '\n';
    return 0;
}
\end{example}

\iosample{}{
942082305 3292560774 1773871898 \\
6 6 6 1 5 3 1 6 4 2 \\
}

第一行在任何編譯器上都一樣，因為 C++ 標準把 \texttt{mt19937} 的演算法規定死了；第二
行經過 \texttt{uniform\_int\_distribution}，標準沒有規定換算的細節，換一套標準函式庫
可能不同，但同一個環境裡一定重現得出來。 C 的 \texttt{rand()} 不適合用在這裡：它的
最大值 \texttt{RAND\_MAX} 在有些環境下只有 $32767$；\texttt{rand() \% k} 在 $k$ 不整除
範圍時，小的數字出現得比較多；種子還要另外用 \texttt{srand} 設，整支程式共用一個狀
態。

主題票選的產生器從命令列讀種子，產生一組合法的小輸入。 種子不同，輸入就不同；出錯時
用同一個種子就能重現：

\begin{example}[label=ex:vote-gen]{cpp}{(vote\_gen.cpp) 主題票選的隨機輸入產生器}
#include <iostream>
#include <random>
#include <string>
using namespace std;

// 用法：./vote_gen 種子
// 產生 1 到 3 組小資料，最後一行是 0 0
int main(int argc, char *argv[]) {
    if (argc < 2) {
        cerr << "usage: vote_gen seed\n";
        return 1;
    }
    mt19937 rng(stoul(argv[1]));
    auto randInt = [&](int lo, int hi) {   // lo 到 hi 之間的整數
        return uniform_int_distribution<int>(lo, hi)(rng);
    };
    int groups = randInt(1, 3);
    for (int g = 0; g < groups; g++) {
        int n = randInt(0, 5), k = randInt(1, 4);
        cout << n << ' ' << k << '\n';
        for (int i = 0; i < n; i++) {
            cout << randInt(1, k) << (i + 1 < n ? ' ' : '\n');
        }
    }
    cout << "0 0\n";
    return 0;
}
\end{example}

\begin{itemize}
\item \texttt{main} 的兩個參數是命令列引數：\texttt{argc} 是個數，\texttt{argv[1]} 是
程式名稱後面的第一個字串。 執行 \texttt{./vote\_gen 7} 時，\texttt{argv[1]} 是
\texttt{"7"}，用 \texttt{stoul} 轉成整數當種子。
\item 範圍故意設得很小：最多三組、每組最多五張票、四個主題。 小輸入出錯時容易手算，
而且同票、沒有票、\texttt{n = 0} 這些情況會經常出現。
\end{itemize}

最後是把它們串起來的腳本。 下面是 Linux 和 macOS 的 bash 寫法；Windows 可以在 WSL
或 Git Bash 裡執行同一個腳本：

\begin{code}[label=code:check-sh]{bash}{check.sh：用暴力解比對}
#!/bin/bash
# 用法：./check.sh 要測的程式 暴力解 產生器 次數
for ((i = 1; i <= $4; i++)); do
    ./$3 $i > test.in
    ./$1 < test.in > mine.out
    ./$2 < test.in > brute.out
    if ! diff -q mine.out brute.out > /dev/null; then
        echo "第 $i 組不一樣"
        cat test.in
        exit 1
    fi
done
echo "$4 組都一樣"
\end{code}

先 \texttt{chmod +x check.sh} 讓它可以執行，再用 \texttt{./check.sh vote vote\_brute
vote\_gen 1000} 跑一千組。 範例~\ref{ex:vote} 的 \texttt{vote} 會印出
\texttt{1000 組都一樣}。

來看它抓到錯誤的樣子。 有人用全域的計數陣列，每組結束時把用過的格子清掉，可是清的
範圍寫錯了：

\begin{code}[label=code:vote-bug]{cpp}{清錯範圍的計數（錯誤示範）}
int cnt[500001];                          // 全域陣列，一開始全是 0
// ……讀票、找最大值、輸出，和 vote.cpp 相同……
for (int c = 1; c <= n; c++) {            // 應該是 c <= k
    cnt[c] = 0;
}
\end{code}

把 $k$ 寫成 $n$ 是很容易打錯的地方，而且兩個範例都會通過：範例裡每一組清掉的範圍剛
好都蓋住了投過票的主題。 用腳本比對，在本書測試的環境下第 46 組就出現不同：

\begin{code}{text}{check.sh 印出的輸入}
第 46 組不一樣
1 3
3
1 4
4
1 2
2
0 0
\end{code}

暴力解印出 \texttt{3}、\texttt{4}、\texttt{2}，這支程式印出 \texttt{3}、\texttt{3}、
\texttt{2}。 第一組 $n = 1$，只清掉 \texttt{cnt[1]}，主題 $3$ 的那一票留到了第二組，
和主題 $4$ 同票，於是輸出較小的 $3$。 這組輸入只有七行，手算一遍就能看出錯在哪裡。

用暴力解比對時，幾件事要注意：

\begin{itemize}
\item \textbf{暴力解要直接。} 照題目的字面意思枚舉，不要在暴力解裡用同一個「聰明的觀
察」。 兩邊都用了同一個錯的想法，就會一起錯，比對不出來。
\item \textbf{兩邊不同時，先手算。} 有時候錯的是暴力解，或是產生器產生了題目不允許的輸
入。
\item \textbf{比對全部相同不等於證明。} 範圍小的隨機資料抓不到溢位和浮點數誤差，也抓
不到太慢。 這些要靠上一小節的最大輸入，以及下面的工具。
\end{itemize}

迴圈次數的例題也能這樣做：暴力解就是第~\ref{sec:ex-loop-count}~小節照字面跑的那支程式，產生器隨機產生
$n \le 1000$ 的詢問。 不過它的浮點數錯誤只在 $n$ 超過 $2^{53}$ 時才會出現，暴力解根本
跑不到那裡，只能靠手造的輸入（$n = 2^{59} - 1$）抓。 兩種測試方法要一起用。

\subsection{讓編譯器幫忙}

練習時用下面的選項編譯：

\begin{code}{bash}{練習時的編譯選項}
g++ -std=c++17 -O2 -Wall -Wextra -fsanitize=address,undefined vote.cpp -o vote
\end{code}

\texttt{-Wall -Wextra} 讓編譯器對可疑的寫法提出警告，例如比較有號和無號的整數、變數
沒有初始化就使用。 \texttt{-fsanitize=address,undefined} 會在執行時檢查陣列越界和有
號整數溢位，出事時印出是哪一行（銜接書第~10~章 10.1 節有示範）。 第~\ref{sec:overflow}~
節的溢位程式這樣編譯後執行，會印出 \texttt{runtime error: signed integer overflow:
100000 * 100000 cannot be represented in type 'int'} 這樣的訊息，前面還有檔名和行號。
加了這些選項，程式會慢好幾倍，所以只在測試時用；量時間時拿掉。

量時間用 \texttt{time} 指令，例如 \texttt{time ./vote < big.in > /dev/null}，看
\texttt{real} 那一行。 \texttt{big.in} 可以用另一個產生器做出來：把
\texttt{vote\_gen} 的範圍改成題目的上限，就是一個最大輸入的產生器。

\begin{apcsnote}
實作題由系統用測試資料評分，考場上看不到這些測試資料，只有題目的範例。 練習時就養
成習慣：寫完先跑範例，再手造幾筆邊界輸入；時間夠的話，寫一個暴力解在小範圍比對。 考
試時不一定有時間寫完整的腳本，但手造最小、最大、剛好相等這三種輸入，幾分鐘就做得完。
\end{apcsnote}

\section{本章常見錯誤}
\label{sec:prep-pitfalls}

前面各節的錯誤，歸納起來是五類。 寫每一題時，交出去之前逐項檢查一次。

\begin{center}
\begin{tabular}{p{0.2\linewidth}p{0.48\linewidth}p{0.2\linewidth}}
\toprule
錯誤 & 例子 & 見 \\
\midrule
估錯工作量 & 看到兩層迴圈就說 $O(n^2)$ 或 $n \log n$；只看 $n$ 沒看詢問數 $q$；每組都清空整個大陣列 &
第~\ref{sec:ex-loop-count}、\ref{sec:from-limits}、\ref{sec:ex-vote}~節 \\
溢位 & \texttt{int} 乘完才轉 \texttt{long long}；\texttt{1 << b}；比大小時直接相乘 &
第~\ref{sec:overflow}~節 \\
浮點數誤差 & 用 \texttt{sqrt}、\texttt{cbrt}、\texttt{log2} 求整數答案，沒有修正 &
第~\ref{sec:ex-loop-count}、\ref{sec:ex-max-size}~節 \\
輸入輸出 & 沒加速、用 \texttt{endl}；\texttt{>>} 後面直接 \texttt{getline}；結束標記判斷錯；多印空白 &
第~\ref{sec:io}~節 \\
只用範例測試 & 清錯範圍、浮點數誤差、太慢，範例都會通過 &
第~\ref{sec:testing}~節 \\
\bottomrule
\end{tabular}
\end{center}

\section{本章練習}
\label{sec:ch1-exercises}

\begin{drills}[label=drills:ch1]
\item 下面兩段程式，$n = 10$ 時 \texttt{cnt++} 各執行幾次？ 一般的 $n$ 時，複雜度是多少？
\begin{code}{cpp}{練習 1}
// (a)
for (int i = 1; i <= n; i++)
    for (int j = i; j <= n; j += i)
        cnt++;
// (b)
for (int i = n; i >= 1; i /= 2)
    for (int j = 1; j <= n; j++)
        cnt++;
\end{code}

\item 程式甲處理 $n$ 筆資料要 $2n^2$ 步，程式乙要 $200n + 5000$ 步。
\begin{enumerate}
\item $n = 50$ 和 $n = 1000$ 時，哪一支比較快？
\item 找出最小的 $n$，使得從這個 $n$ 開始，程式乙一直都嚴格比程式甲快。
\end{enumerate}

\item 下面每一題，判斷想用的做法來不來得及（用每秒 $10^8$ 次估），並說明理由。
\begin{enumerate}
\item $n \le 2 \times 10^5$ 個數、$q \le 2 \times 10^5$ 個詢問，每個詢問把陣列從頭掃
一遍。
\item 最多 $10^4$ 組資料，每組 $n \le 1000$，每組用 $O(n^2)$ 的做法。
\item 同 (b)，但題目另外保證所有組的 $n$ 加起來不超過 $10^5$。
\item $n \le 20$，枚舉所有子集合，每個子集合花 $O(n)$ 檢查。
\end{enumerate}

\item 下面的程式片段中，$n \le 10^5$，$0 \le$ \texttt{a[i]} $\le 10^6$。 哪幾行可能
溢位？ 怎麼改？
\begin{code}{cpp}{練習 4}
long long sum = 0, sq = 0;
for (int i = 0; i < n; i++) {
    sum += a[i];                       // (1)
    sq += a[i] * a[i];                 // (2)
}
int avg = sum / n;                     // (3)
long long total = n * (n + 1) / 2;     // (4)
\end{code}

\item 為第~\ref{sec:ex-loop-count}~小節的例題設計五筆手造的輸入，每一筆寫出它要檢查的
是哪一種錯誤，以及正確的輸出。
\end{drills}

\section{實作練習}
\label{sec:ch1-practice}

下面是 Aki-Judge 上和本章相關的題目。 難度分四階：暖身是直接套用本章的寫法；基本要
自己想一下怎麼估計或怎麼讀輸入；實戰的題敘比較長，要先把題目轉換成本章的形狀；挑戰
需要再多一步觀察。

\begin{center}
\begin{tabular}{llc}
\toprule
小節 & 題目 & 難度 \\
\midrule
估計次數與複雜度 & \akiref{stage-runs-under-the-entry-cap}{依入口上限計算的連續關卡} & 基本 \\
 & \akiref{items-ordered-exactly-once}{只出現一次的值} & 實戰 \\
 & \akiref{rectangles-with-an-even-number-of-stamps}{章數為偶數的矩形} & 實戰 \\
 & \akiref{partner-pairs-at-two-scales}{兩種規模的搭檔配對} & 實戰 \\
 & \akiref{faster-from-here-on}{從此以後都比較快} & 挑戰 \\
\midrule
整數運算與溢位 & \akiref{sheets-for-a-saddle-stitched-booklet}{騎馬釘小冊子的紙張數} & 暖身 \\
 & \akiref{position-on-a-circular-track}{環形跑道上的位置} & 暖身 \\
 & \akiref{splitting-the-class-trip-bill}{班遊費用分攤} & 基本 \\
 & \akiref{download-progress-bar}{下載進度條} & 基本 \\
 & \akiref{stairs-or-the-express-elevator}{樓梯還是快速電梯} & 實戰 \\
 & \akiref{per-purchase-vs-monthly-loyalty-points}{逐筆給點與月結給點的差距} & 挑戰 \\
\midrule
輸入與輸出 & \akiref{jump-rope-record-breaks}{跳繩計數器的破紀錄次數} & 暖身 \\
 & \akiref{ballots-of-unknown-count}{不知道有幾行的投票紀錄} & 暖身 \\
 & \akiref{steps-walked-during-each-break}{每天各節下課的步數} & 暖身 \\
 & \akiref{drink-shop-sales-per-shift}{飲料店每班的營業額} & 暖身 \\
 & \akiref{study-group-sign-in-sheets}{讀書會簽到表} & 基本 \\
 & \akiref{study-room-seat-log}{自習室的座位登記} & 實戰 \\
\bottomrule
\end{tabular}
\end{center}

本章的三題例題也在 Aki-Judge 上，讀完各節之後可以先自己寫一次再對照。

下面這組是 STL 的複習題，用到的是銜接書第~3、8、9~章的 \texttt{vector}、
\texttt{pair}、\texttt{sort} 與比較函式、\texttt{set}／\texttt{map}。 從第~2~章起，
這些容器會直接拿來用，不再說明。 寫起來卡住的話，先回頭翻銜接書對應的那一章。

\begin{center}
\begin{tabular}{llc}
\toprule
小節 & 題目 & 難度 \\
\midrule
排序與比較函式 & \akiref{sorting-the-camp-roster}{營隊報名名冊排序} & 暖身 \\
 & \akiref{graduation-photo-rows}{畢業照的分排站位} & 暖身 \\
 & \akiref{league-ranks-by-win-rate}{校際聯賽勝率名次} & 基本 \\
 & \akiref{sorting-file-names-the-human-way}{照人的直覺排檔名} & 基本 \\
 & \akiref{clicking-column-headers}{試算表連點欄位標題} & 實戰 \\
 & \akiref{winter-break-reading-marathon-standings}{寒假閱讀馬拉松結算} & 實戰 \\
 & \akiref{order-of-the-lighthouse-sweep}{燈塔光束掃過的順序} & 挑戰 \\
\midrule
\texttt{set} 與 \texttt{map} & \akiref{hallway-light-switch-log}{走廊燈的開關紀錄} & 暖身 \\
 & \akiref{the-calculators-variable-table}{小算盤的變數表} & 暖身 \\
 & \akiref{class-bonus-and-personal-bonus}{全班加分與個人加分} & 基本 \\
 & \akiref{cone-squares-on-the-field}{操場上的方格標記} & 基本 \\
 & \akiref{transposed-or-mirrored}{移調與倒影} & 實戰 \\
 & \akiref{matching-prefixes-of-two-attendance-lists}{兩份點名單的前綴配對} & 實戰 \\
 & \akiref{left-and-right-socks-in-a-row}{左右腳都配成雙的區段} & 挑戰 \\
\bottomrule
\end{tabular}
\end{center}

\section{讀完你應該會}
\label{sec:ch1-checklist}

\begin{itemize}
\item 給一段內層範圍跟著外層變、或變數每次乘除一個數的迴圈，能列表算出它的精確執行次
數，並寫出它的複雜度。
\item 看到題目的限制，能把 $n$、$q$、組數都代進自己的做法，判斷來不來得及；題目限制
「總和」時，能讓每一組的工作量跟那一組的大小成正比。
\item 寫程式前能估出答案和中間值的最大值，在需要的地方用 \texttt{long long}、
\texttt{1LL}，並用除法比較可能溢位的乘積。
\item 會用 \texttt{while (cin >> x)} 讀到檔案結束、用 \texttt{getline} 加
\texttt{istringstream} 逐行拆資料，並正確處理多組資料的結束標記。
\item 能為一題寫出暴力解、隨機產生器和比對腳本，找出一組讓兩支程式輸出不同的小輸入。
\end{itemize}

\textbf{延伸閱讀}（書末會整理完整的參考資料）

\begin{itemize}
\item Antti Laaksonen，\textit{Competitive Programmer's Handbook}，第 1、2 章：輸入輸
出、整數型別和時間複雜度的簡短整理，並用最大子陣列和示範一個做法怎麼一步步變快。
\item Steven Halim 等，\textit{Competitive Programming 4}，Book 1 的 1.3 節：怎麼從
限制估計複雜度，以及測試時該準備哪些輸入。
\item Steven S. Skiena，\textit{The Algorithm Design Manual}，第 3 版，第 2 章：
big-O 的定義和常見的成長速度，寫得比本章嚴謹。
\end{itemize}

\addcontentsline{toc}{section}{練習解答}
\section*{練習解答}

\begin{enumerate}
\item
\begin{enumerate}
\item $27$ 次。 外層是 \texttt{i} 時，\texttt{j} 是 $i, 2i, 3i, \ldots$，共
$\lfloor n / i \rfloor$ 次，$n = 10$ 時是 $10 + 5 + 3 + 2 + 2 + 1 + 1 + 1 + 1 + 1
= 27$。 一般的 $n$ 是 $\lfloor n/1 \rfloor + \lfloor n/2 \rfloor + \cdots +
\lfloor n/n \rfloor$，大約是 $n$ 乘上 $1 + 1/2 + \cdots + 1/n$，而這個和大約是
$\ln n$，所以是 $O(n \log n)$。
\item $40$ 次。 外層 \texttt{i} 是 $10, 5, 2, 1$，共 $4$ 輪，每輪內層 $10$ 次。 一般
的 $n$ 是 $n \times (\lfloor \log_2 n \rfloor + 1)$，$O(n \log n)$。
\end{enumerate}

\item
\begin{enumerate}
\item $n = 50$ 時甲要 $5000$ 步，乙要 $15000$ 步，甲比較快。 $n = 1000$ 時甲要
$2 \times 10^6$ 步，乙要 $205000$ 步，乙比較快。
\item $121$。 $n = 120$ 時甲 $28800$、乙 $29000$，甲還是比較快；$n = 121$ 時甲
$29282$、乙 $29200$，乙比較快。 兩者的差 $2n^2 - 200n - 5000$ 在 $n \ge 50$ 之後只
會越來越大，所以從 $121$ 起乙一直比較快。
\end{enumerate}

\item
\begin{enumerate}
\item 不行。 $n \times q = 4 \times 10^{10}$。 要換成第~5~章的前綴和這類先處理好、
每個詢問 $O(1)$ 或 $O(\log n)$ 的做法。
\item 不行。 每組最多 $10^6$ 步，$10^4$ 組是 $10^{10}$。
\item 可以。 每組 $n_i^2 \le n_i \times 1000$，加起來不超過
$1000 \times \sum n_i = 1000 \times 10^5 = 10^8$。 總和的限制讓「很多組都很大」的情況
不會發生。
\item 可以。 $2^{20} \times 20 \approx 2 \times 10^7$。
\end{enumerate}

\item 第 (2) 行和第 (4) 行。
\begin{itemize}
\item (2) \texttt{a[i] * a[i]} 用 \texttt{int} 乘，最大 $10^{12}$，溢位。 改成
\texttt{(long long)a[i] * a[i]}。 加起來最多 $10^{17}$，\texttt{sq} 存得下。
\item (1) 和 (3) 沒問題：\texttt{sum} 是 \texttt{long long}，最多 $10^{11}$；平均值不超過 $10^6$。
\item (4) \texttt{n * (n + 1)} 用 \texttt{int} 乘，$n = 10^5$ 時大約 $10^{10}$，溢位。
改成 \texttt{(long long)n * (n + 1) / 2}。
\end{itemize}

\item 一種答法（輸入的第一行 $q$ 省略）：
\begin{itemize}
\item \texttt{3 1}：$n$ 最小，片段 3 的 \texttt{j} 從 $0$ 開始，答案 $2$。 漏掉
$+1$ 會印 $1$。
\item \texttt{1 1}：片段 1 只有一格，答案 $1$。
\item \texttt{2 7} 和 \texttt{2 8}：$n$ 剛好在 $2$ 的次方前後，答案 $7$ 和 $15$，檢查
迴圈條件的 $\le$。
\item 片段 1 取最大的 $n$（\texttt{1} 後面接 $10^{18}$）：檢查時間和溢位，答案是
\texttt{1999999999999999976}。
\item \texttt{2 576460752303423487}：$n = 2^{59} - 1$，檢查有沒有用浮點數算輪數，答案
是 $576460752303423487$。
\end{itemize}
\end{enumerate}

% 佔位：本章尚未開始。book.tex 裡先註解掉。
\chapter{排序與有序資料}

本章尚未開始。

% 佔位：本章尚未開始。book.tex 裡先註解掉。
\chapter{二分搜尋答案}

本章尚未開始。

% 佔位：本章尚未開始。book.tex 裡先註解掉。
\chapter{枚舉、模擬與回溯}

本章尚未開始。

% 佔位：本章尚未開始。book.tex 裡先註解掉。
\chapter{前綴和、雙指標與滑動視窗}

本章尚未開始。

% 佔位：本章尚未開始。book.tex 裡先註解掉。
\chapter{堆疊、佇列與優先佇列}

本章尚未開始。

\chapter{貪心}
\label{ch:greedy}

有一類題目，答案可以一步一步決定：每一步挑當下看起來最好的那一個，挑了就
不再更改，最後得到的剛好就是最佳解。 這種做法叫\textbf{貪心}（greedy）。
程式通常很短，常常只是一次排序加上一個迴圈。 難的地方在別處：看起來合理的
挑法往往有好幾種，其中大多數是錯的，而錯的那幾種通常也能通過題目給的範例。

這一章先講兩件事：什麼時候可以考慮貪心，以及怎麼確認一個貪心是對的（交換
論證）或是錯的（找反例）。 接著用三題例題練三種常見的形狀：兩群東西排序後
配對、從一堆區間裡挑出最多個互不重疊的、把區間的端點排好之後由左往右掃描。

用到的語法在《C++ 程式設計：給學過 C 的人》（以下簡稱銜接書）都講過：
\texttt{sort} 與比較函式、lambda、\texttt{next\_permutation} 在第~9~章，
\texttt{pair} 與結構化綁定在第~8~章，解題程式的骨架和從資料範圍估計時間在
第~10~章 10.1 節。

\section{選了就不回頭}
\label{sec:greedy-idea}

先看一個小問題。 書包最多能裝 $W$ 公斤，桌上有 $n$ 本書，第 $i$ 本重
$w_i$ 公斤，最多能帶幾本？

直覺的做法是從最輕的開始放，放到下一本放不下為止。 例如 $W = 10$，五本書
重 $4, 2, 7, 3, 5$。 排好是 $2, 3, 4, 5, 7$，放進 $2$、$3$、$4$ 共 $9$
公斤，再放 $5$ 就超過了，所以帶 $3$ 本。

這個做法有貪心的兩個特徵：

\begin{itemize}
\item 有一個固定的順序（由輕到重），照這個順序一個一個決定。
\item 每一步只看眼前：這本放得下就放，放了就不拿出來，也不回頭考慮別的組合。
\end{itemize}

要保證找到最佳解，最直接的方法是枚舉所有的子集合。 $n$ 本書有 $2^n$ 種帶
法，$n = 40$ 就超過一兆種。 貪心只要排序加一次掃描，是 $O(n \log n)$。 速
度差這麼多，代價是要回答一個問題：只看眼前，為什麼不會錯過更好的組合？

\subsection{什麼時候想到貪心}

下面幾種情況出現時，可以先試試貪心：

\begin{itemize}
\item 題目問「最多幾個」「最少幾次」「總和最小是多少」，而 $n$ 到
$2 \times 10^5$，照銜接書 10.1 節的表只容得下 $O(n \log n)$，枚舉所有可
能一定來不及。
\item 解可以看成一連串的選擇。 排成某種順序之後，每一次選擇好像只跟「目前做
到哪裡」有關，例如目前裝了多重、目前的時間到了哪裡。
\item 手算幾個小例子，發現最佳解總是很規律：最輕的那幾本、最早結束的那一個、
大的配大的。
\end{itemize}

這些只是線索。 想到一個挑法之後，要同時做兩件事：想辦法說明它為什麼對，也試
著找它出錯的例子。 哪一邊先有結果，就知道該不該用它。

\subsection{交換論證}

說明貪心正確，最常用的方法叫\textbf{交換論證}（exchange argument）。 做法
是隨便拿一個最佳解來，如果它和貪心選的不一樣，就把它的一部分換成貪心選的，
說明換完之後仍然合法，而且沒有變差。 一直換下去，最佳解就變成了貪心的解，
所以貪心的解也是最佳的。

用書包的例子走一次。 設最佳的帶法帶了 $k$ 本書，把所有的書由輕到重排好，
最輕的 $k$ 本叫做 $L$。

\begin{enumerate}
\item 如果最佳的帶法不是 $L$，那就有一本 $L$ 裡的書 $a$ 沒有被帶，也有一本
不在 $L$ 裡的書 $b$ 被帶了。 $b$ 不在最輕的 $k$ 本裡，所以 $w_a \le w_b$。
\item 把 $b$ 拿出來，換成 $a$。 本數還是 $k$，總重量的變化是
$w_a - w_b \le 0$，原本放得下，換完也放得下。
\item 每換一次，這個帶法和 $L$ 相同的書就多一本，最多換 $k$ 次就變成 $L$。
所以 $L$ 整組放得下。
\end{enumerate}

貪心從最輕的開始放，$L$ 的 $k$ 本一定都放得進去，所以至少帶 $k$ 本；$k$
已經是最多，貪心帶的就剛好是 $k$ 本。

交換論證可以照固定的格式寫：

\begin{itemize}
\item 寫出要比較的兩個東西：貪心的選擇，和任意一個最佳解。
\item 找一個兩者不一樣的地方，把最佳解的那一部分換成貪心的選擇，說明題目的
限制仍然滿足，而且答案沒有變差。
\item 說明換的動作會停下來，例如每換一次，兩者相同的部分就多一點。
\end{itemize}

第二步是重點，也是最容易寫得含糊的地方。 「換成比較輕的，應該比較好」不算
說明。 要寫出換完之後哪個量變了、變了多少，以及為什麼不會違反題目的限制。

\subsection{先找反例}

證明一直寫不出來的時候，這個貪心常常就是錯的。 這時可以反過來，找一組讓它
出錯的小輸入，叫做\textbf{反例}（counterexample）。 一組反例就足以說明一
條規則不能用。

例子：郵局只有 $1$ 元、$5$ 元、$6$ 元三種郵票，要剛好貼滿 $n$ 元，最少要
幾張？ 貪心的想法是每次先貼面額最大、而且還貼得下的那一種。 $n = 12$ 時貼兩
張 $6$ 元，是對的。 $n = 10$ 時先貼一張 $6$ 元，剩下 $4$ 元只能貼四張 $1$
元，共 $5$ 張；可是兩張 $5$ 元就夠了。

反例不一定猜得到。 比較可靠的方法是寫一個保證正確、但是很慢的\textbf{暴力
解}，在小範圍裡和貪心逐一比對，讓電腦去找：

\begin{example}[label=ex:stamps]{cpp}{(stamps.cpp) 讓電腦找貪心的反例}
#include <iostream>
using namespace std;

// 面額 6、5、1 元，每次先貼面額最大、還貼得下的郵票
int greedyCount(int amount) {
    int count = 0;
    for (int value : {6, 5, 1}) {
        count += amount / value;
        amount %= value;
    }
    return count;
}

// 枚舉 6 元和 5 元各用幾張，剩下的用 1 元補
int bestCount(int amount) {
    int best = amount;                    // 全部用 1 元
    for (int six = 0; six * 6 <= amount; six++) {
        for (int five = 0; six * 6 + five * 5 <= amount; five++) {
            int ones = amount - six * 6 - five * 5;
            best = min(best, six + five + ones);
        }
    }
    return best;
}

int main() {
    for (int amount = 1; amount <= 20; amount++) {
        int g = greedyCount(amount), b = bestCount(amount);
        if (g != b) {
            cout << amount << ' ' << g << ' ' << b << '\n';
        }
    }
    return 0;
}
\end{example}

\iosample{}{
10 5 2 \\
15 5 3 \\
16 6 3 \\
20 5 4 \\
}

每一行是「金額、貪心的張數、最少的張數」。 $20$ 元以內有四個金額出錯，最
小的是 $10$ 元。 \texttt{bestCount} 用雙層迴圈枚舉 $6$ 元和 $5$ 元各用幾
張，金額一大就很慢，但拿來檢查小範圍剛好。

這個寫法之後會一直用到：有了一個貪心的想法，先寫暴力解，在幾千組隨機的小輸
入上和貪心比對。 都一樣，再花時間想證明；出現一組不一樣，就省下了寫證明和交
出錯誤程式的時間。 比對全部相同不等於證明，但已經擋掉了大部分錯誤的想法。

順帶一提，面額 $1, 5, 10, 50$ 的硬幣用同樣的貪心找零是對的，換成
$1, 5, 6$ 就錯了。 同一條規則對不對，取決於題目的數字有什麼性質，沒辦法從
「看起來合理」判斷。 一般的面額要用第~9~章的動態規劃。

\subsection{常見錯誤}

\begin{itemize}
\item \textbf{只用範例測資檢查貪心}：範例是為了說明題意而設計的，很多錯的貪
心都會通過。 郵票的題目如果範例是 $12$ 元，錯的貪心也會印出正確的 $2$。 交
出去之前，至少用暴力解在小範圍比對一次。
\item \textbf{拿幾個例子當證明}：「試了五組都對」不能說明它永遠對。 反過來，
一組反例就足以推翻一條規則。
\item \textbf{暴力解自己也錯}：比對時兩邊不一樣，先手算那組輸入，確定是哪一
邊錯。 暴力解要寫得越直接越好，例如老老實實枚舉所有可能。 在暴力解裡又用了
一次同樣的貪心想法，兩邊會一起錯，比對就沒有意義了。
\end{itemize}

\section{排序後配對}
\label{sec:greedy-matching}

兩群東西要一對一配對時，常見的貪心是兩邊都排序，再依名次對齊：最小的配最
小的，第二小的配第二小的。 這一節的例題就是這個形狀。 重點有兩個：寫出它為
什麼對，以及認出一個看起來也很合理、其實會出錯的做法。

\subsection{例題：全勝所需的最少加練次數}
\label{sec:ex-training}

\textbf{題意}（\akiref{fewest-drills-for-a-clean-sweep}{全勝所需的最少加練次數}）

\begin{itemize}
\item 我方有 $n$ 位社員，對方有 $m$ 位選手（$m \le n$），每人有一個戰力值。
\item 對方每位選手都要和我方一位社員比一場，每位社員最多比一場。 戰力嚴格大
於對手才算贏，平手算輸。
\item 比賽前可以加練，每加練一次，一位社員的戰力加 $1$。 誰上場、誰打誰都可
以自己安排。
\item 求讓 $m$ 場全贏的最少加練次數。 $m \le n \le 2 \times 10^5$，戰力在
$0$ 到 $10^9$ 之間。
\end{itemize}

例如社員的戰力是 $4, 9, 1, 6$，對手是 $7, 3, 8$。 派 $4, 6, 9$ 上場，
$4$ 打 $3$、$6$ 打 $7$、$9$ 打 $8$，只有 $6$ 要加練 $2$ 次變成 $8$，答案
是 $2$。 如果讓 $9$ 打 $7$，$6$ 就得打 $8$，要加練 $3$ 次。

\textbf{想法}

先把一場比賽的成本寫清楚：戰力 $a$ 的社員要贏戰力 $b$ 的對手，得變成至少
$b + 1$，所以要加練 $\max(0, b + 1 - a)$ 次。 總次數是每一場的成本加起來，
問題變成怎麼挑人、怎麼配對，讓總和最小。

最直接的做法是枚舉：挑哪 $m$ 個人上場、上場的人怎麼排。 只看 $m = n$ 的情
況就有 $n!$ 種排法，$n = 10$ 已經三百多萬種，$n$ 到 $2 \times 10^5$ 完全不
可能。 不過這個枚舉仍然有用：$n \le 8$ 時用 \texttt{next\_permutation} 枚
舉所有排列，就是對拍用的暴力解。

手算幾個例子，可以看出兩個規律。

\textbf{觀察一：上場的是最強的 $m$ 位。} 假設在某個安排裡，戰力 $a$ 的社員
上場，戰力 $a' \ge a$ 的社員卻坐在場邊。 讓場邊的人換上去打同一位對手，成本
從 $\max(0, b + 1 - a)$ 變成 $\max(0, b + 1 - a')$，不會變多，其他場也不
受影響。 一直換下去，上場的就是最強的 $m$ 位。

\textbf{觀察二：依名次配對。} 把上場的 $m$ 位和對手都由弱到強排好，第 $i$
弱的社員打第 $i$ 弱的對手。 用交換論證說明。 假設某個安排裡有一組
\textbf{交叉}：社員 $x \le y$，對手 $p \le q$，卻是 $x$ 打 $q$、$y$ 打
$p$。 把兩人的對手互換，讓 $x$ 打 $p$、$y$ 打 $q$，要說明成本不會變多。

四種組合各自「還差多少」是
\[
A = q + 1 - x,\qquad D = p + 1 - y,\qquad B = p + 1 - x,\qquad C = q + 1 - y.
\]
交叉時用到 $A$ 和 $D$，互換後用到 $B$ 和 $C$。 因為 $x \le y$、
$p \le q$，$A$ 是四個裡最大的，$D$ 是最小的，而且 $B + C = A + D$（兩邊都
是 $p + q + 2 - x - y$）。 成本是每個數先取 $\max(0, \cdot)$ 再相加，分三
種情況：

\begin{itemize}
\item $D \ge 0$：四個數都不小於 $0$，互換前後的成本都是 $A + D = B + C$，
一樣多。
\item $A \le 0$：四個數都不大於 $0$，前後都是 $0$。
\item $D < 0 < A$：交叉時的成本是 $A$。 互換之後，如果 $B$、$C$ 都是正的，
成本是 $B + C = A + D < A$；只有一個是正的，那一個不超過 $A$；都不是正的，
成本是 $0$。 三種都不超過 $A$。
\end{itemize}

所以把一組交叉換開，總成本不會增加。 只要還不是依名次配對，就一定找得到一
組交叉；每換開一組，交叉的組數就會減少，最後一定變成依名次配對。 和觀察一
合起來：挑最強的 $m$ 位，依名次配對，就是最少的加練次數。

\textbf{程式}

\begin{example}[label=ex:training]{cpp}{(training.cpp) 全勝所需的最少加練次數}
#include <algorithm>
#include <iostream>
#include <vector>
using namespace std;

int main() {
    ios::sync_with_stdio(false);
    cin.tie(nullptr);
    int n, m;
    cin >> n >> m;
    vector<long long> a(n), b(m);
    for (long long &x : a) {
        cin >> x;
    }
    for (long long &x : b) {
        cin >> x;
    }
    sort(a.begin(), a.end());
    sort(b.begin(), b.end());
    long long total = 0;
    for (int i = 0; i < m; i++) {
        // 上場的是最強的 m 位：a[n - m] 到 a[n - 1]
        long long need = b[i] + 1 - a[n - m + i];   // 要比對手多 1
        if (need > 0) {
            total += need;
        }
    }
    cout << total << '\n';
    return 0;
}
\end{example}

\iosample{
4 3 \\
4 9 1 6 \\
7 3 8 \\
}{
2 \\
}

\begin{itemize}
\item 兩個陣列都由小到大排序後，最強的 $m$ 位社員是 \texttt{a[n - m]} 到
\texttt{a[n - 1]}，其中第 $i$ 弱的是 \texttt{a[n - m + i]}，對上第 $i$ 弱
的對手 \texttt{b[i]}。
\item \texttt{need} 不是正的，表示本來就會贏，不用加練。 這兩行也可以寫成
\texttt{total += max(0LL, need);}。 寫成 \texttt{max(0, need)} 會編譯失敗，
因為 \texttt{0} 是 \texttt{int}、\texttt{need} 是 \texttt{long long}，
\texttt{max} 的兩個引數必須同型別（銜接書第~8~章的函式模板）。
\item 戰力存成 \texttt{long long}，\texttt{b[i] + 1 - a[...]} 就直接用
\texttt{long long} 算。 總和最大大約是 $2 \times 10^5 \times 10^9 =
2 \times 10^{14}$，一定要用 \texttt{long long}。
\end{itemize}

寫完之後，用前面說的暴力解（枚舉排列，前 $m$ 個人依序打 $m$ 位對手）在
$n \le 8$ 的隨機資料上比對，可以再多一層確認。

\textbf{複雜度}

兩次排序是 $O(n \log n + m \log m)$，配對一次迴圈 $O(m)$。 $n = 2 \times
10^5$ 時排序大約 $4 \times 10^6$ 次比較，最大的測資在本書測試的環境下約
$0.03$ 秒。

\subsection{常見錯誤}

\begin{itemize}
\item \textbf{先湊最多勝場，再補練}：一個很自然的想法是先不加練，盡量多贏幾
場，剩下贏不了的對手再找人加練。 例如對手由弱到強，每位都派「贏得了他的人裡
最弱的那位」，派不出人的對手留到最後，由剩下的社員加練去打。 輸入
\begin{code}{text}{這種做法會錯的輸入}
3 2
3 0 0
2 1
\end{code}
它先用唯一贏得了的 $3$ 打掉 $1$，剩下 $0$ 要打 $2$，加練 $3$ 次，程式印出
\texttt{3}。 正確答案是 $2$：讓 $3$ 打 $2$，一位 $0$ 加練 $2$ 次去打 $1$。
「勝場最多」和「加練最少」是兩個不同的目標，前一個的最佳安排，不一定是後一
個的最佳安排。
\item \textbf{派最弱的 $m$ 位上場}：把 \texttt{a[n - m + i]} 寫成
\texttt{a[i]}，範例的輸入會派 $1, 4, 6$ 上場，印出 \texttt{10}。 排序之後
最弱的在前面，要從後面數 $m$ 個。
\item \textbf{平手當成贏}：成本寫成 \texttt{b[i] - a[n - m + i]}，少了
$+1$，範例的輸入印出 \texttt{1}。 $6$ 打 $7$ 只加練到 $7$，是平手，不算贏。
\item \textbf{總和用 \texttt{int}}：\texttt{total} 改成 \texttt{int}，輸入
$n = m = 2 \times 10^5$，社員全是 $0$、對手全是 $10^9$，正確答案是
$200000200000$。 在本書測試的環境下程式印出 \texttt{553094464}，加了
\texttt{-Wall -Wextra} 也沒有任何警告。
\end{itemize}

\section{區間排程}
\label{sec:interval-scheduling}

很多題目的資料是一段一段的時間：借教室、排值班、安排表演。 一段時間用
\textbf{區間}表示。 這一節用半開區間 $[s, e)$：從時刻 $s$ 開始，到了 $e$
那一刻已經離開，所以 $[2, 4)$ 和 $[4, 6)$ 剛好接在一起，不算重疊。 兩個區
間 $[a, b)$ 和 $[c, d)$ \textbf{重疊}，若且唯若 $a < d$ 而且 $c < b$，也就
是兩段都在對方結束之前就開始了。

最常見的問題是：給一堆區間，最多能挑出幾個兩兩不重疊的？ 這叫\textbf{區間排
程}（interval scheduling）。 能想到的挑法有好幾種，聽起來都有道理。 這一節
的例題要你把其中三種都做出來，正好拿來比較。

\subsection{例題：三種挑選規則的比較}
\label{sec:ex-three-rules}

\textbf{題意}（\akiref{comparing-three-selection-rules}{三種挑選規則的比較}）

\begin{itemize}
\item 社團辦公室收到 $n$ 筆借用申請，第 $i$ 筆要用 $[s_i, e_i)$。
\item 三種規則都是先把申請排成某個次序，再依序檢查：和已經選到的每一筆都不重
疊就選，否則跳過。
\item 規則 1：開始早的在前；開始相同時結束早的在前；再相同就編號小的在前。
\item 規則 2：長度 $e - s$ 短的在前；長度相同時開始早的在前；再相同就編號小
的在前。
\item 規則 3：結束早的在前；結束相同時開始\textbf{晚}的在前；再相同就編號小
的在前。
\item 輸出三種規則各選到幾筆。 $n \le 2000$，時刻在 $0$ 到 $10^9$ 之間。
\end{itemize}

例如四筆申請 $[0, 20)$、$[1, 5)$、$[4, 7)$、$[6, 11)$。 規則 1 先選
$[0, 20)$，後面每一筆都和它重疊，只選到 $1$ 筆。 規則 2 先選最短的
$[4, 7)$，它和 $[1, 5)$、$[6, 11)$ 都重疊，也只選到 $1$ 筆。 規則 3 選
$[1, 5)$ 和 $[6, 11)$，選到 $2$ 筆。

\textbf{想法}

三條規則的模擬照題目寫就好，要想的只有「和已經選到的每一筆都不重疊」怎麼檢
查。 最直接的是把選到的申請存在一個 \texttt{vector} 裡，每來一筆就和它們全
部比一次。 最壞的情況是每一筆都被選到，第 $i$ 筆要比 $i - 1$ 次，總共
$n(n - 1)/2 \approx 2 \times 10^6$ 次，$n \le 2000$ 綽綽有餘。 三條規則共
用同一個檢查函式，只差在排序的方式。

規則 1 和規則 3 其實不用和全部比。 以規則 3 為例，申請依結束時間排序，已經
選到的每一筆都比目前這筆早結束（或同時），其中最後選到的那一筆結束得最晚；
目前這筆只要不早於它的結束時間開始，就和所有選到的都不重疊。 規則 2 就不能
這樣省：依長度排序時，選到的申請散在時間軸上各處，最後選的那一筆不一定靠
近目前這一筆（見本節的常見錯誤）。 $n$ 只有 $2000$，統一和全部比對，最不
容易寫錯。

程式寫得出來，但這題真正要學的是：哪一條規則能保證選到最多？

\textbf{規則 1 會錯。} 反例：$[0, 10)$、$[1, 2)$、$[3, 4)$。 規則 1 先選
$[0, 10)$，另外兩筆都被擋掉，選到 $1$ 筆；不選它，可以選 $[1, 2)$ 和
$[3, 4)$ 兩筆。 早開始的申請可能很長，一筆就佔掉一大段時間。

\textbf{規則 2 也會錯。} 反例：$[0, 4)$、$[3, 6)$、$[5, 9)$。 最短的
$[3, 6)$ 剛好跨在另外兩筆的交界上，選了它就同時擋掉兩筆，只選到 $1$ 筆；最
多其實可以選 $[0, 4)$ 和 $[5, 9)$ 兩筆。 短的申請佔的時間少，但位置不好時
照樣擋路。

\textbf{規則 3 是對的。} 用交換論證。 設最多能選 $k$ 筆，任取一個選到 $k$
筆的選法，依結束時間排好是 $o_1, o_2, \ldots, o_k$。 規則 3 選的第一筆
$g_1$ 是所有申請裡結束最早的，所以 $g_1$ 的結束時間不晚於 $o_1$ 的結束時
間。 把 $o_1$ 換成 $g_1$：$o_2, \ldots, o_k$ 都在 $o_1$ 結束之後才開始，當
然也在 $g_1$ 結束之後才開始，不會和 $g_1$ 重疊。 換完還是 $k$ 筆、兩兩不重
疊，而且第一筆和規則 3 一樣。

接下來，規則 3 跳過所有和 $g_1$ 重疊的申請，在「開始時間不早於 $g_1$ 結束」
的申請裡選結束最早的，當作第二筆。 換過的選法剩下的 $o_2, \ldots, o_k$ 也
都在這些申請裡，同樣的理由可以把 $o_2$ 換成規則 3 的第二筆。 一直換下去，
$k$ 筆都換成規則 3 選的，所以規則 3 至少選到 $k$ 筆，也就是最多。

把三條規則放在一起看，差別在於選了一筆之後「留給後面多少時間」。 規則 3 每
次都讓剩下的時間從最早的地方開始，前兩條規則都沒有這個保證。

\begin{bridgerecap}[多個欄位的比較函式]
\texttt{sort} 的第三個引數是比較函式，收兩個元素 \texttt{x}、\texttt{y}，
當 \texttt{x} 應該排在 \texttt{y} 前面時回傳 \texttt{true}（銜接書第~9~章）。
比好幾個欄位時，第一個欄位不同就比第一個欄位，相同才往下比。 比較函式必須
是嚴格的「小於」：兩個元素相等時要回傳 \texttt{false}，所以不能用
\texttt{<=}。 \texttt{sort} 不保證相等的元素維持原來的順序；要讓結果只有一
種，就要比到每兩個元素都分得出先後，例如最後比編號。
\end{bridgerecap}

\textbf{程式}

\begin{example}[label=ex:three-rules]{cpp}{(three\_rules.cpp) 三種挑選規則的比較}
#include <algorithm>
#include <iostream>
#include <vector>
using namespace std;

struct Req {
    long long s, e;   // 時段 [s, e)
    int id;
};

bool overlap(const Req &x, const Req &y) {
    return x.s < y.e && y.s < x.e;
}

bool byStart(const Req &x, const Req &y) {      // 規則 1
    if (x.s != y.s) {
        return x.s < y.s;
    }
    if (x.e != y.e) {
        return x.e < y.e;
    }
    return x.id < y.id;
}

bool byLength(const Req &x, const Req &y) {     // 規則 2
    long long lx = x.e - x.s, ly = y.e - y.s;
    if (lx != ly) {
        return lx < ly;
    }
    if (x.s != y.s) {
        return x.s < y.s;
    }
    return x.id < y.id;
}

bool byEnd(const Req &x, const Req &y) {        // 規則 3
    if (x.e != y.e) {
        return x.e < y.e;
    }
    if (x.s != y.s) {
        return x.s > y.s;                       // 開始晚的在前
    }
    return x.id < y.id;
}

// v 已經照某個規則排好：依序檢查，和已選的每一筆都不重疊才選
int pickInOrder(const vector<Req> &v) {
    vector<Req> chosen;
    for (const Req &r : v) {
        bool ok = true;
        for (const Req &c : chosen) {
            if (overlap(r, c)) {
                ok = false;
                break;
            }
        }
        if (ok) {
            chosen.push_back(r);
        }
    }
    return chosen.size();
}

int main() {
    ios::sync_with_stdio(false);
    cin.tie(nullptr);
    int n;
    cin >> n;
    vector<Req> v(n);
    for (int i = 0; i < n; i++) {
        cin >> v[i].s >> v[i].e;
        v[i].id = i;
    }
    sort(v.begin(), v.end(), byStart);
    int r1 = pickInOrder(v);
    sort(v.begin(), v.end(), byLength);
    int r2 = pickInOrder(v);
    sort(v.begin(), v.end(), byEnd);
    int r3 = pickInOrder(v);
    cout << r1 << ' ' << r2 << ' ' << r3 << '\n';
    return 0;
}
\end{example}

\iosample{
4 \\
0 20 \\
1 5 \\
4 7 \\
6 11 \\
}{
1 1 2 \\
}

\begin{itemize}
\item \texttt{overlap} 直接照重疊的定義寫。 用 \texttt{<} 而不是
\texttt{<=}，端點相接就不算重疊。
\item 三個比較函式最後都比 \texttt{id}，任兩筆申請都分得出先後，排序的結果
只有一種。 所以同一個 \texttt{v} 可以連續重排三次，前一次的順序不會影響後
一次。
\item \texttt{pickInOrder} 收的是 \texttt{const} 參考，只讀不改，也不複製。
它回傳 \texttt{chosen.size()}，型別是 \texttt{size\_t}，回傳時轉成
\texttt{int}；筆數最多 $2000$，不會有問題。
\end{itemize}

\textbf{複雜度}

三次排序各 $O(n \log n)$，三次 \texttt{pickInOrder} 各最多 $n(n-1)/2$ 次
\texttt{overlap}，總共 $O(n^2)$。 $n = 2000$ 時大約 $6 \times 10^6$ 次比
較，最大的測資在本書測試的環境下不到 $0.01$ 秒。

\subsection{常見錯誤}

\begin{itemize}
\item \textbf{規則 2 只和最後選到的那一筆比}：把規則 1、3 的省略用在規則 2。
輸入
\begin{code}{text}{規則 2 只比最後一筆會錯的輸入}
3
0 2
5 8
1 5
\end{code}
正確的輸出是 \texttt{2 2 2}，這樣寫會印出 \texttt{2 3 2}。 規則 2 依序看
$[0, 2)$、$[5, 8)$、$[1, 5)$：前兩筆都選；$[1, 5)$ 和最後選的 $[5, 8)$ 只
是相接，於是被選了，可是它和更早選的 $[0, 2)$ 重疊。
\item \textbf{相接當成重疊}：\texttt{overlap} 寫成 \texttt{x.s <= y.e \&\&
y.s <= x.e}。 輸入 $[2, 4)$、$[0, 7)$、$[4, 6)$、$[3, 5)$、$[6, 9)$
（第一行是 \texttt{5}）時，正確的輸出是 \texttt{1 3 3}，這樣寫會印出
\texttt{1 2 2}：規則 3 選了 $[2, 4)$ 之後，$[4, 6)$ 被當成重疊而跳過。
讀題時先確認區間是半開還是兩端都算，再決定用 \texttt{<} 還是 \texttt{<=}。
\item \textbf{同分的次序沒照題目}：規則 1 只比開始時間，同開始的照輸入順序
（例如用 \texttt{stable\_sort}）。 輸入 $[0, 5)$、$[0, 1)$、$[1, 2)$ 時，
規則 1 應該先看 $[0, 1)$，選到 $2$ 筆；照輸入順序會先選到很長的 $[0, 5)$，
只選到 $1$ 筆，程式印出 \texttt{1 2 2}。 規則 3 的同分次序不影響筆數（交換
論證裡沒有用到它），規則 1、2 的會，題目怎麼規定就怎麼寫。
\end{itemize}

\section{區間覆蓋與端點掃描}
\label{sec:cover-sweep}

上一節從區間裡挑出互不重疊的。 這一節看另外兩種常見的區間題：用最少的區間
把一整段蓋滿，以及把區間拆成端點，排序後由左往右掃一遍。 兩種都是「排好
順序、依序處理、處理過的不再回頭」，和前兩節是同一種寫法。

\subsection{區間覆蓋}

問題：從時刻 $x$ 到 $y$ 一直都要有人值班。 有 $n$ 位志工，第 $i$ 位能待的
時段是 $[s_i, t_i)$。 最少要找幾位，才能把 $[x, y)$ 整段蓋滿，中間沒有空
檔？ 下一位只要在前一位離開的那一刻或更早到，就接得上。

貪心的做法：記住目前已經顧到哪個時刻 \texttt{reach}，一開始是 $x$。 在所有
「開始時刻 $\le$ \texttt{reach}」的志工裡，選離開最晚的那一位，把
\texttt{reach} 推到他離開的時刻，重複到 \texttt{reach} $\ge y$。 如果接得上
的志工裡沒有人能把 \texttt{reach} 往後推，就是蓋不滿。

這次用另一種方式說明它是對的：證明貪心一路領先。 設貪心選了 $k$ 位之後顧到
$R_k$。 任取一種接班的方式，它用了 $k$ 位之後顧到 $r_k$。 要說明對每個 $k$
都有 $r_k \le R_k$。

\begin{itemize}
\item $k = 1$：兩邊的第一位都必須在 $x$ 或更早到。 貪心在這些人裡選離開最晚
的，所以 $r_1 \le R_1$。
\item 假設 $r_k \le R_k$。 另一種方式的第 $k + 1$ 位必須在 $r_k$ 或更早到，
所以也在 $R_k$ 或更早到，是貪心第 $k + 1$ 步會看到的人。 貪心在這些人裡選離
開最晚的，所以 $r_{k+1} \le R_{k+1}$。
\end{itemize}

如果某種方式用 $k$ 位就蓋滿了，$y \le r_k \le R_k$，貪心用 $k$ 位也蓋滿了，
所以貪心用的人數最少。 跟交換論證比起來，這種寫法不改動最佳解，只一步一步
比較兩邊各走到哪裡，適合「每一步都在推進某個量」的題目。

\begin{example}[label=ex:cover]{cpp}{(cover.cpp) 用最少的志工蓋滿一段時間}
#include <algorithm>
#include <iostream>
#include <utility>
#include <vector>
using namespace std;

int main() {
    int n;
    long long x, y;
    cin >> n >> x >> y;
    vector<pair<long long, long long>> iv(n);   // [s, t)
    for (auto &[s, t] : iv) {
        cin >> s >> t;
    }
    sort(iv.begin(), iv.end());                 // 依開始時刻排序
    long long reach = x;                        // 目前顧到哪裡
    int count = 0, i = 0;
    while (reach < y) {
        long long far = reach;
        while (i < n && iv[i].first <= reach) { // 接得上的都看一遍
            far = max(far, iv[i].second);
            i++;
        }
        if (far == reach) {                     // 沒有人能往後推
            count = -1;
            break;
        }
        reach = far;                            // 選離開最晚的那一位
        count++;
    }
    cout << count << '\n';
    return 0;
}
\end{example}

\iosample{
7 0 10 \\
0 3 \\
1 5 \\
2 4 \\
4 8 \\
5 6 \\
7 10 \\
6 9 \\
}{
4 \\
}

範例要蓋滿 $[0, 10)$。 一開始接得上的只有 $[0, 3)$，顧到 $3$；接著
$[1, 5)$、$[2, 4)$ 接得上，選 $[1, 5)$ 顧到 $5$；再來 $[4, 8)$、$[5, 6)$
裡選 $[4, 8)$ 顧到 $8$；最後 $[6, 9)$、$[7, 10)$ 裡選 $[7, 10)$，共 $4$ 位。

\begin{itemize}
\item \texttt{for (auto \&[s, t] : iv)} 是結構化綁定加參考（銜接書第~8~章），
\texttt{s}、\texttt{t} 直接是 \texttt{iv} 裡每個 \texttt{pair} 的兩個成
員，讀進來就存好了。
\item \texttt{i} 只往前走、從不歸零。 已經看過的志工開始時刻都 $\le$ 目前的
\texttt{reach}，以後的 \texttt{reach} 只會更大，他們不會再有用，最好的那位
已經記在 \texttt{far} 裡。 所以內層迴圈加起來最多跑 $n$ 次，整支程式是排序
的 $O(n \log n)$。
\item \texttt{far == reach} 表示接得上的人都沒辦法往後推，這時一定要停下來。
\end{itemize}

\subsection{例題：全程出席的會議時段}
\label{sec:ex-meeting}

\textbf{題意}（\akiref{a-meeting-slot-nobody-leaves}{全程出席的會議時段}）

\begin{itemize}
\item 社團要開一場連續 $D$ 分鐘的線上討論會。 從第 $x$ 分鐘開始的話，會佔用
第 $x, x + 1, \ldots, x + D - 1$ 分鐘。
\item 有 $n$ 個小組，第 $i$ 組有 $m_i$ 人，在第 $s_i$ 到第 $t_i$ 分鐘（兩端
都算）有空。 討論會的每一分鐘都有空，這一組才會全組參加；只要有一分鐘沒空，
整組都不來。
\item 求參加人數最多是多少，以及達到最多人數的最早的 $x$。 沒有任何一組能
來時輸出 \texttt{0 0}。
\item $n \le 2 \times 10^5$，$D$ 和時刻都到 $10^9$，$m_i \le 10^5$。
\end{itemize}

這題的時間是一分鐘一格，兩端都算，和上一小節的半開區間不同，下面的式子都照
這個規定寫。

例如 $D = 3$，四組是（$5$ 人，第 $1$～$5$ 分鐘）、（$2$ 人，第 $4$～$9$ 分
鐘）、（$4$ 人，第 $2$～$6$ 分鐘）、（$3$ 人，第 $6$～$7$ 分鐘）。 從第 $2$
分鐘開始，第一組和第三組都能來，共 $9$ 人；第 $3$ 分鐘開始也是 $9$ 人，但要
輸出比較早的 $2$。

\textbf{想法}

最直接的做法是枚舉 $x$，每個 $x$ 檢查每一組來不來。 $x$ 可能到 $10^9$，不
能每一分鐘都試。 只試每組的 $s_i$ 當作候選，有 $n$ 個候選、每個檢查 $n$ 組，
是 $O(n^2)$，$n = 2 \times 10^5$ 時是 $4 \times 10^{10}$，來不及。

\textbf{第一步：換成「可以從哪裡開始」。} 第 $i$ 組會來，等於討論會的第一分
鐘和最後一分鐘都在他們的空檔裡：$s_i \le x$ 而且 $x + D - 1 \le t_i$。 也就
是說，$x$ 落在區間 $[s_i,\ t_i - D + 1]$ 裡（兩端都算）。 如果
$t_i - D + 1 < s_i$，空檔不到 $D$ 分鐘，這一組怎樣都來不了，直接丟掉。

換過之後，問題變成：數線上有一些區間，每個區間有一個人數，找一個整數點 $x$，
讓蓋住它的區間人數加起來最大。 注意不能直接拿原本的空閒時間 $[s_i, t_i]$
來算重疊：上面的例子裡，第 $4$ 分鐘前三組都有空，合計 $11$ 人，可是從第 $4$
分鐘開始的討論會要開到第 $6$ 分鐘，第一組第 $6$ 分鐘已經沒空了。

\textbf{第二步：拆成端點，排序後掃描。} 每個區間 $[s_i,\ t_i - D + 1]$ 拆成
兩個\textbf{事件}：在時刻 $s_i$ 人數加 $m_i$；在時刻 $t_i - D + 2$（第一個不
能開始的時刻）人數減 $m_i$。 把所有事件依時刻排序，由左往右把變化量累加起來，
累加到某個事件時的總和，就是「從這個時刻開始開會，會來幾個人」，一直維持到
下一個事件之前。 最大值一定出現在某個事件的時刻上，所以只要看這 $2n$ 個點。

同一個時刻同時有加有減時，要\textbf{先減後加}。 在時刻 $T$ 減掉的組，最晚只
能從 $T - 1$ 開始，在 $T$ 已經不能來了；先加後減的話，會有一瞬間把兩邊的人
都算進去。 \texttt{pair} 的預設比較先比 \texttt{first}，相同再比
\texttt{second}，負數排在正數前面，剛好就是先減後加。

題目要最早的 $x$。 由左往右掃描，總和\textbf{嚴格大於}目前的最大值時才更新，
記下的就是第一次達到最大值的時刻。

用上面的例子跑一次。 第四組 $7 - 3 + 1 = 5 < 6$，丟掉；其他三組的可開始區
間是 $[1, 3]$、$[4, 7]$、$[2, 4]$。 事件排好之後：

\begin{center}
\begin{tabular}{cccccc}
\toprule
時刻 & 1 & 2 & 4 & 4 & 5 \\
變化 & $+5$ & $+4$ & $-5$ & $+2$ & $-4$ \\
總和 & 5 & 9 & 4 & 6 & 2 \\
\bottomrule
\end{tabular}
\end{center}

最後還有時刻 $8$ 的 $-2$，總和回到 $0$。 最大是 $9$，第一次出現在時刻 $2$。

這題放在貪心這一章，是因為寫法和前面一樣：把東西排成一個順序，由左往右處理，
處理過的不再回頭。 第~5~章的差分是同一個想法：在區間的開頭加、在結尾的下
一格減。

\textbf{程式}

\begin{example}[label=ex:meeting]{cpp}{(meeting.cpp) 全程出席的會議時段}
#include <algorithm>
#include <iostream>
#include <utility>
#include <vector>
using namespace std;

int main() {
    ios::sync_with_stdio(false);
    cin.tie(nullptr);
    int n;
    long long d;
    cin >> n >> d;
    vector<pair<long long, long long>> ev;   // (時刻, 人數的變化)
    for (int i = 0; i < n; i++) {
        long long m, s, t;
        cin >> m >> s >> t;
        long long last = t - d + 1;          // 最晚可以從 last 開始
        if (last < s) {
            continue;                        // 空檔不到 d 分鐘
        }
        ev.push_back({s, m});                // 從 s 起可以開始
        ev.push_back({last + 1, -m});        // 從 last + 1 起不行了
    }
    sort(ev.begin(), ev.end());              // 同時刻負的在前
    long long cur = 0, best = 0, bestX = 0;
    for (auto [at, delta] : ev) {
        cur += delta;
        if (cur > best) {                    // 嚴格變大才更新
            best = cur;
            bestX = at;
        }
    }
    cout << best << ' ' << bestX << '\n';
    return 0;
}
\end{example}

\iosample{
4 3 \\
5 1 5 \\
2 4 9 \\
4 2 6 \\
3 6 7 \\
}{
9 2 \\
}

\begin{itemize}
\item \texttt{best} 和 \texttt{bestX} 一開始都是 $0$。 沒有任何一組能來時，
迴圈裡不會更新，剛好印出題目要的 \texttt{0 0}。
\item 人數的總和最多是 $2 \times 10^5 \times 10^5 = 2 \times 10^{10}$，超過
\texttt{int}，所以變化量和總和都用 \texttt{long long}。
\item \texttt{t - d + 1} 可能是負數（$d$ 比 $t$ 大很多時），用
\texttt{long long} 算不會出問題，而且負數一定小於 $s$，這組會被丟掉。
\end{itemize}

\textbf{複雜度}

$2n$ 個事件排序是 $O(n \log n)$，掃描 $O(n)$。 最大的測資在本書測試的環境
下約 $0.05$ 秒。

\subsection{常見錯誤}

\begin{itemize}
\item \textbf{沒有把空檔縮短}：直接用 $[s_i, t_i]$ 拆事件，找最多組同時有
空的時刻。 範例的輸入會印出 \texttt{11 4}，就是上面說的第 $4$ 分鐘。
\item \textbf{可開始區間少算一格}：最晚的開始時刻寫成 \texttt{t - d}，把兩
端都算的空檔當成了半開區間，剛好塞得下的組被漏掉。 輸入
\begin{code}{text}{少算一格會錯的輸入}
4 4
3 0 3
6 10 12
2 2 9
2 6 10
\end{code}
正確答案是 \texttt{4 6}：從第 $6$ 分鐘開始開到第 $9$ 分鐘，第三組的空檔剛好
到第 $9$ 分鐘。 少算一格會印出 \texttt{2 2}。
\item \textbf{用 \texttt{>=} 更新答案}：輸入 \texttt{2 1}、\texttt{3 0 0}、
\texttt{3 5 5}（三行）時，第 $0$ 分鐘和第 $5$ 分鐘都是 $3$ 人，應該輸出
\texttt{3 0}。 用 \texttt{>=} 時，後面相等的時刻會蓋掉前面的，印出
\texttt{3 5}。
\item \textbf{同一時刻先加後減}：自己寫比較函式，同時刻時把正的排前面。 輸入
\texttt{2 1}、\texttt{2 0 1}、\texttt{3 2 3} 時，第一組最晚從 $1$ 開始，第
二組最早從 $2$ 開始，兩組不可能都來，答案是 \texttt{3 2}。 先加後減會在時刻
$2$ 短暫算出 $5$ 人，印出 \texttt{5 2}。
\item \textbf{人數用 \texttt{int}}：輸入 $2 \times 10^5$ 組、每組 $10^5$ 人、
都在第 $0$～$10$ 分鐘有空，$D = 1$，正確答案是 \texttt{20000000000 0}。 把
人數相關的變數都改成 \texttt{int}，在本書測試的環境下印出
\texttt{2147465408 0}。
\end{itemize}

\section{本章常見錯誤}
\label{sec:greedy-pitfalls}

前面各節的錯誤，歸納起來是四類。 寫貪心的題目時，交出去之前逐項檢查一次。

\begin{center}
\begin{tabular}{p{0.2\linewidth}p{0.44\linewidth}p{0.24\linewidth}}
\toprule
錯誤 & 例子 & 見 \\
\midrule
沒有理由就貪 & 郵票先貼大面額；加練題先湊最多勝場 &
第~\ref{sec:greedy-idea}、\ref{sec:ex-training}~節 \\
排序的鍵選錯 & 區間排程依開始時間或長度排序 & 第~\ref{sec:ex-three-rules}~節 \\
相等時的順序 & 規則 1 同開始時的次序；掃描時同一時刻先加後減；用
\texttt{>=} 更新答案 & 第~\ref{sec:ex-three-rules}、\ref{sec:ex-meeting}~節 \\
端點的開閉 & 半開區間用了 \texttt{<=}；兩端都算的空檔算成 \texttt{t - d} &
第~\ref{sec:ex-three-rules}、\ref{sec:ex-meeting}~節 \\
\bottomrule
\end{tabular}
\end{center}

另外兩個和貪心無關、但在這一章每一題都出現的錯誤：總和超過 \texttt{int}，
以及只用範例測資測試。 前者在動手前先估答案最大多少，後者用暴力解在小範圍對
拍。

\section{本章練習}
\label{sec:ch7-exercises}

\begin{drills}[label=drills:ch7]
\item 校外教學要划雙人獨木舟。 每艘最多坐兩人，兩人的體重加起來不能超過
$W$；每個人的體重都不超過 $W$。 求最少要幾艘。 下面兩個函式都先把體重由小到
大排序：
\begin{code}{cpp}{練習 1 的兩種寫法}
int boatsA(vector<int> w, int cap) {
    sort(w.begin(), w.end());
    int i = 0, j = (int)w.size() - 1, boats = 0;
    while (i <= j) {
        if (i < j && w[i] + w[j] <= cap) {
            i++;                    // 最輕的跟最重的同船
        }
        j--;                        // 最重的一定上船
        boats++;
    }
    return boats;
}

int boatsB(vector<int> w, int cap) {
    sort(w.begin(), w.end());
    int n = w.size(), i = 0, boats = 0;
    while (i < n) {
        if (i + 1 < n && w[i] + w[i + 1] <= cap) {
            i += 2;                 // 相鄰兩人同船
        } else {
            i += 1;
        }
        boats++;
    }
    return boats;
}
\end{code}
其中一個是對的，另一個會錯。 找出錯的那一個並給一組反例；對的那一個，用交換
論證說明為什麼對。

\item 下面每一條貪心規則，判斷它是否一定能得到最佳解。 會錯的，給一組反例。
\begin{enumerate}
\item 郵票只有 $1$ 元、$7$ 元、$10$ 元，要剛好貼滿 $n$ 元且張數最少，每次先
貼面額最大、還貼得下的。（找出最小的反例金額。）
\item 背包最多裝重量 $W$，每件物品有重量和價值，每件最多拿一次，要總價值最大。
依「價值除以重量」由大到小，放得下就放。
\item 區間排程（第~\ref{sec:interval-scheduling}~節），依開始時間由晚到早排序，
和已經選到的都不重疊就選。
\item 第~\ref{sec:ex-training}~節的加練問題，對手由強到弱，每位都派剩下的社
員裡最強的那一位去打。
\end{enumerate}

\item 估計下列程式的執行次數或時間複雜度。
\begin{enumerate}
\item 範例~\ref{ex:cover} 的內層 \texttt{while} 迴圈（\texttt{far = max(...)}
那一行）在整支程式裡最多執行幾次？ 外層迴圈呢？
\item 範例~\ref{ex:three-rules} 的 $n = 2000$ 時，一次 \texttt{pickInOrder}
最多呼叫 \texttt{overlap} 幾次？ 什麼樣的輸入會達到這個次數？
\item 下面的程式片段，\texttt{a} 已經由小到大排好。 最壞情況下
\texttt{j++} 執行幾次？ 和 (a) 比較，差在哪裡？
\begin{code}{cpp}{練習 3(c)}
for (int i = 0; i < n; i++) {
    int j = i;
    while (j < n && a[j] <= a[i] + k) {
        j++;
    }
    best = max(best, j - i);
}
\end{code}
\end{enumerate}

\item 有八筆申請：\#1 $[1, 4)$、\#2 $[3, 5)$、\#3 $[0, 6)$、\#4 $[5, 7)$、
\#5 $[3, 9)$、\#6 $[5, 9)$、\#7 $[6, 10)$、\#8 $[8, 11)$。
\begin{enumerate}
\item 寫出規則 3 的檢查順序，以及依序選到哪幾筆。
\item 範例~\ref{ex:three-rules} 對這組輸入會印出什麼？
\end{enumerate}

\item 有人把範例~\ref{ex:meeting} 的排序改成只比時刻：
\begin{code}{cpp}{練習 5}
sort(ev.begin(), ev.end(),
     [](const pair<long long, long long> &p,
        const pair<long long, long long> &q) {
         return p.first < q.first;
     });
\end{code}
範例的輸入仍然印出 \texttt{9 2}。 這樣改有什麼問題？ 給一組可能會錯的輸入。

\item 有 $n$ 位同學排隊問助教問題，第 $i$ 位要花 $t_i$ 分鐘。 每位同學的等待
時間是排在他前面的人花的時間總和。 證明依 $t_i$ 由小到大排隊，所有人的等待
時間總和最小。 提示：比較相鄰兩個人互換前後，總和差多少。
\end{drills}

\section{實作練習}
\label{sec:ch7-practice}

下面是 Aki-Judge 上和本章相關的題目，依小節分組。 難度分四階：暖身是直接套用
本章的寫法；基本要自己想出排序的鍵或證明；實戰的題敘比較長，要先把題目轉換成
本章的形狀；挑戰會和其他章的技巧合在一起。

\begin{center}
\begin{tabular}{llc}
\toprule
小節 & 題目 & 難度 \\
\midrule
排序與交換論證 & \akiref{full-crates-only}{整箱出貨} & 暖身 \\
 & \akiref{assigning-cleaning-zones}{打掃區域分配} & 暖身 \\
 & \akiref{rush-jobs-for-the-3d-printer}{3D 列印機的臨時插件} & 基本 \\
 & \akiref{ordering-multiply-add-and-add-multiply-cards}{乘加卡與加乘卡的出牌順序} & 基本 \\
 & \akiref{total-waiting-time-at-the-checkup-stations}{健康檢查站的總等待時間} & 基本 \\
 & \akiref{enhancement-stone-allocation}{強化石的使用次數} & 基本 \\
 & \akiref{which-report-goes-first}{哪一份可以先寫} & 實戰 \\
 & \akiref{the-shop-that-keeps-raising-prices}{越買越貴的商店} & 實戰 \\
 & \akiref{fewest-swaps-to-a-least-complaint-line}{把隊伍調成最少抱怨} & 挑戰 \\
\midrule
區間排程 & \akiref{walking-between-fair-booths}{園遊會攤位的移動時間} & 暖身 \\
 & \akiref{the-talk-you-must-attend}{非聽不可的那一場講座} & 基本 \\
 & \akiref{highlighting-keywords-in-notes}{筆記上的螢光筆標記} & 基本 \\
 & \akiref{earliest-time-to-collect-k-stamps}{集滿 K 枚印章的最早時刻} & 基本 \\
 & \akiref{overnight-duty-shifts}{跨夜的值班時段} & 實戰 \\
 & \akiref{most-on-target-periods}{切出最多個達標時段} & 實戰 \\
\midrule
區間覆蓋與掃描 & \akiref{fewest-recruitment-posts}{招生貼文的最少則數} & 暖身 \\
 & \akiref{saving-power-on-the-hallway-lights}{走廊感應燈省電開法} & 暖身 \\
 & \akiref{volunteers-for-the-tagged-shelf-spots}{處理所有待整理書位的志工} & 基本 \\
 & \akiref{relay-shifts-at-the-announcement-booth}{廣播台的值班接力} & 基本 \\
 & \akiref{fewest-trips-for-the-health-check-van}{健檢巡迴車的最少出車次數} & 實戰 \\
\bottomrule
\end{tabular}
\end{center}

本章的三題例題也在 Aki-Judge 上，讀完各節之後可以先自己寫一次再對照。

\addcontentsline{toc}{section}{練習解答}
\section*{練習解答}

\begin{enumerate}
\item \texttt{boatsB} 會錯。 $W = 10$，體重 $1, 2, 8, 9$：\texttt{boatsB}
讓 $1$ 和 $2$ 同船，$8$、$9$ 各自一艘，共 $3$ 艘；$1$ 配 $9$、$2$ 配 $8$ 只
要 $2$ 艘。 它把兩個輕的人放在一起，重的人就找不到伴了。

\texttt{boatsA} 是對的。 看目前最重的人 $h$ 和最輕的人 $\ell$。
\begin{itemize}
\item $h + \ell > W$：$h$ 和誰都不能同船，任何分法裡他都獨自一艘。
\texttt{boatsA} 也讓他獨自一艘。
\item $h + \ell \le W$：任取一個最佳的分法。 如果 $h$ 獨自一艘，把 $\ell$
移過來，艘數不會變多。 如果 $h$ 和 $z$ 同船、$\ell$ 獨自一艘，把 $\ell$ 和
$z$ 互換。 如果 $h$ 和 $z$ 同船、$\ell$ 和 $u$ 同船，也把 $\ell$ 和 $z$ 互
換：新的兩艘是 $h + \ell \le W$，以及 $z + u \le z + h \le W$（$u$ 不比
$h$ 重）。 換完艘數不變，而且 $h$ 和 $\ell$ 同船，和 \texttt{boatsA} 一樣。
\end{itemize}
兩種情況下，\texttt{boatsA} 對 $h$ 的處理都和某個最佳分法一致；剩下的人是同
樣的問題，同樣的理由一直用下去。

\item
\begin{enumerate}
\item 會錯，最小的反例是 $14$ 元：貪心貼 $10 + 1 + 1 + 1 + 1$ 共 $5$ 張，
兩張 $7$ 元就夠了。 把範例~\ref{ex:stamps} 的面額改成 $10, 7, 1$、範圍改到
$30$ 跑一次，第一行就是 \texttt{14 5 2}。
\item 會錯。 $W = 10$，三件物品（重量 $6$、價值 $7$）、（$5$, $5$）、
（$5$, $5$）。 第一件的價值除以重量最大，先放進去之後只剩 $4$，別的都放不下，
總價值 $7$；拿兩件重量 $5$ 的是 $10$。 這題要用第~9~章的動態規劃（銜接書
第~10~章練習 3 的 0/1 背包）。
\item 對。 把時間軸左右翻轉（每個時刻 $x$ 換成 $-x$），開始最晚的就變成結
束最早的，這條規則就變成規則 3；翻轉不會改變哪兩筆重疊，所以選到的筆數和規
則 3 一樣多。
\item 對。 最強的對手配最強的社員，第二強的配第二強的，一路下去，用到的剛好
是最強的 $m$ 位社員，而且依名次配對，和範例~\ref{ex:training} 的安排完全相
同，只是從強的那一端開始配。
\end{enumerate}

\item
\begin{enumerate}
\item \texttt{far = max(...)} 最多執行 $n$ 次。 每執行一次，\texttt{i} 就加
$1$，而 \texttt{i} 從 $0$ 開始、從不減少、不超過 $n$。 外層迴圈除了最後一
次以外，每一次都要看到至少一位新的志工才能把 \texttt{reach} 往後推（舊的志
工都已經算在之前的 \texttt{far} 裡），所以最多 $n + 1$ 次。 排序之後的部分
是 $O(n)$，整支程式是 $O(n \log n)$。
\item 最多 $n(n - 1)/2 = 1999000$ 次。 每一筆申請都被選到時，第 $i$ 筆要和
前面選到的 $i - 1$ 筆全部比過，中途不會 \texttt{break}。 兩兩不重疊的輸入
就會這樣，例如 $[0, 1)$、$[1, 2)$、\ldots、$[1999, 2000)$。
\item 所有 \texttt{a[i]} 都相同時，每個 \texttt{i} 的 \texttt{j} 都從
\texttt{i} 走到 \texttt{n}，共 $n + (n - 1) + \cdots + 1 = n(n + 1)/2$ 次，
是 $O(n^2)$。 (a) 的 \texttt{i} 不會歸零，這裡的 \texttt{j} 每一輪都從
\texttt{i} 重新開始。 \texttt{a} 排好之後，\texttt{i} 變大時 \texttt{j} 能
走到的地方不會變小，所以把 \texttt{int j = i;} 移到迴圈外面、寫成
\texttt{int j = 0;}，\texttt{j} 只往前走，總共 $O(n)$。 這是第~5~章雙指標的
寫法。
\end{enumerate}

\item
\begin{enumerate}
\item 依結束時間排序，同結束時開始晚的在前：\#1、\#2、\#3、\#4、\#6、\#5、
\#7、\#8。 選 \#1 $[1, 4)$；\#2、\#3 都和 \#1 重疊；選 \#4 $[5, 7)$；\#6、
\#5、\#7 都和 \#4 重疊；選 \#8 $[8, 11)$。 選到 \#1、\#4、\#8 三筆。
\item \texttt{2 3 3}。 規則 1 先選 $[0, 6)$，接著 $[6, 10)$，$[8, 11)$ 和它
重疊，共 $2$ 筆。 規則 2 依長度看 \#2、\#4、\#1、\#8、\ldots，選到 \#2
$[3, 5)$、\#4 $[5, 7)$、\#8 $[8, 11)$ 三筆。
\end{enumerate}

\item 只比時刻，同一個時刻的加和減誰先誰後就沒有規定，可能先加後減，把不能
同時來的兩組算在一起。 輸入
\begin{code}{text}{練習 5 會錯的輸入}
2 1
3 2 3
2 0 1
\end{code}
正確答案是 \texttt{3 2}，在本書測試的環境下這樣改會印出 \texttt{5 2}。 把兩
組的順序對調，又會印出正確的答案：結果取決於輸入順序和 \texttt{sort} 內部
怎麼做，所以範例剛好過了也不代表對。 用 \texttt{pair} 預設的比較，或在時刻
相同時比 \texttt{second}，就固定是先減後加。

\item 任取一種排法，假設有相鄰的兩人 $i$、$j$，$i$ 在前，而且 $t_i > t_j$。
設排在他們前面的人總共花了 $P$ 分鐘。 互換前，$i$ 等 $P$、$j$ 等
$P + t_i$；互換後，$j$ 等 $P$、$i$ 等 $P + t_j$。 排在他們前面的人不受影
響；排在後面的人，前面多的仍然是 $t_i + t_j$，也不受影響。 所以總和的變化是
$(2P + t_j) - (2P + t_i) = t_j - t_i < 0$，互換之後嚴格變小。 只要還有相鄰
的兩人是長的在前，就能再變小，所以總和最小的排法一定沒有這種情形，也就是
由小到大排好（時間相同的人誰先都一樣）。 這種只比較相鄰兩個的寫法，是排序
類貪心最常用的交換論證。
\end{enumerate}

% 佔位：本章尚未開始。book.tex 裡先註解掉。
\chapter{遞迴與分治}

本章尚未開始。

\chapter{動態規劃（一）：一維}
\label{ch:dp1d}

郵局的櫃台只剩 $1$ 元、$5$ 元、$6$ 元三種郵票。 一封信要剛好貼滿 $n$ 元，
最少要貼幾張？ 第~7~章 7.1 節試過「先貼面額最大的」，
$10$ 元時貼出 $6 + 1 + 1 + 1 + 1$ 共五張，可是兩張 $5$ 元就夠了。 往下讀之
前先想一個問題：如果有人已經替你找到 $n$ 元的最佳貼法，把他最後貼上去的那一
張撕掉，剩下的郵票是什麼樣子？

第~7~章 7.4 節還留下另外兩個貪心解不了的問題：每場講座有不同
的價值時，要讓總價值最大；東西不能切開的背包。 三個問題有同一個難處：前面
的選擇會影響後面還能拿什麼，而哪個選擇比較好，要看到後面才知道。 貪心只看
眼前，所以會錯；把所有選法都試一遍又太慢。 這一章的\textbf{動態規劃}把兩邊
的優點合在一起：像暴力解一樣考慮所有選法，但同樣的小問題只算一次。

這一章是全書第一個動態規劃的章，從頭講起。 先用郵票走一遍完整的過程：暴力
遞迴、把算過的答案記下來、改成由小到大填表，再整理出設計動態規劃時要回答的
四件事。 接著換一個角度，把每個動態規劃都畫成一張圖，計算順序、複雜度、
「為什麼不會無限遞迴」都可以從圖上讀出來。 三題例題分別練計數、一個位置
兩種狀態、一個位置三種狀態；最後回到第~7~章的講座，說明貪心和動態規劃的
分界。 背包要用兩維的表格，留到《C++ 演算法解題 2》第~1~章。

銜接書第~10~章 10.3 節用最長遞增子序列示範過一次動態規劃，這一章不假設讀者
記得。 用到的語法只有遞迴函式（銜接書第~10~章 10.2 節）、\texttt{vector}、
\texttt{long long}，以及第~\ref{sec:dp1-greedy}~節的 \texttt{upper\_bound}
（銜接書第~9~章）。 沒讀過銜接書的讀者，只要會寫遞迴函式、會用
\texttt{vector} 當陣列，就能讀這一章。

\section{從郵票開始}
\label{sec:dp1-stamps}

\subsection{先寫暴力解}

把章首的問題寫清楚：用 $1$、$5$、$6$ 元的郵票貼滿 $n$ 元，每種張數不限，
最少要幾張？ 用 $\mathrm{best}(n)$ 表示答案。

回答章首的問題。 $n$ 元的最佳貼法裡，最後貼上的那一張只有三種可能：$1$ 元、
$5$ 元或 $6$ 元。 假設是 $5$ 元，撕掉它之後，剩下的郵票剛好貼滿 $n - 5$ 元，
而且一定是 $n - 5$ 元的最佳貼法：如果 $n - 5$ 元有張數更少的貼法，換過去再補
上這張 $5$ 元，$n$ 元就有更好的貼法，和「最佳」矛盾。 所以最後一張是 $5$ 元
時，總張數是 $\mathrm{best}(n - 5) + 1$。

不知道最後一張是哪一種，就三種都試，取最小的：
\[
\mathrm{best}(n) = 1 + \min\bigl(\mathrm{best}(n - 1),\ \mathrm{best}(n - 5),\
\mathrm{best}(n - 6)\bigr),
\]
其中 $n - 5$、$n - 6$ 小於 $0$ 的那一項不算（貼不下那一張）。 $0$ 元不用貼，
$\mathrm{best}(0) = 0$。

這個式子把「$n$ 元的答案」寫成「比較小的金額的答案」，照著寫就是一支遞迴程
式。 它試遍了所有可能的最後一張，再遞迴地試遍倒數第二張、倒數第三張……，等
於枚舉了所有貼法，是一個暴力解。 下面的程式多用一個全域變數數出函式一共被
呼叫幾次：

\begin{example}[label=ex:stamps-rec]{cpp}{(stamps\_rec.cpp) 郵票問題的暴力遞迴}
#include <iostream>
using namespace std;

long long calls = 0;   // best 一共被呼叫幾次

// 貼滿 n 元最少要幾張郵票（面額 1、5、6 元）
int best(int n) {
    calls++;
    if (n == 0) {
        return 0;                     // 0 元不用貼
    }
    int res = best(n - 1) + 1;        // 最後一張是 1 元
    if (n >= 5) {
        res = min(res, best(n - 5) + 1);
    }
    if (n >= 6) {
        res = min(res, best(n - 6) + 1);
    }
    return res;
}

int main() {
    for (int n : {10, 20, 30, 40, 50}) {
        calls = 0;
        int b = best(n);
        cout << n << ' ' << b << ' ' << calls << '\n';
    }
    return 0;
}
\end{example}

\iosample{}{
10 2 48 \\
20 4 1625 \\
30 5 53936 \\
40 7 1793825 \\
50 9 59637888 \\
}

每一行是「金額、最少張數、呼叫次數」。 答案是對的，$10$ 元是兩張，和第~7~章
的結論一樣。 問題在第三欄：金額每多 $10$，呼叫次數大約變成 $33$ 倍。 照這個
速度，$n = 100$ 要呼叫大約 $2 \times 10^{15}$ 次，在本書測試的環境下要跑好幾
天。

\subsection{同樣的問題算了好幾次}

慢的原因從呼叫的過程看得出來。 下圖是 $\mathrm{best}(7)$ 的呼叫樹（第~8~章）：
每個方格是一次呼叫，方格裡的數字是參數 $n$，往右或往下連到的是它呼叫的函式。
$\mathrm{best}(7)$ 呼叫 $6$、$2$、$1$；$\mathrm{best}(6)$ 呼叫 $5$、$1$、$0$；
依此類推。

\begin{center}
\begin{tikzpicture}[x=1.05cm, y=0.72cm,
    nd/.style={cell, minimum width=0.6cm, minimum height=0.55cm},
    hd/.style={hlcell, minimum width=0.6cm, minimum height=0.55cm},
    pd/.style={draw=optStroke, thick, fill=optFill, minimum width=0.6cm,
               minimum height=0.55cm, font=\ttfamily},
    ed/.style={draw=black!55}]
% 第一列：7-6-5-4-3-2-1-0
\node[nd] (a7) at (0,0) {7};
\node[nd] (a6) at (1,0) {6};
\node[nd] (a5) at (2,0) {5};
\node[nd] (a4) at (3,0) {4};
\node[nd] (a3) at (4,0) {3};
\node[hd] (a2) at (5,0) {2};
\node[pd] (a1) at (6,0) {1};
\node[nd] (a0) at (7,0) {0};
\foreach \p/\c in {a7/a6, a6/a5, a5/a4, a4/a3, a3/a2, a2/a1, a1/a0} {
    \draw[ed] (\p) -- (\c);
}
% best(5) 直接呼叫的第二個
\node[nd] (b0) at (3,-1) {0};
% best(6) 直接呼叫的第二、三個
\node[pd] (c1) at (2,-2) {1};
\node[nd] (c0) at (3,-2) {0};
\node[nd] (d0) at (2,-3) {0};
% best(7) 直接呼叫的第二、三個
\node[hd] (e2) at (1,-4) {2};
\node[pd] (e1) at (2,-4) {1};
\node[nd] (e0) at (3,-4) {0};
\node[pd] (f1) at (1,-5) {1};
\node[nd] (f0) at (2,-5) {0};
\draw[ed] (a5.south) |- (b0.west);
\draw[ed] (a6.south) |- (c1.west);
\draw[ed] (a6.south) |- (d0.west);
\draw[ed] (c1) -- (c0);
\draw[ed] (a7.south) |- (e2.west);
\draw[ed] (a7.south) |- (f1.west);
\draw[ed] (e2) -- (e1);
\draw[ed] (e1) -- (e0);
\draw[ed] (f1) -- (f0);
\end{tikzpicture}
\diagcaption{$\mathrm{best}(7)$ 的呼叫樹，一共 $17$ 次呼叫。 橘色的
$\mathrm{best}(2)$ 被呼叫 $2$ 次，粉紅色的 $\mathrm{best}(1)$ 被呼叫 $4$
次，$\mathrm{best}(0)$ 被呼叫 $6$ 次。}
\label{fig:stamps-tree}
\end{center}

參數只有 $0$ 到 $7$ 八種，卻呼叫了 $17$ 次，因為同一個參數被算了好幾次，每次
都從頭遞迴下去，得到一模一樣的答案。 $n$ 越大，重複越嚴重：$\mathrm{best}(n)$
底下有 $\mathrm{best}(n - 1)$ 和 $\mathrm{best}(n - 5)$ 兩棵大樹，而
$\mathrm{best}(n - 1)$ 底下又有一棵 $\mathrm{best}(n - 6)$，和
$\mathrm{best}(n)$ 直接呼叫的那棵完全相同。 樹的大小因此隨 $n$ 指數成長，可
是真正不同的問題只有 $n + 1$ 個。

\subsection{記住算過的答案}

第~8~章用過記憶化：開一個表格，每個參數第一次算完就把答案記下來，下
次再被呼叫時直接查表回傳。 用在郵票上：

\begin{example}[label=ex:stamps-memo]{cpp}{(stamps\_memo.cpp) 加上記憶化}
#include <iostream>
#include <vector>
using namespace std;

long long calls = 0;
vector<int> memo;                     // memo[n] == -1 表示還沒算過

int best(int n) {
    calls++;
    if (memo[n] != -1) {
        return memo[n];               // 算過了，直接查
    }
    int res;
    if (n == 0) {
        res = 0;
    } else {
        res = best(n - 1) + 1;
        if (n >= 5) {
            res = min(res, best(n - 5) + 1);
        }
        if (n >= 6) {
            res = min(res, best(n - 6) + 1);
        }
    }
    memo[n] = res;                    // 記下來再回傳
    return res;
}

int main() {
    int n;
    cin >> n;
    memo.assign(n + 1, -1);
    int b = best(n);
    cout << b << ' ' << calls << '\n';
    return 0;
}
\end{example}

\iosample{
50 \\
}{
9 142 \\
}

同樣是 $n = 50$，呼叫次數從將近六千萬次降到 $142$ 次。 每個參數只會真正計算
一次，每次計算最多再呼叫三次，所以總呼叫次數不超過 $1 + 3n$，是 $O(n)$。

\begin{itemize}
\item 答案一定不是負的，所以用 $-1$ 表示「還沒算過」。 如果答案可能是任何整
數，就要另外開一個 \texttt{bool} 陣列記錄算過了沒有。
\item 先查表、再計算、算完立刻存進 \texttt{memo}，三個步驟的位置都不能換。
例如把 \texttt{memo[n] = res;} 漏掉，程式仍然對，只是又變回暴力解的速度。
\item 遞迴有深度的限制。 \texttt{best(n)} 會一路呼叫 \texttt{best(n - 1)}、
\texttt{best(n - 2)}……，深度大約是 $n$。 在本書測試的環境下（堆疊
$8$~MB），這支程式的 $n$ 到大約 $17$ 萬就會因為堆疊用完而當掉。
\end{itemize}

\subsection{由小到大填表}

記憶化的程式是從 $\mathrm{best}(n)$ 開始往下問，問到底才開始填表。 換個方向
想：算 $\mathrm{best}(x)$ 只需要 $\mathrm{best}(x - 1)$、$\mathrm{best}(x - 5)$、
$\mathrm{best}(x - 6)$，都比 $x$ 小。 所以從 $x = 0$ 開始，由小到大一格一格
填，填到 $x$ 的時候，需要的三格一定已經填好了，根本不用遞迴：

\begin{example}[label=ex:stamps-table]{cpp}{(stamps\_table.cpp) 由小到大填表，並印出一種貼法}
#include <iostream>
#include <vector>
using namespace std;

int main() {
    int n;
    cin >> n;
    const int stamps[3] = {1, 5, 6};
    vector<int> best(n + 1), last(n + 1);   // last[x]：x 元的最佳貼法最後貼哪一張
    best[0] = 0;
    for (int x = 1; x <= n; x++) {
        best[x] = x + 1;                    // 比任何可能的答案都大
        for (int c : stamps) {
            if (c <= x && best[x - c] + 1 < best[x]) {
                best[x] = best[x - c] + 1;
                last[x] = c;
            }
        }
    }
    cout << best[n] << '\n';
    for (int x = n; x > 0; x -= last[x]) {  // 沿著 last 往回走，印出一種貼法
        cout << last[x] << (x - last[x] > 0 ? ' ' : '\n');
    }
    return 0;
}
\end{example}

\iosample{
23 \\
}{
4 \\
5 6 6 6 \\
}

下表是 $x = 0$ 到 $12$ 的計算過程。 「看三格」那三欄是三種最後一張各自得到的
張數，取最小的填進 \texttt{best}；最小值有好幾個時，程式記下先看到的那一個。

\begin{center}
\begin{tabular}{c|ccc|cc}
\toprule
$x$ & 最後 $1$ 元 & 最後 $5$ 元 & 最後 $6$ 元 & \texttt{best[x]} & \texttt{last[x]} \\
\midrule
0 & -- & -- & -- & 0 & -- \\
1 & 1 & -- & -- & 1 & 1 \\
2 & 2 & -- & -- & 2 & 1 \\
3 & 3 & -- & -- & 3 & 1 \\
4 & 4 & -- & -- & 4 & 1 \\
5 & 5 & 1 & -- & 1 & 5 \\
6 & 2 & 2 & 1 & 1 & 6 \\
7 & 2 & 3 & 2 & 2 & 1 \\
8 & 3 & 4 & 3 & 3 & 1 \\
9 & 4 & 5 & 4 & 4 & 1 \\
10 & 5 & 2 & 5 & 2 & 5 \\
11 & 3 & 2 & 2 & 2 & 5 \\
12 & 3 & 3 & 2 & 2 & 6 \\
\bottomrule
\end{tabular}
\end{center}

例如 $x = 10$：最後一張是 $1$ 元時要 $\mathrm{best}(9) + 1 = 5$ 張，是 $5$ 元
時要 $\mathrm{best}(5) + 1 = 2$ 張，是 $6$ 元時要 $\mathrm{best}(4) + 1 = 5$
張，最少是 $2$。 貪心在這裡錯，是因為它只試了「最後一張是 $6$ 元」這一種。

\begin{itemize}
\item \texttt{best[x] = x + 1} 是「目前還不知道，先放一個一定會被比下去的
數」。 有 $1$ 元郵票，$x$ 元最多 $x$ 張，所以 $x + 1$ 夠大。 面額裡沒有 $1$
元時，有些金額根本貼不出來，這個寫法就不對了，本章練習第~2~題會看到它錯在
哪裡。
\item \texttt{last[x]} 記下最佳貼法的最後一張。 從 $n$ 開始，每次減掉
\texttt{last[x]}，就沿著當初的選擇一路走回 $0$，印出的是從最後一張往前的
順序。 範例的 $23$ 元是 $5 + 6 + 6 + 6$。 表格裡只存數字，要知道「怎麼
做到」，就得另外記下每一格是從哪裡來的。
\item 雙層迴圈，外層 $n$ 次，內層 $3$ 次，時間 $O(n)$。 沒有遞迴，也就沒有
深度的限制，$n = 10^6$ 一樣瞬間算完。
\end{itemize}

\subsection{動態規劃的四件事}
\label{sec:dp1-four}

回頭看郵票走過的路：先找一個「大問題用小問題的答案算出來」的式子，照著寫成
遞迴是暴力解，加上記憶化或改成填表就快了。 這種做法叫動態規劃。 設計一個動
態規劃，要回答四件事：

\begin{enumerate}
\item \textbf{狀態}（state）：表格的一格代表什麼。 一定要用一句完整的中文寫
出來，例如「\texttt{best[x]}：剛好貼滿 $x$ 元最少要幾張」。 說不清楚的狀態，
後面三件事都會跟著錯。
\item \textbf{轉移}（transition）：一格怎麼由其他格算出來。 常用的想法是看
「最後一步」有哪幾種可能，每一種可能對應到一個比較小的狀態，再取最小、最大
或加總。
\item \textbf{初始值}：最小的那幾格，沒辦法再往下拆，直接寫出答案。 郵票是
$\mathrm{best}(0) = 0$。
\item \textbf{計算順序}：填到某一格時，轉移用到的格子必須已經填好。 郵票是
$x$ 由小到大。
\end{enumerate}

答案在表格的哪一格，也要確認：郵票是 \texttt{best[n]}，有的題目要取整張表的
最大值，有的要把幾格加起來。

記憶化和由小到大填表，算的是同一個式子、同一張表，只是填的順序不同：記憶化
從要的答案往回問，碰到哪格填哪格；填表事先排好順序，從頭填到尾。 兩者是同
一個動態規劃的兩種寫法。 解題時建議先用中文把四件事寫下來，再寫成填表的迴
圈；計算順序不好安排時（例如第~8~章那種參數會跳來跳去的遞迴），才用記憶化。

\subsection{常見錯誤}

\begin{itemize}
\item \textbf{只記答案，沒記算過了沒}：用 $0$ 表示「還沒算過」，但某些參數的
答案本來就是 $0$。 郵票的 \texttt{memo[0]} 就是 $0$，每次查到它都會以為沒算
過，重新計算。 這個例子剛好只是多算幾次，碰到答案常常是 $0$ 的題目，速度會
退回暴力解。 選一個不可能是答案的值，例如 $-1$。
\item \textbf{遞迴太深}：記憶化的程式在 $n$ 大的時候會當掉，而且範例和小的
測試資料都看不出來。 上面的 \texttt{stamps\_memo.cpp} 輸入 $200000$，在本書
測試的環境下程式當掉，什麼都沒印（堆疊用完）。 $n$ 到 $10^5$ 以上又能排出
計算順序時，寫成填表的迴圈。
\item \textbf{還沒填就先用}：計算順序排錯，用到的格子還是初始的 $0$。 例如把
郵票的迴圈改成 \texttt{x} 由大到小，\texttt{best[x - 1]} 這些格子都還沒算，
$n = 10$ 會印出 \texttt{1}，貼法印成十張 $1$ 元。
\end{itemize}

\section{把動態規劃畫成圖}
\label{sec:dp1-dag}

這一節換一個角度看同一件事。 把每個狀態畫成一個點，如果算 B 的時候要用到 A，
就從 A 畫一個箭頭到 B。 郵票的點是金額 $0, 1, 2, \ldots$，從 $x$ 有三個箭頭
分別指向 $x + 1$、$x + 5$、$x + 6$，意思是「在 $x$ 元的貼法後面再貼一張」。
下圖畫出 $0$ 到 $8$：

\begin{center}
\begin{tikzpicture}[x=1.25cm, y=1cm,
    st/.style={circle, draw=black!60, thick, minimum size=0.7cm, inner sep=0pt,
               font=\ttfamily\small},
    ar/.style={-{Stealth[length=1.8mm]}, draw=black!45},
    hl/.style={-{Stealth[length=2mm]}, very thick, draw=memStroke}]
\foreach \x in {0,...,8} {
    \node[st] (s\x) at (\x,0) {\x};
}
\foreach \x/\v in {0/0,1/1,2/2,3/3,4/4,5/1,6/1,7/2,8/3} {
    \node[idxlbl] at (\x,-1.15) {\v};
}
\node[idxlbl, left] at (-0.4,-1.15) {best};
\foreach \x in {0,1,2,3,4,5,7} {
    \pgfmathtruncatemacro{\y}{\x+1}
    \draw[ar] (s\x) -- (s\y);
}
\foreach \x in {0,...,3} {
    \pgfmathtruncatemacro{\y}{\x+5}
    \draw[ar] (s\x) to[bend left=32] (s\y);
}
\foreach \x in {1,2} {
    \pgfmathtruncatemacro{\y}{\x+6}
    \draw[ar] (s\x) to[bend right=32] (s\y);
}
\draw[hl] (s0) to[bend right=32] (s6);
\draw[hl] (s6) -- (s7);
\end{tikzpicture}
\diagcaption{郵票問題的狀態圖（金額 $0$ 到 $8$）。 直的箭頭是貼 $1$ 元，上面
的弧是 $5$ 元，下面的弧是 $6$ 元。 點下方的小字是 $\mathrm{best}$。 橘色是
$7$ 元的一種最佳貼法：$6 + 1$。}
\label{fig:stamps-dag}
\end{center}

從 $0$ 走到 $x$ 的每一條路，都對應一種貼法：走過的箭頭依序就是貼上的郵票。
所以「貼滿 $x$ 元最少幾張」就是「從 $0$ 走到 $x$ 最少要走幾個箭頭」，轉移
式裡的 $\min$ 是在比較「最後一個箭頭從哪一點進來」。 很多看起來不一樣的動態
規劃問題，放到圖上都是同一類問題：

\begin{center}
\begin{tabular}{ll}
\toprule
題目問的 & 在圖上是 \\
\midrule
最少（成本、次數） & 起點到終點最短的路 \\
最多（價值、分數） & 起點到終點最長的路 \\
有幾種方法 & 起點到終點有幾條不同的路 \\
做不做得到 & 起點走不走得到終點 \\
\bottomrule
\end{tabular}
\end{center}

例如問「貼滿 $6$ 元有幾種貼法（先後順序不同算不同）」，就是數從 $0$ 到 $6$
有幾條路：$1 \times 6$、$1 + 5$、$5 + 1$、$6$，共 $4$ 條。 轉移只要把
$\min$ 換成加法：$\mathrm{ways}(x) = \mathrm{ways}(x - 1) + \mathrm{ways}(x - 5)
+ \mathrm{ways}(x - 6)$，$\mathrm{ways}(0) = 1$。

這張圖有一個重要的性質：箭頭一定從小的金額指向大的金額，沿著箭頭走永遠不會
回到走過的點。 這種有方向、沒有\textbf{環}（從一點出發沿著箭頭走，又回到
原點）的圖，叫\textbf{有向無環圖}（directed acyclic graph，DAG）。 前面的幾
件事，從圖上都看得出來：

\begin{itemize}
\item \textbf{遞迴為什麼會停}：遞迴是沿著箭頭倒著走，圖上沒有環，倒著走也不
會繞回原地，最後一定走到沒有箭頭進來的點，也就是初始值。
\item \textbf{計算順序}：每個點都要排在「所有指向它的點」後面。 只要找得到這
樣的順序，就能由小到大填表。 郵票的箭頭都往右，從左到右就是一種順序；第~10
~章的格子圖，箭頭往右和往下，一列一列填也可以。
\item \textbf{記憶化和填表的差別}：記憶化從終點倒著走，只會碰到「走得到終點」
的那些點；填表把所有點都算一次。 大部分的題目兩者算的點差不多，少數題目的
狀態很多、真正用到的很少，記憶化會快很多。
\item \textbf{複雜度}：每個點算一次，每次看所有進來的箭頭，所以時間是「點數
加箭頭數」。 實際估計時用「狀態數乘上每個狀態要看幾個前一狀態」：郵票是
$n \times 3$。 這個估法對整章的題目都適用。
\end{itemize}

反過來，如果某個問題畫出來的圖有環，例如「計數器上有 $+5$ 和 $-3$ 兩個按鈕，
從 $0$ 按到 $n$ 最少幾下」，$x$ 可以走到 $x + 5$ 再走回 $x + 2$、再走回
$x - 1$……，沒有一個點能保證排在另一個點前面，轉移式照寫成遞迴會無限呼叫下
去。 這類問題要用第~11、12~章的廣度優先搜尋，本章練習第~4~題會再談。

動態規劃的程式裡從來不用真的把圖建出來，圖只是幫你想的工具。 設計狀態和轉
移時，在紙上畫幾個點和箭頭，檢查三件事：箭頭的意思是什麼、有沒有環、終點是
哪一點。 《C++ 演算法解題 2》第~6~章會在真正的圖上做同樣的事。

第~8~章追蹤遞迴時，從最小的參數開始往上列表，算回傳值或呼叫次數；現在看得出
來，那就是在這張圖上由小到大填表。 圖~\ref{fig:stamps-tree} 的 $17$ 次也可以
這樣算：$\mathrm{calls}(n) = 1 +$（它直接呼叫的那幾個參數的
$\mathrm{calls}$ 加起來），從 $n = 0$ 填到 $7$。

\section{計數：一格一格推過去}
\label{sec:dp1-count}

郵票問的是最少，這一節的例題問「有幾種」，同時練習兩件計數題一定會碰到的事：
答案太大要取餘數，以及轉移可以「往回拉」也可以「往前推」。

\subsection{例題：跳格子遊戲的走法數}
\label{sec:ex-hopscotch}

\textbf{題意}（\akiref{hopscotch-route-count}{跳格子遊戲的走法數}）

\begin{itemize}
\item 地上有位置 $0, 1, \ldots, n$，從位置 $0$ 出發，要剛好跳到位置 $n$。
\item 位置 $i$（$0 \le i < n$）寫著數字 $a_i$，$1 \le a_i \le 3$。 站在位置
$i$ 時，可以往前跳 $1$ 到 $a_i$ 格；跳過 $n$ 的跳法不能用。
\item 經過的落點不同就算不同的走法。 求走法數除以 $10^9 + 7$ 的餘數。
$n \le 2 \times 10^5$。
\end{itemize}

例如 $n = 4$，$a = (2, 1, 3, 1)$。 從 $0$ 可以到 $1$、$2$；從 $1$ 只能到
$2$；從 $2$ 可以到 $3$、$4$（到 $5$ 會跳過終點）；從 $3$ 只能到 $4$。 走法
有 $0 \to 1 \to 2 \to 3 \to 4$、$0 \to 1 \to 2 \to 4$、$0 \to 2 \to 3 \to 4$、
$0 \to 2 \to 4$，共 $4$ 種。

\textbf{想法}

先畫圖。 點是位置 $0$ 到 $n$，從 $i$ 畫箭頭到 $i + 1, \ldots, i + a_i$（不超
過 $n$）。 一種走法就是一條從 $0$ 到 $n$ 的路，題目問的是路的條數。

\begin{center}
\begin{tikzpicture}[x=1.6cm, y=1cm,
    st/.style={circle, draw=black!60, thick, minimum size=0.7cm, inner sep=0pt,
               font=\ttfamily\small},
    ar/.style={-{Stealth[length=1.8mm]}, draw=black!55}]
\foreach \x in {0,...,4} {
    \node[st] (p\x) at (\x,0) {\x};
}
\foreach \x/\v in {0/1,1/1,2/2,3/2,4/4} {
    \node[idxlbl] at (\x,-1.0) {\v};
}
\node[idxlbl, left] at (-0.4,-1.0) {ways};
\foreach \x/\a in {0/2,1/1,2/3,3/1} {
    \node[font=\scriptsize, text=black!60] at (\x,0.62) {$a=\a$};
}
\draw[ar] (p0) -- (p1);
\draw[ar] (p1) -- (p2);
\draw[ar] (p2) -- (p3);
\draw[ar] (p3) -- (p4);
\draw[ar] (p0) to[bend right=35] (p2);
\draw[ar] (p2) to[bend right=35] (p4);
\end{tikzpicture}
\diagcaption{範例的狀態圖。 點下方是抵達該位置的走法數。}
\label{fig:hop-dag}
\end{center}

最直接的做法是暴力遞迴：抵達位置 $i$ 的走法數，等於所有「一步能跳到 $i$」的
位置 $j$ 的走法數加起來。 每個位置最多跳 $3$ 格，所以 $j$ 只可能是 $i - 1$、
$i - 2$、$i - 3$，條件是 $j + a_j \ge i$。 這和郵票一樣有大量重複：每格都寫
$3$ 時，$n = 30$ 在本書測試的環境下要 $1.1$ 秒，$n = 35$ 要 $15$ 秒，每多
$5$ 格大約慢 $14$ 倍。

照四件事寫：

\begin{itemize}
\item 狀態：\texttt{ways[i]} 是「從位置 $0$ 跳到位置 $i$ 的走法數」。
\item 轉移：看最後一跳從哪裡起跳。 \texttt{ways[i]} 是所有滿足
$j + a_j \ge i$ 的 \texttt{ways[j]} 的和，$j$ 是 $i - 3$ 到 $i - 1$。
\item 初始值：\texttt{ways[0] = 1}。 站在起點、還沒跳，算一種走法。 這一格如
果是 $0$，後面每一格加來加去都是 $0$。
\item 計算順序：箭頭都往右，$i$ 由小到大。 答案是 \texttt{ways[n]}。
\end{itemize}

上面的轉移是站在 $i$ 往回看：「誰能跳到我？」 這種寫法叫\textbf{往回拉}。
這題往回看要檢查 $j + a_j \ge i$，有點麻煩；反過來站在 $j$ 往前看比較直接：
「我能跳到誰？」 \texttt{ways[j]} 算好之後，把它加到 \texttt{ways[j + 1]} 到
\texttt{ways[j + a\_j]} 上，這種寫法叫\textbf{往前推}。 兩種寫法加的是同一批
箭頭，只是一個在箭頭的終點加、一個在起點加。 往前推要注意一件事：處理到 $j$
時，所有指向 $j$ 的箭頭都必須已經推過了，\texttt{ways[j]} 才是完整的。 箭頭
都往右，$j$ 由小到大就滿足這個條件。

範例用往前推的過程：

\begin{center}
\begin{tabular}{cc|ccccc}
\toprule
處理位置 & $a_j$ & \texttt{ways[0]} & \texttt{[1]} & \texttt{[2]} & \texttt{[3]} & \texttt{[4]} \\
\midrule
開始 & & 1 & 0 & 0 & 0 & 0 \\
0 & 2 & 1 & 1 & 1 & 0 & 0 \\
1 & 1 & 1 & 1 & 2 & 0 & 0 \\
2 & 3 & 1 & 1 & 2 & 2 & 2 \\
3 & 1 & 1 & 1 & 2 & 2 & 4 \\
\bottomrule
\end{tabular}
\end{center}

處理位置 $2$ 時，$a_2 = 3$ 本來可以推到 $5$，超過終點的那一格不推。

\textbf{取餘數}

每格都寫 $3$ 時，走法數每多一格大約變成 $1.84$ 倍，$n = 74$ 就超過
\texttt{unsigned long long} 能存的範圍，所以題目只要餘數。 餘數的好處是可以
邊算邊取：$(x + y) \bmod M$ 等於 $\bigl((x \bmod M) + (y \bmod M)\bigr) \bmod M$，
乘法也一樣。 所以每做一次加法就取一次餘數，表格裡的數永遠小於
$M = 10^9 + 7$，兩個相加小於 $2.1 \times 10^9$，剛好還在 \texttt{int} 的
上限 $2^{31} - 1 \approx 2.147 \times 10^9$ 以內。 只要多加一項或做一次乘法
就會超過，所以表格一律存成 \texttt{long long}。 減法不能這樣直接做，《C++ 演算法解題 2》第~1~章會碰到。

同樣的模型也出現在別的地方：一份工作流程，每個步驟做完之後可以跳過後面幾個
步驟，問從頭到尾有幾種不同的執行順序；一條公車路線，每一站能轉搭的區間車可
以坐到後面幾站，問從起點到終點有幾種坐法。

\textbf{程式}

\begin{example}[label=ex:hopscotch]{cpp}{(hopscotch.cpp) 跳格子遊戲的走法數}
#include <iostream>
#include <vector>
using namespace std;

int main() {
    ios::sync_with_stdio(false);
    cin.tie(nullptr);
    const long long MOD = 1000000007;
    int n;
    cin >> n;
    vector<int> a(n);
    for (int &x : a) {
        cin >> x;
    }
    vector<long long> ways(n + 1, 0);   // ways[i]：抵達位置 i 的走法數
    ways[0] = 1;                        // 站在起點：一種（什麼都還沒跳）
    for (int i = 0; i < n; i++) {
        for (int d = 1; d <= a[i] && i + d <= n; d++) {
            ways[i + d] = (ways[i + d] + ways[i]) % MOD;   // 從 i 推給 i + d
        }
    }
    cout << ways[n] << '\n';
    return 0;
}
\end{example}

\iosample{
4 \\
2 1 3 1 \\
}{
4 \\
}

\begin{itemize}
\item 內層迴圈的條件 \texttt{i + d <= n} 同時做了兩件事：不把跳過終點的走法算
進來，也不寫到 \texttt{ways} 的範圍外面。
\item 外層迴圈只跑到 $n - 1$。 位置 $n$ 是終點，不用再往外推。
\end{itemize}

\textbf{複雜度}

每個位置往前推最多 $3$ 次，時間 $O(n)$，表格 $O(n)$。 最大的測試資料在本書
測試的環境下約 $0.04$ 秒。

\subsection{常見錯誤}

\begin{itemize}
\item \textbf{跳過終點也算抵達}：以為「跳到 $n$ 或更遠都算到了」，把
\texttt{ways} 開大兩格、拿掉 \texttt{i + d <= n}，最後把
\texttt{ways[n]}、\texttt{ways[n + 1]}、\texttt{ways[n + 2]} 加起來。 範例會
印出 \texttt{6}，多算了從位置 $2$ 直接跳到 $5$ 的兩種。 題目說「剛好跳到」，
就只能算 \texttt{ways[n]}。
\item \textbf{最後才取餘數}：表格用 \texttt{unsigned long long}，只在輸出時
\texttt{\% MOD}。 $n = 73$ 以內還看不出問題；輸入 $n = 80$、每格都是 $3$，
正確答案是 \texttt{997451737}，這樣寫會印出 \texttt{462595541}，中間早就溢
位了，也沒有任何警告。
\item \textbf{起點設成 $0$}：\texttt{ways[0] = 0}，整張表都是 $0$，範例印出
\texttt{0}。 計數題的起點幾乎都是 $1$：「什麼都還沒做」也是一種做法。
\item \textbf{沒有記憶化的遞迴}：照「往回拉」的式子寫成遞迴，範例和小的測試
資料都會過，$n$ 大一點就超過時間限制。
\end{itemize}

\section{選或不選：一個位置兩種狀態}
\label{sec:dp1-two}

前面兩個問題，狀態都只要一個數字（金額、位置）。 很多題目光知道「做到第
$i$ 個」不夠，還得知道第 $i$ 個是怎麼處理的，下一步才知道能做什麼。 這一節
的例題是最常見的形狀：每個東西選或不選，相鄰的不能都選。

\subsection{例題：招牌燈展示模式的總亮度}
\label{sec:ex-sign}

\textbf{題意}（\akiref{total-brightness-of-the-sign-demo}{招牌燈展示模式的總亮度}）

\begin{itemize}
\item 招牌上有一排 $n$ 盞燈，第 $i$ 盞亮起時亮度是 $a_i$。 相鄰的兩盞不能同
時亮；第 $1$ 盞和第 $n$ 盞不相鄰。
\item 替每盞燈決定亮或不亮，沒有相鄰兩盞同時亮，就是一種合法的亮法，全部不亮
也算。 一種亮法的亮度是亮著的燈的亮度總和。
\item 展示模式會把每一種合法的亮法各顯示一次。 求所有合法亮法的亮度加起來，
除以 $10^9 + 7$ 的餘數。 $n \le 2 \times 10^5$，$0 \le a_i \le 10^9$。
\end{itemize}

例如三盞燈的亮度是 $2, 5, 3$。 合法的亮法有全不亮、只亮第 $1$ 盞、只亮第
$2$ 盞、只亮第 $3$ 盞、亮第 $1$ 和第 $3$ 盞，共 $5$ 種，亮度是 $0 + 2 + 5 +
3 + 5 = 15$。

\textbf{想法}

暴力解是枚舉 $2^n$ 種亮或不亮的組合，挑出合法的加起來。 $n = 30$ 就有十億種，
$n$ 到 $2 \times 10^5$ 不可能。 這題一次要處理兩件事：合法亮法的條件，以及
「所有亮法的總和」。 分成三步，一次加一件事。

\textbf{第一步：只數有幾種亮法。} 試著用一個數字當狀態：$c_i$ 是「前 $i$ 盞燈
有幾種合法亮法」。 想從 $c_{i-1}$ 算 $c_i$ 時會卡住：第 $i$ 盞能不能亮，要看第
$i - 1$ 盞有沒有亮，而 $c_{i-1}$ 把兩種情況混在一起了。 缺的資訊放進狀態裡：

\begin{itemize}
\item \texttt{cOff[i]}：前 $i$ 盞的合法亮法中，第 $i$ 盞\textbf{不亮}的有幾種。
\item \texttt{cOn[i]}：前 $i$ 盞的合法亮法中，第 $i$ 盞\textbf{亮著}的有幾種。
\end{itemize}

轉移看第 $i$ 盞：不亮的話，第 $i - 1$ 盞亮不亮都可以；亮的話，第 $i - 1$ 盞
一定不亮。
\[
\texttt{cOff[i]} = \texttt{cOff[i-1]} + \texttt{cOn[i-1]},\qquad
\texttt{cOn[i]} = \texttt{cOff[i-1]}.
\]
初始值是「還沒有燈」：\texttt{cOff[0] = 1}（一種亮法，什麼都沒有，也可以看成
最後一盞不亮），\texttt{cOn[0] = 0}。 算下去：

\begin{center}
\begin{tabular}{c|cccccc}
\toprule
$i$ & 0 & 1 & 2 & 3 & 4 & 5 \\
\midrule
\texttt{cOff[i]} & 1 & 1 & 2 & 3 & 5 & 8 \\
\texttt{cOn[i]} & 0 & 1 & 1 & 2 & 3 & 5 \\
合計 & 1 & 2 & 3 & 5 & 8 & 13 \\
\bottomrule
\end{tabular}
\end{center}

三盞燈有 $5$ 種，和上面列舉的一樣。 畫成圖的話，每個 $i$ 有兩個點（第 $i$ 盞
亮、不亮），箭頭只從第 $i - 1$ 層指到第 $i$ 層，「亮」指向「亮」的箭頭不存在。

\textbf{第二步：最亮的一種亮法。} 同樣的兩個狀態，問的換成最大值，就是另一
個常見的題目：\texttt{mOff[i]}、\texttt{mOn[i]} 是第 $i$ 盞不亮、亮著時，前
$i$ 盞最大的亮度。 轉移把加法換成 $\max$，亮的那一邊加上 $a_i$：
\[
\texttt{mOff[i]} = \max(\texttt{mOff[i-1]}, \texttt{mOn[i-1]}),\qquad
\texttt{mOn[i]} = \texttt{mOff[i-1]} + a_i.
\]
$2, 5, 3$ 依序算出 $(0, 2)$、$(2, 5)$、$(5, 5)$，最亮是 $5$（只亮中間，或亮
兩邊）。 狀態一樣、轉移的形狀一樣，只是「怎麼合併」從加法換成 $\max$；這就
是上一節表格裡「路的條數」和「最長的路」的差別。

\textbf{第三步：所有亮法的亮度總和。} 題目要的是把每一種亮法的亮度都加起來。
只記總和夠不夠？ 設 \texttt{sOn[i]} 是第 $i$ 盞亮著的那些亮法的亮度總和。
第 $i$ 盞亮的亮法，都是「第 $i - 1$ 盞不亮的某一種亮法，再點亮第 $i$ 盞」。
前面有 \texttt{cOff[i-1]} 種這樣的亮法，每一種都多了 $a_i$，所以總和要加上
$a_i$ 乘以\textbf{種數}：
\[
\texttt{sOn[i]} = \texttt{sOff[i-1]} + a_i \times \texttt{cOff[i-1]},\qquad
\texttt{sOff[i]} = \texttt{sOff[i-1]} + \texttt{sOn[i-1]}.
\]
只記總和不夠，因為新加的 $a_i$ 要加幾次取決於種數。 所以每個狀態同時帶兩個
數字：有幾種、總和多少。 答案是 \texttt{sOff[n] + sOn[n]}。 範例的計算過程：

\begin{center}
\begin{tabular}{cc|cccc}
\toprule
$i$ & $a_i$ & \texttt{cOff} & \texttt{sOff} & \texttt{cOn} & \texttt{sOn} \\
\midrule
0 & -- & 1 & 0 & 0 & 0 \\
1 & 2 & 1 & 0 & 1 & $0 + 2 \times 1 = 2$ \\
2 & 5 & 2 & 2 & 1 & $0 + 5 \times 1 = 5$ \\
3 & 3 & 3 & 7 & 2 & $2 + 3 \times 2 = 8$ \\
\bottomrule
\end{tabular}
\end{center}

$7 + 8 = 15$。 $i = 3$ 那一列的 $3 \times 2$：第 $3$ 盞亮的亮法有兩種（只亮
第 $3$ 盞、亮第 $1$ 和第 $3$ 盞），各自多了亮度 $3$。

同樣的模型也出現在別的地方：一排攤位，相鄰的不能都開，求所有開法的營收總
和；或在一條街上選擇哪些路口裝監視器，相鄰路口不重複裝。 要「對所有選法加
總」時，記得同時帶著選法的數量。

\textbf{程式}

\begin{example}[label=ex:sign]{cpp}{(sign.cpp) 招牌燈展示模式的總亮度}
#include <iostream>
#include <vector>
using namespace std;

int main() {
    ios::sync_with_stdio(false);
    cin.tie(nullptr);
    const long long MOD = 1000000007;
    int n;
    cin >> n;
    // 前 i 盞、第 i 盞不亮（Off）或亮著（On）：c 是亮法數，s 是亮度總和
    vector<long long> cOff(n + 1), sOff(n + 1), cOn(n + 1), sOn(n + 1);
    cOff[0] = 1;                        // 還沒有燈：一種亮法，亮度 0
    sOff[0] = cOn[0] = sOn[0] = 0;
    for (int i = 1; i <= n; i++) {
        long long a;
        cin >> a;
        a %= MOD;
        cOff[i] = (cOff[i - 1] + cOn[i - 1]) % MOD;
        sOff[i] = (sOff[i - 1] + sOn[i - 1]) % MOD;
        cOn[i] = cOff[i - 1];                            // 前一盞必須不亮
        sOn[i] = (sOff[i - 1] + cOff[i - 1] * a) % MOD;  // 每一種都多了 a
    }
    cout << (sOff[n] + sOn[n]) % MOD << '\n';
    return 0;
}
\end{example}

\iosample{
3 \\
2 5 3 \\
}{
15 \\
}

\begin{itemize}
\item \texttt{a \%= MOD} 先把 $a_i$ 縮到 $M$ 以內。 $a_i$ 本來就小於 $M$，這
行在這題不會改變任何值，但習慣上乘法的兩邊都先取餘數。
\item \texttt{cOff[i - 1] * a} 的兩個數都小於 $M \approx 10^9$，乘積小於
$1.1 \times 10^{18}$，在 \texttt{long long} 的上限 $9.2 \times 10^{18}$ 以內。
兩個都是 \texttt{long long}，乘法就用 \texttt{long long} 做。
\item 每一格只用到前一格。 四個 \texttt{vector} 可以換成八個變數，只記「上一
格」和「這一格」，記憶體從 $O(n)$ 變成 $O(1)$。 這種寫法叫滾動陣列，第~10~章
在二維的表格上會用到。 這一章為了讓表格看得清楚，都用整個陣列。
\end{itemize}

\textbf{複雜度}

$O(n)$。 最大的測試資料在本書測試的環境下約 $0.05$ 秒。

\subsection{常見錯誤}

\begin{itemize}
\item \textbf{亮度只加一次}：寫成 \texttt{sOn[i] = sOff[i - 1] + a}，忘了前面
有好幾種亮法都要加。 範例印出 \texttt{12}。
\item \textbf{種數的初始值是 $0$}：\texttt{cOff[0] = 0}，所有的種數都是
$0$，總和也一直是 $0$，範例印出 \texttt{0}。 和跳格子一樣，「什麼都沒有」也
是一種。
\item \textbf{用 \texttt{int} 存}：把所有變數都改成 \texttt{int}，乘法
\texttt{cOff * a} 會溢位。 範例 2（四盞燈，亮度都是 $10^9$）的正確答案是
\texttt{999999937}，這樣寫在本書測試的環境下印出 \texttt{705032669}。
\item \textbf{最後才取餘數}：表格裡的加法、乘法都不取餘數，只在輸出時
\texttt{\% MOD}。 $100$ 盞亮度都是 $10^9$ 的燈，正確答案是
\texttt{726391504}，這樣寫印出 \texttt{-66304235}：中間的數早就超過
\texttt{long long}，變成負的。
\end{itemize}

\section{狀態要記得夠多：大夜班}
\label{sec:dp1-three}

上一節的「亮或不亮」是兩種狀態。 這一節的例題每一天有三種選擇，而且限制同
時牽涉到前一天和後一天，要想清楚「後一天」的限制要怎麼在由前往後的計算裡檢
查。

\subsection{例題：前後都要休息的大夜班}
\label{sec:ex-night}

\textbf{題意}（\akiref{rest-days-around-the-night-shift}{前後都要休息的大夜班}）

\begin{itemize}
\item 接下來 $n$ 天，每天從白班、大夜班、休息三種選一種。 第 $i$ 天上白班賺
$a_i$，上大夜班賺 $b_i$，休息賺 $0$。
\item 上大夜班的前一天和後一天都必須休息；第 $1$ 天沒有前一天、第 $n$ 天沒有
後一天，那一邊就不用管。 一個休息日可以同時是兩個大夜班的「後一天」和「前一
天」。 白班沒有限制。
\item 求最高的總薪水。 $n \le 2 \times 10^5$，$a_i \le 10^9$，
$b_i \le 3 \times 10^9$。
\end{itemize}

例如 $3$ 天，白班都是 $4$，大夜班是 $1, 10, 1$。 第 $2$ 天上大夜班可以賺
$10$，但第 $1$、$3$ 天就得休息；三天都上白班是 $12$，比較多。

\textbf{想法}

先試貪心：挑大夜班薪水最高的那天先排。 上面的例子就錯了，$10$ 很大，卻擋掉
兩天白班。 暴力解試遍每天的三種選擇，共 $3^n$ 種，$n = 20$ 就超過三十億。

照上一節的經驗，先問：排第 $i$ 天時，要知道前面的什麼？

\begin{itemize}
\item 第 $i$ 天想上大夜班：第 $i - 1$ 天必須休息。
\item 第 $i$ 天想上白班或大夜班：第 $i - 1$ 天不能是大夜班（大夜班的後一天要
休息）。
\end{itemize}

兩條都只和「第 $i - 1$ 天做什麼」有關。 「後一天必須休息」這條限制看起來是往
後看的，但它可以改成從後一天往前檢查：輪到第 $i + 1$ 天時，前一天是大夜班就
只能休息。 所以狀態是「前 $i$ 天排好、第 $i$ 天做的是哪一種」時，前面的薪水
最多是多少：

\begin{itemize}
\item \texttt{rest[i]}、\texttt{day[i]}、\texttt{night[i]}：排好前 $i$ 天，第
$i$ 天休息、上白班、上大夜班時，最高的總薪水。
\end{itemize}

哪一種可以接在哪一種後面，畫成圖是：

\begin{center}
\begin{tikzpicture}[x=1cm, y=1cm,
    sn/.style={circle, draw=black!60, thick, minimum size=1.1cm, font=\small},
    ar/.style={-{Stealth[length=2mm]}, thick, draw=black!60}]
\node[sn] (R) at (0,0) {休息};
\node[sn] (D) at (-2.6,-1.6) {白班};
\node[sn, draw=memStroke, fill=memFill] (G) at (2.6,-1.6) {大夜};
\draw[ar] (R) to[out=120, in=60, looseness=5] (R);
\draw[ar] (D) to[out=210, in=150, looseness=5] (D);
\draw[ar] (R) to[bend right=15] (D);
\draw[ar] (D) to[bend right=15] (R);
\draw[ar] (R) to[bend left=15] (G);
\draw[ar] (G) to[bend left=15] (R);
\end{tikzpicture}
\diagcaption{前一天（箭頭起點）做完之後，後一天（箭頭終點）可以做什麼。 大夜
班只和休息相連。}
\label{fig:night-states}
\end{center}

這張圖本身有環（休息可以接休息），但它描述的是「相鄰兩天」的關係。 真正的
狀態是（第幾天, 做什麼），箭頭一律從第 $i - 1$ 天指到第 $i$ 天，所以仍然是
有向無環圖，可以照天數由小到大算。 把圖上進入每個點的箭頭寫成式子，就是轉
移：
\[
\begin{aligned}
\texttt{rest[i]}  &= \max(\texttt{rest[i-1]}, \texttt{day[i-1]}, \texttt{night[i-1]}),\\
\texttt{day[i]}   &= \max(\texttt{rest[i-1]}, \texttt{day[i-1]}) + a_i,\\
\texttt{night[i]} &= \texttt{rest[i-1]} + b_i.
\end{aligned}
\]

兩個大夜班共用中間的休息日，在這個寫法裡自動成立：大夜、休息、大夜就是沿著
圖上的「大夜 $\to$ 休息 $\to$ 大夜」走。

\textbf{初始值}：第 $1$ 天沒有前一天，把第 $0$ 天當成休息，\texttt{rest[0] = 0}；
第 $0$ 天不可能上班，\texttt{day[0]} 和 \texttt{night[0]} 設成一個很小的數，
表示「這種結尾不存在」，取 $\max$ 時一定會輸。 這個很小的數不能直接用
\texttt{LLONG\_MIN}：它只要再加上一點薪水就會溢位。 用
\texttt{LLONG\_MIN / 4}，加上幾個薪水也還是很小的負數。

\textbf{答案}：第 $n$ 天三種都可以（第 $n$ 天的大夜班沒有後一天），取三者的
最大值。

範例 1 是 $7$ 天，白班 $3, 3, 2, 2, 6, 2, 1$，大夜班 $10, 5, 9, 8, 4, 12, 8$：

\begin{center}
\begin{tabular}{ccc|ccc}
\toprule
天 & $a_i$ & $b_i$ & \texttt{rest} & \texttt{day} & \texttt{night} \\
\midrule
0 & & & 0 & -- & -- \\
1 & 3 & 10 & 0 & 3 & 10 \\
2 & 3 & 5 & 10 & 6 & 5 \\
3 & 2 & 9 & 10 & 12 & 19 \\
4 & 2 & 8 & 19 & 14 & 18 \\
5 & 6 & 4 & 19 & 25 & 23 \\
6 & 2 & 12 & 25 & 27 & 31 \\
7 & 1 & 8 & 31 & 28 & 33 \\
\bottomrule
\end{tabular}
\end{center}

答案是 $33$。 從最後一格往回追：第 $7$ 天大夜班 $33 = 25 + 8$，來自第 $6$ 天
休息的 $25$；$25$ 來自第 $5$ 天白班的 $25 = 19 + 6$；依此類推，班表是「大夜、
休息、大夜、休息、白班、休息、大夜」。 貪心先挑最高的第 $6$ 天大夜班
（$12$），最多只能拿到 $31$。

同樣的模型也出現在別的地方：工廠的機台做高負載的工作之前和之後都要停機保養；
運動員的高強度訓練日前後要排恢復日。 只要「今天能做什麼只看昨天做了什麼」，
就把「昨天做了什麼」放進狀態。

\textbf{程式}

\begin{example}[label=ex:night]{cpp}{(night.cpp) 前後都要休息的大夜班}
#include <algorithm>
#include <climits>
#include <iostream>
#include <vector>
using namespace std;

int main() {
    ios::sync_with_stdio(false);
    cin.tie(nullptr);
    int n;
    cin >> n;
    vector<long long> a(n + 1), b(n + 1);
    for (int i = 1; i <= n; i++) {
        cin >> a[i];
    }
    for (int i = 1; i <= n; i++) {
        cin >> b[i];
    }
    const long long NONE = LLONG_MIN / 4;   // 不可能的結尾
    // 排好前 i 天、第 i 天是休息／白班／大夜班時的最高總薪水
    vector<long long> rest(n + 1), day(n + 1), night(n + 1);
    rest[0] = 0;                        // 第 0 天當作休息
    day[0] = night[0] = NONE;
    for (int i = 1; i <= n; i++) {
        rest[i] = max({rest[i - 1], day[i - 1], night[i - 1]});
        day[i] = max(rest[i - 1], day[i - 1]) + a[i];      // 前一天不能是大夜
        night[i] = rest[i - 1] + b[i];                     // 前一天一定休息
    }
    cout << max({rest[n], day[n], night[n]}) << '\n';
    return 0;
}
\end{example}

\iosample{
7 \\
3 3 2 2 6 2 1 \\
10 5 9 8 4 12 8 \\
}{
33 \\
}

\begin{itemize}
\item \texttt{max(\{x, y, z\})} 用大括號一次比三個數，要
\texttt{\#include <algorithm>}（銜接書第~8~章的 \texttt{initializer\_list}）。
\item $b_i$ 最大 $3 \times 10^9$，本身就超過 \texttt{int}，讀進來就要用
\texttt{long long}。 總和最多大約 $2 \times 10^5 \times 3 \times 10^9 =
6 \times 10^{14}$。
\item 這一題的薪水都不是負的，把 \texttt{NONE} 寫成 $0$ 其實也算得出正確答案
（多出來的「第 $0$ 天上班」不會讓任何結果變大）。 但習慣上不可能的狀態就標
成不可能，薪水可以是負數、或者題目問最小值時，才不會出錯。
\end{itemize}

\textbf{複雜度}

$O(n)$。 最大的測試資料在本書測試的環境下約 $0.06$ 秒，大部分花在讀入。

\subsection{常見錯誤}

\begin{itemize}
\item \textbf{把大夜班當成三天一塊}：把「休息、大夜、休息」當成一個固定佔三
天的區塊，用一個狀態 $f(i)$ 表示前 $i$ 天的最高薪水。 這樣兩個大夜班沒辦法
共用中間的休息日。 輸入
\begin{code}{text}{三天一塊會錯的輸入}
5
0 0 0 0 0
5 0 5 0 5
\end{code}
正確答案是 $15$（第 $1$、$3$、$5$ 天大夜班），這種寫法只得到 \texttt{10}。
\item \textbf{只檢查後一天}：把 \texttt{night[i]} 寫成
\texttt{max(rest[i - 1], day[i - 1]) + b[i]}，讓大夜班也能接在白班後面。 範例
2（$4, 4, 4$／$1, 10, 1$）會印出 \texttt{14}：第 $1$ 天白班、第 $2$ 天大夜
班、第 $3$ 天休息，違反了前一天要休息。 範例 1 印出 \texttt{37}。
\item \textbf{答案只看休息}：輸出 \texttt{rest[n]}，以為最後一天休息一定最
好。 範例 1 印出 \texttt{31}，最後一天的大夜班被漏掉了。
\item \textbf{用 \texttt{int} 讀}：\texttt{a}、\texttt{b} 宣告成
\texttt{vector<int>}。 輸入有 $3000000000$ 時讀取失敗，後面的數字都讀不進來。
兩天、白班 $1, 1$、大夜班 $3000000000, 0$ 的輸入，正確答案是
\texttt{3000000000}，這樣寫印出 \texttt{2147483647}。
\end{itemize}

\section{狀態的幾種常見寫法}
\label{sec:dp1-shapes}

三題例題的狀態都是「做到第 $i$ 個」再加上一點資訊。 一維的動態規劃，狀態大
多是下面四種之一。 讀題時先想是哪一種，再補上缺的資訊。

\begin{center}
\begin{tabular}{p{0.25\linewidth}p{0.36\linewidth}p{0.29\linewidth}}
\toprule
狀態 & 轉移在看什麼 & 本章的例子 \\
\midrule
前 $i$ 個（或金額 $x$）的最佳值、方法數 & 最後一步是哪一種 & 郵票、跳格子 \\
前 $i$ 個，而且第 $i$ 個是某種情況 & 第 $i$ 個和第 $i - 1$ 個能不能相鄰 & 招牌燈、大夜班 \\
以第 $i$ 個結尾的連續一段 & 接在前一段後面，或從這裡重新開始 & 下面的最大區間和 \\
前 $i$ 個切成好幾段 & 最後一段從哪裡開始 & 下面的分段 \\
\bottomrule
\end{tabular}
\end{center}

\textbf{以第 $i$ 個結尾。} 一家店記下 $7$ 天的淨利 $3, -5, 4, -1, 2, -6, 1$，
要找連續幾天、淨利總和最大的一段。 第~5~章用前綴和解過這個最大區間和：固定
右端，減掉前面最小的前綴和。 這裡換成動態規劃的角度。 「前 $i$ 天裡最好的一
段」不好轉移：第
$i + 1$ 天能不能接上去，要看那一段是不是剛好結束在第 $i$ 天。 所以把狀態定成
「\texttt{endAt[i]}：最後一天剛好是第 $i$ 天的各段裡，最大的總和」。 這樣的一
段只有兩種：接在以第 $i - 1$ 天結尾的最好的那段後面，或從第 $i$ 天自己開始。
\begin{code}{cpp}{以第 i 天結尾的最大總和}
endAt[0] = a[0];
for (int i = 1; i < n; i++) {
    endAt[i] = a[i] + max(endAt[i - 1], 0LL);   // 前面是負的就不要接
}
// 答案是所有 endAt[i] 裡最大的
\end{code}
例子裡 \texttt{endAt} 依序是 $3, -2, 4, 3, 5, -1, 1$，最大是 $5$，就是第
$3$ 到 $5$ 天的 $4 - 1 + 2$。 這種狀態的答案不在最後一格，要取全部的最大值。

\textbf{切成好幾段。} 題目要把一列東西依序切成幾段，每段有個成本，例如每段
最多 $C$ 個、每段的成本和它的長度或內容有關。 狀態是「前 $i$ 個已經切好的最
小成本」，轉移看最後一段是第 $j + 1$ 到第 $i$ 個：
$\mathrm{dp}(i) = \min_j \bigl(\mathrm{dp}(j) + \mathrm{cost}(j + 1, i)\bigr)$。
$j$ 有好幾個可以選，每個狀態要看的前一狀態變多了，複雜度是狀態數乘上 $j$ 的
個數。 有些題目會限制最後一段的長度，可以選的 $j$ 就少了。

\section{回到第 7 章：貪心和動態規劃}
\label{sec:dp1-greedy}

第~7~章 7.4 節的三個問題，郵票已經解決了。 這一節解決第一個：
每場講座有價值時，要讓聽到的總價值最大。 背包留到《C++ 演算法解題 2》第~1~章。

\subsection{有價值的講座}

第 7 章的反例是 $[0, 3)$ 價值 $1$、$[2, 6)$ 價值 $10$：照結束時間先選
$[0, 3)$，總價值只有 $1$。 換一組大一點的例子，六場講座：

\begin{center}
\begin{tikzpicture}[x=0.9cm, y=0.55cm]
\ivbar{6}{1}{3}{A（3）}
\ivbar[ivpick]{5}{2}{5}{B（6）}
\ivbar{4}{4}{7}{C（5）}
\ivbar[ivpick]{3}{6}{9}{D（4）}
\ivbar{2}{5}{10}{E（7）}
\ivbar[ivpick]{1}{9}{12}{F（4）}
\ivaxis{0}{12}{0.2}
\end{tikzpicture}
\diagcaption{六場講座，括號裡是價值。 橘色的 B、D、F 是總價值最大的選法，共
$14$。}
\label{fig:valued-talks}
\end{center}

第~7~章的規則 3（結束早的先選）選到 A、C、F，總價值 $12$。 換成「價值高的
先選」，會先選 E（$7$），接著 B（$6$），F 和 E 重疊，總價值 $13$。 兩條規則
都不是最好的。

\textbf{狀態}。 先把講座依結束時刻排好：A、B、C、D、E、F。 狀態是
「\texttt{best[i]}：只看最早結束的前 $i$ 場，能拿到的最大總價值」。

\textbf{轉移}。 看第 $i$ 場聽不聽：

\begin{itemize}
\item 不聽：答案和只看前 $i - 1$ 場一樣，\texttt{best[i - 1]}。
\item 聽：前面只能再選和它不重疊的講座，也就是結束時刻不晚於它開始時刻的那
些。 講座已經依結束時刻排好，這些講座剛好是最前面的某 $p$ 場（$p$ 隨 $i$ 而
定）。 價值是 \texttt{best[p]} 加上這一場的價值。
\end{itemize}

兩者取大的。 初始值 \texttt{best[0] = 0}，$i$ 由小到大算，答案是
\texttt{best[n]}。 每個 $i$ 都只有兩個箭頭進來：從 $i - 1$ 和從 $p$。

找 $p$ 就是在排好的結束時刻裡，數有幾個 $\le$ 這一場的開始時刻，正好是
\texttt{upper\_bound} 做的事（第~3~章、銜接書第~9~章）。
講座是半開區間，結束在 $s$ 的講座和從 $s$ 開始的講座不重疊，所以用
\texttt{upper\_bound}（找第一個 $> s$ 的），相接的那一場要算進來。

\begin{center}
\begin{tabular}{cccc|cc|c}
\toprule
$i$ & 講座 & 時段 & 價值 & $p$ & 聽：\texttt{best[p]} $+$ 價值 & \texttt{best[i]} \\
\midrule
1 & A & $[1, 3)$ & 3 & 0 & $0 + 3 = 3$ & 3 \\
2 & B & $[2, 5)$ & 6 & 0 & $0 + 6 = 6$ & 6 \\
3 & C & $[4, 7)$ & 5 & 1 & $3 + 5 = 8$ & 8 \\
4 & D & $[6, 9)$ & 4 & 2 & $6 + 4 = 10$ & 10 \\
5 & E & $[5, 10)$ & 7 & 2 & $6 + 7 = 13$ & 13 \\
6 & F & $[9, 12)$ & 4 & 4 & $10 + 4 = 14$ & 14 \\
\bottomrule
\end{tabular}
\end{center}

例如 F 從 $9$ 開始，結束時刻 $\le 9$ 的是 A、B、C、D 四場，$p = 4$；聽 F 的話
前面最多拿 \texttt{best[4] = 10}（B 和 D），加上 F 是 $14$，比不聽的 $13$ 多。

\begin{example}[label=ex:lectures]{cpp}{(lectures.cpp) 有價值的講座}
#include <algorithm>
#include <iostream>
#include <vector>
using namespace std;

struct Talk {
    long long s, e, v;   // 時段 [s, e)，價值 v
};

int main() {
    int n;
    cin >> n;
    vector<Talk> t(n);
    for (Talk &x : t) {
        cin >> x.s >> x.e >> x.v;
    }
    sort(t.begin(), t.end(), [](const Talk &x, const Talk &y) {
        return x.e < y.e;
    });
    vector<long long> ends(n);
    for (int i = 0; i < n; i++) {
        ends[i] = t[i].e;
    }
    // best[i]：只看結束最早的前 i 場，能拿到的最大總價值
    vector<long long> best(n + 1, 0);
    for (int i = 1; i <= n; i++) {
        const Talk &cur = t[i - 1];
        // 前 i - 1 場裡，結束時刻 <= cur.s 的有 p 場，它們剛好是最前面的 p 場
        int p = upper_bound(ends.begin(), ends.begin() + (i - 1), cur.s)
                - ends.begin();
        best[i] = max(best[i - 1], best[p] + cur.v);   // 不聽／聽這一場
    }
    cout << best[n] << '\n';
    return 0;
}
\end{example}

\iosample{
6 \\
1 3 3 \\
2 5 6 \\
4 7 5 \\
6 9 4 \\
5 10 7 \\
9 12 4 \\
}{
14 \\
}

\begin{itemize}
\item 表格的第 $i$ 格對應排序後的 \texttt{t[i - 1]}，多一格 \texttt{best[0]}
當「一場都不看」，轉移裡就不用特別處理 $p = 0$。
\item 只在前 $i - 1$ 場裡找：第 $i$ 場自己的結束時刻一定大於開始時刻，本來
就不會被算進去，限制範圍只是讓意思更清楚。
\item 排序 $O(n \log n)$，每場一次二分搜 $O(\log n)$，總共 $O(n \log n)$。
\item 把 \texttt{upper\_bound} 換成 \texttt{lower\_bound}，相接的講座會被當成
重疊：E 接不上 B、F 接不上 D，範例印出 \texttt{12}。
\end{itemize}

為什麼這次排序沒有讓貪心變對？ 排序只是替狀態排出計算順序：「前 $i$ 場」要
有意義，聽第 $i$ 場時能接的講座就得剛好是前面一整段。 真正做決定的是轉移裡
的 $\max$：聽和不聽兩種都算過，取好的那一個。 貪心則是排好之後直接決定，不
回頭比較。

\subsection{不能切開的背包}

背包的問題是：最多裝重量 $W$，每件物品有重量和價值，求最大總價值。 試著照
這一章的方法定狀態：「只看前 $i$ 件物品的最大總價值」。 決定第 $i$ 件要不要
放時，要知道背包還剩多少空間，而這取決於前面放了哪些東西，前 $i - 1$ 件的最
大價值沒有告訴我們這件事。 缺的資訊要放進狀態：「前 $i$ 件物品、用掉的重量
是 $w$ 時的最大價值」，狀態變成兩個維度。 這是《C++ 演算法解題 2》第~1~章的主題。

\subsection{什麼時候用動態規劃}

\begin{itemize}
\item 題目問「最少」「最多」「有幾種」「做不做得到」，而所有選法一個一個試會
太多（$2^n$、$3^n$、$n!$）。
\item 選擇可以排成一個順序依序做（由左到右、由小到大），而且做到一半時，「之
後還能怎麼做」只取決於少量的資訊：目前的位置、金額、最後一個的情況。 這些
資訊就是狀態。
\item 試過的貪心找得到反例，或者對拍時和暴力解不一樣。
\item 估計「狀態數 $\times$ 每個狀態看幾個前一狀態」，在時間限制內。 照銜接
書 10.1 節的表，$n = 2 \times 10^5$ 時，每個狀態只能看常數個或
$O(\log n)$ 個前一狀態。
\end{itemize}

看起來像、但不能直接用的情況：

\begin{itemize}
\item 狀態之間的箭頭有環，例如可以往回走的移動、可以加也可以減的按鈕。 找
不到計算順序，要用第~11、12~章的廣度優先搜尋。
\item 「之後怎麼做」取決於整個過去，例如要記得前面選過哪些東西（不是幾個，是
哪些）。 狀態數是指數，在 $n$ 很小時可以用第~4~章的回溯。
\item 資料範圍大到連「狀態數 $\times$ 轉移數」都超過，而且找得到正確的貪心規
則，例如沒有價值的區間排程。 能證明的貪心比動態規劃快，程式也短。
\end{itemize}

第~7~章 7.4 節的實用判斷可以再補一句：範圍允許時，先寫動態規
劃。 動態規劃考慮了所有選法，對了就是對了，不需要交換論證；貪心快，但要證明
或至少對拍過。

\begin{apcsnote}
動態規劃的題目，正式的程式和暴力解常常只差在「同樣的狀態有沒有重算」，所以
暴力遞迴很容易寫出來，可以用來和填表的程式對拍。 時間不夠寫填表時，暴力遞
迴加上記憶化通常已經夠快；但遞迴深度到十幾萬層時可能會當掉，交出去之前用最
大的 $n$ 自己跑一次。
\end{apcsnote}

\section{本章常見錯誤}
\label{sec:dp1-pitfalls}

各節的錯誤，歸納起來是五類。 寫完動態規劃，交出去之前逐項檢查。

\begin{center}
\begin{tabular}{p{0.2\linewidth}p{0.46\linewidth}p{0.22\linewidth}}
\toprule
錯誤 & 例子 & 見 \\
\midrule
狀態少了資訊 & 招牌燈只記「前 $i$ 盞有幾種」；大夜班只記「前 $i$ 天最多多少」，
或把大夜班當成三天一塊 & 第~\ref{sec:ex-sign}、\ref{sec:ex-night}~節 \\
轉移漏了條件 & 大夜班只檢查後一天；跳格子跳過終點也算；亮度只加一次 &
第~\ref{sec:ex-hopscotch}、\ref{sec:ex-sign}、\ref{sec:ex-night}~節 \\
初始值錯 & 計數題的起點設成 $0$；不可能的狀態設成 $0$ 或「不夠大」的數；用
\texttt{LLONG\_MIN} 再加數字 & 第~\ref{sec:dp1-stamps}、\ref{sec:ex-night}~節、練習第~2~題 \\
計算順序錯 & 由大到小填，用到還沒算的格子；記憶化遞迴太深 &
第~\ref{sec:dp1-stamps}~節 \\
答案取錯格 & 只看 \texttt{rest[n]}；以第 $i$ 個結尾的狀態只看最後一格 &
第~\ref{sec:ex-night}、\ref{sec:dp1-shapes}~節 \\
\bottomrule
\end{tabular}
\end{center}

另外兩個和動態規劃本身無關、但幾乎每一題都會碰到的錯誤：取餘數只在最後做，
以及總和、乘積超過 \texttt{int}。 前者每做一次加法或乘法就取一次；後者在動
手前先估表格裡最大的數。

\section{本章練習}
\label{sec:dp1-exercises}

\begin{drills}[label=drills:dp1]
\item 下面的函式：
\begin{code}{cpp}{練習 1}
int g(int n) {
    if (n <= 2) {
        return n;
    }
    return g(n - 1) + 2 * g(n - 3);
}
\end{code}
\begin{enumerate}
\item 列出 $g(0)$ 到 $g(7)$ 的值。
\item 呼叫 \texttt{g(7)}，\texttt{g} 一共被呼叫幾次（包括 \texttt{g(7)} 自
己）？
\item 照範例~\ref{ex:stamps-memo} 的方式加上記憶化（查表的呼叫也算一次），
\texttt{g(7)} 一共被呼叫幾次？
\end{enumerate}

\item 郵局的 $1$ 元郵票也賣完了，只剩 $5$ 元和 $6$ 元。
\begin{enumerate}
\item $1$ 到 $30$ 元裡，哪些金額貼不出來？
\item 把範例~\ref{ex:stamps-table} 的 \texttt{stamps} 改成 $\{5, 6\}$，其他不
變，\texttt{best[7]} 會算出多少？ 為什麼 \texttt{x + 1} 在這裡不能當「一定
會被比下去的數」？
\item 怎麼改才對？ 貼不出來時輸出 $-1$。
\end{enumerate}

\item 回到 $1$、$5$、$6$ 元三種郵票。
\begin{enumerate}
\item 貼的先後順序不同算不同的貼法（$1 + 5$ 和 $5 + 1$ 不同），貼滿 $12$ 元
有幾種貼法？ 寫出狀態、轉移和初始值。
\item 如果只看每種用了幾張，不管先後（$1 + 5$ 和 $5 + 1$ 算同一種），(a) 的
轉移會多算。 為什麼？
\end{enumerate}

\item 計數器顯示 $0$，有兩個按鈕：$+5$ 和 $-3$。 要讓計數器顯示 $n$，最少按
幾下？ 有人照郵票的想法寫出
$f(x) = 1 + \min\bigl(f(x - 5), f(x + 3)\bigr)$，$f(0) = 0$，照著寫成遞迴。
這支程式會怎樣？ 用狀態圖說明原因。

\item 大夜班的規定改成：前一天要休息，\textbf{後兩天}都要休息。 寫出狀態和
轉移。（提示：休息有兩種，一種是大夜班後的第一天，隔天還得休息。）

\item 招牌燈的規定改成：可以有相鄰的燈同時亮，但不能有連續三盞都亮。
\begin{enumerate}
\item 寫出計算「合法亮法有幾種」的狀態和轉移。
\item $n = 5$ 時有幾種？
\end{enumerate}

\item 圖~\ref{fig:valued-talks} 的講座，如果 E 的價值從 $7$ 改成 $9$，寫出新
的 \texttt{best} 表和答案。

\item 有人想直接用「\texttt{pre[i]}：前 $i$ 天裡，總和最大的連續一段」當狀態
來求最大區間和。 給一個例子說明，只知道 \texttt{pre[i - 1]} 和 $a_i$，算不
出 \texttt{pre[i]}。
\end{drills}

\section{實作練習}
\label{sec:dp1-practice}

下面是 Aki-Judge 上和本章相關的題目，依第~\ref{sec:dp1-shapes}~節的四種狀態
分組。 難度分四階：暖身是直接套用本章的寫法；基本要自己定出狀態；實戰的題敘
比較長，要先把題目轉換成本章的形狀；挑戰會和其他章的技巧合在一起。

\begin{center}
\begin{tabular}{llc}
\toprule
狀態 & 題目 & 難度 \\
\midrule
前 $i$ 個／金額 $x$ & \akiref{springboard-track}{跳板賽道} & 暖身 \\
 & \akiref{doubling-portal-levels}{翻倍傳送門的闖關花費} & 暖身 \\
 & \akiref{fewest-presses-on-the-counter}{計數器按鈕的最少按法} & 暖身 \\
 & \akiref{value-and-call-count-of-a-recursive-function}{遞迴函式的值與呼叫次數} & 暖身 \\
 & \akiref{stickers-along-the-puddled-trail}{避開水窪收集貼紙} & 暖身 \\
 & \akiref{day-passes-and-two-day-passes}{游泳池的單日票與雙日票} & 基本 \\
 & \akiref{fewest-hops-down-a-cracked-corridor}{避開破地板的最少跳躍次數} & 基本 \\
 & \akiref{countdown-card-game-counting-plays}{倒數桌遊的走法數} & 基本 \\
 & \akiref{hopscotch-product-score}{跳格子的乘積分數} & 基本 \\
 & \akiref{when-the-clerks-snack-rule-fails}{零食包裝的貪心反例} & 基本 \\
 & \akiref{stepping-stones-the-smallest-jump-code}{踏石過溪的跳法代號} & 基本 \\
 & \akiref{catering-and-cleanup-days}{外燴後的整備日} & 基本 \\
 & \akiref{splitting-the-exam-stack}{考卷分疊批改的最低成本} & 基本 \\
 & \akiref{hiking-with-mandatory-water-stops}{不能錯過補水站的健行} & 挑戰 \\
\midrule
第 $i$ 個的情況 & \akiref{seating-plans-in-the-study-room}{自習室座位的安排數} & 暖身 \\
 & \akiref{total-cost-of-the-bead-sample-book}{串珠樣品冊的總成本} & 暖身 \\
 & \akiref{least-effort-with-bike-gears}{單車變速的最小總體力} & 暖身 \\
 & \akiref{poster-rows-without-equal-neighbors}{相鄰不同寬的海報排法} & 基本 \\
 & \akiref{corridor-tiles-around-floor-drains}{避開排水孔的走廊地磚} & 基本 \\
 & \akiref{reading-the-shelf-robots-log}{貨架機器人的移動紀錄} & 基本 \\
 & \akiref{back-to-back-review-bonus}{三科複習的連讀加分} & 基本 \\
 & \akiref{adjacent-channel-interference}{相鄰頻道的干擾} & 基本 \\
 & \akiref{best-seam-score-of-a-tile-row}{磁磚接縫的最大總分} & 基本 \\
 & \akiref{walk-or-ride-a-shared-bike}{走路還是騎公共自行車} & 基本 \\
 & \akiref{training-plans-with-a-cool-down}{以慢跑收操的訓練排法} & 基本 \\
 & \akiref{sprinklers-along-the-garden-beds}{菜園灑水頭的最大淨收成} & 基本 \\
 & \akiref{water-fountains-on-a-loop-trail}{環狀步道的飲水機} & 基本 \\
 & \akiref{to-redeem-or-not-to-redeem}{會員點數該不該折抵} & 實戰 \\
 & \akiref{best-part-time-earnings-with-a-day-off}{請假日的打工最佳收入} & 挑戰 \\
 & \akiref{band-frequencies-without-exact-multiples}{不能剛好差 m 倍的樂器頻率} & 挑戰 \\
\bottomrule
\end{tabular}
\end{center}

\begin{center}
\begin{tabular}{llc}
\toprule
狀態 & 題目 & 難度 \\
\midrule
以第 $i$ 個結尾 & \akiref{best-run-for-the-pop-up-shop}{快閃店的最佳檔期} & 暖身 \\
 & \akiref{quiz-score-lead}{小考領先幅度} & 暖身 \\
 & \akiref{exhibition-run-covering-a-guest-day}{涵蓋來賓日的最佳展期} & 基本 \\
 & \akiref{best-stretch-past-a-lookout}{經過觀景台的最佳路段} & 基本 \\
 & \akiref{best-stretch-in-a-compressed-ledger}{壓縮帳本的最佳時段} & 基本 \\
 & \akiref{bike-ride-between-rental-stations}{借還車站之間的河濱騎行} & 基本 \\
 & \akiref{add-cards-and-flip-cards}{加分卡與反轉卡} & 基本 \\
 & \akiref{flip-a-stretch-around-the-table}{圓桌翻牌} & 基本 \\
 & \akiref{a-training-camp-of-at-least-l-days}{至少連續 L 天的訓練營} & 基本 \\
 & \akiref{stalls-between-the-trees}{隔著奇數棵樹的攤位} & 實戰 \\
 & \akiref{food-truck-run-with-limited-moves}{最多換場 k 次的餐車檔期} & 實戰 \\
 & \akiref{on-warm-or-off}{開、保溫還是關} & 實戰 \\
 & \akiref{two-training-camps-not-too-far-apart}{休息不能太久的兩次集訓} & 挑戰 \\
\midrule
切成好幾段 & \akiref{height-gaps-in-graduation-group-photos}{畢業照分組的身高落差} & 基本 \\
 & \akiref{dishwashing-rounds}{洗碗輪次的放洗順序} & 基本 \\
 & \akiref{dispatching-shuttles-in-batches}{接駁車的分批發車} & 實戰 \\
 & \akiref{practice-blocks-for-band-camp}{社團集訓的連續練習區塊} & 實戰 \\
\bottomrule
\end{tabular}
\end{center}

本章的三題例題也在 Aki-Judge 上，讀完各節之後可以先自己寫一次再對照。

\section{讀完你應該會}
\label{sec:dp1-checklist}

\begin{itemize}
\item 拿到一個「最少」「最多」「有幾種」的問題，能用一句中文寫出狀態，並用
「最後一步有哪幾種」寫出轉移、初始值和計算順序。
\item 能把一個暴力遞迴加上記憶化，或改寫成由小到大填表的迴圈，並說得出兩者
算的是同一張表。
\item 能把一個動態規劃畫成狀態圖，從圖上判斷計算順序、估計複雜度，並看出圖
有環時不能直接用。
\item 轉移時發現「只知道前 $i$ 個的答案不夠」，能把缺的資訊（最後一個亮不
亮、昨天做什麼、是否以第 $i$ 個結尾）加進狀態。
\item 計數題能在每次運算後取餘數、正確設定起點為 $1$；最佳化題能用「不可能」
的值標記到不了的狀態。
\end{itemize}

\textbf{延伸閱讀}（書末會整理完整的參考資料）

\begin{itemize}
\item Antti Laaksonen，\textit{Competitive Programmer's Handbook}，第 7 章：
從找零錢開始，短短幾頁走過記憶化、填表與計數，可以很快複習一遍。
\item 劉汝佳，《算法競賽入門經典》第 2 版，9.1、9.2 節：用有向無環圖的角度看
動態規劃，例題很多。
\item Thomas H. Cormen 等，\textit{Introduction to Algorithms}，第 4 版，14.1、
14.3 節：動態規劃的嚴謹說明，包括什麼樣的問題可以拆成子問題。
\end{itemize}

\addcontentsline{toc}{section}{練習解答}
\section*{練習解答}

\begin{enumerate}
\item
\begin{enumerate}
\item 由小到大填：$g(0) = 0$、$g(1) = 1$、$g(2) = 2$、$g(3) = 2 + 0 = 2$、
$g(4) = 2 + 2 = 4$、$g(5) = 4 + 4 = 8$、$g(6) = 8 + 4 = 12$、
$g(7) = 12 + 8 = 20$。
\item 設 $c(n)$ 是呼叫 \texttt{g(n)} 一共引發的呼叫次數，$c(n) = 1$
（$n \le 2$），否則 $c(n) = 1 + c(n - 1) + c(n - 3)$。 依序是
$1, 1, 1, 3, 5, 7, 11, 17$，答案 $17$。 這也是一個動態規劃。
\item $11$ 次。 $g(7)$ 到 $g(3)$ 各真正計算一次，每次再呼叫兩次，共
$5 \times 2 = 10$ 次，加上最外面的 \texttt{g(7)}。
\end{enumerate}

\item
\begin{enumerate}
\item $1, 2, 3, 4, 7, 8, 9, 13, 14, 19$。 $20$ 以上都貼得出來：$20$ 到 $24$
分別是四張郵票的各種組合（$5 \times 4$、$5 \times 3 + 6$、\ldots、
$6 \times 4$），之後每多 $5$ 元再加一張 $5$ 元。
\item 算出 $3$。 \texttt{best[1]} 到 \texttt{best[4]} 貼不出來，留在
\texttt{x + 1}，也就是 $2, 3, 4, 5$；接著 \texttt{best[7]} 從
\texttt{best[1] + 1 = 3} 得到 $3$，看起來像一個正常的答案。 「還沒找到」的
值被當成真的張數往後傳了。 有 $1$ 元郵票時，每個金額都貼得出來，這個值一定
會被蓋掉，所以範例~\ref{ex:stamps-table} 沒事。
\item 用一個明確的「不可能」值，轉移時跳過它：
\begin{code}{cpp}{練習 2(c)}
const int INF = 1e9;
best[0] = 0;
for (int x = 1; x <= n; x++) {
    best[x] = INF;
    for (int c : {5, 6}) {
        if (c <= x && best[x - c] != INF) {
            best[x] = min(best[x], best[x - c] + 1);
        }
    }
}
cout << (best[n] == INF ? -1 : best[n]) << '\n';
\end{code}
\texttt{INF} 不要用 \texttt{INT\_MAX}：只要有一個地方忘了檢查就加 $1$，會溢
位成負數，變成「最小」。
\end{enumerate}

\item
\begin{enumerate}
\item $29$ 種。 狀態 $w(x)$：剛好貼滿 $x$ 元、依序貼上的貼法數；轉移看最後
一張，$w(x) = w(x - 1) + w(x - 5) + w(x - 6)$（小於 $0$ 的項不算）；初始值
$w(0) = 1$。 依序是 $1, 1, 1, 1, 1, 2, 4, 6, 8, 10, 13, 19, 29$。 在狀態圖
上，這是從 $0$ 到 $12$ 的路的條數。
\item 這個轉移把每一種「最後一張」都當成不同的貼法，所以 $1 + 5$ 和 $5 + 1$
分別從 $w(5)$ 和 $w(1)$ 各算一次。 不管先後時，要規定郵票的使用順序（例如先
決定 $1$ 元用幾張，再決定 $5$ 元、$6$ 元），狀態裡要多記「目前考慮到第幾種面
額」，這是《C++ 演算法解題 2》第~1~章的寫法。
\end{enumerate}

\item 會無限遞迴，直到堆疊用完而當掉。 $f(1)$ 要用到 $f(4)$，$f(4)$ 要用到
$f(7)$，$f(7)$ 又要用到 $f(2)$ 和 $f(10)$……參數一直變大，遞迴永遠不會碰到
$f(0)$ 之外的終點。 狀態圖上有 $x \to x + 5$、$x \to x - 3$ 兩種箭頭，會繞
成環，例如 $0 \to 5 \to 2 \to 7 \to 4 \to 1 \to 6 \to 3 \to 0$，找不到一個
「每個點都排在指向它的點後面」的順序。 這種每按一下成本都是 $1$、狀態會繞回
來的問題，用第~11、12~章的廣度優先搜尋，從 $0$ 開始一層一層往外找。

\item 四種狀態：\texttt{rest}（休息，隔天可以做任何事）、\texttt{day}、
\texttt{night}、\texttt{after}（大夜班後的第一個休息日，隔天還得休息）。
\[
\begin{aligned}
\texttt{rest[i]}  &= \max(\texttt{rest[i-1]}, \texttt{day[i-1]}, \texttt{after[i-1]}),\\
\texttt{day[i]}   &= \max(\texttt{rest[i-1]}, \texttt{day[i-1]}) + a_i,\\
\texttt{night[i]} &= \texttt{rest[i-1]} + b_i,\\
\texttt{after[i]} &= \texttt{night[i-1]}.
\end{aligned}
\]
大夜班的隔天只能是 \texttt{after}，\texttt{after} 的隔天只能是 \texttt{rest}，
兩天休息就保證了。 \texttt{night} 只能接在 \texttt{rest} 後面：\texttt{after}
也是休息日，但它的隔天必須休息，不能上大夜班。 初始值、答案都和原題一樣，答
案取四者的最大值。

\item
\begin{enumerate}
\item 狀態是「前 $i$ 盞、最後連續亮了 $k$ 盞」的種數，$k = 0, 1, 2$：
$c_0(i) = c_0(i-1) + c_1(i-1) + c_2(i-1)$，$c_1(i) = c_0(i-1)$，
$c_2(i) = c_1(i-1)$。 初始值 $c_0(0) = 1$，其他是 $0$。
\item $24$ 種。 每個 $i$ 的合計依序是 $2, 4, 7, 13, 24$。
\end{enumerate}

\item 只有 E 那一列變了：聽 E 是 \texttt{best[2]} $+ 9 = 15$，比 $10$ 大，
\texttt{best[5] = 15}。 F 那一列：聽 F 仍然是 $10 + 4 = 14$，不聽是 $15$，
\texttt{best[6] = 15}。 \texttt{best} 依序是 $3, 6, 8, 10, 15, 15$，答案
$15$，選 B 和 E。

\item 看兩組輸入：$(5, -10, 6)$ 和 $(2, 3, 6)$。 兩組的 \texttt{pre[2]}
都是 $5$（第一組只取第 $1$ 天，第二組取第 $1$、$2$ 天），第三天也都是 $6$。
可是第一組的 \texttt{pre[3]} 是 $6$（只取第 $3$ 天，接不上前面的 $5$，因為
中間隔著 $-10$），第二組是 $2 + 3 + 6 = 11$。 同樣的 \texttt{pre[2]} 和
$a_3$ 算出不同的答案，所以光靠它們算不出來：\texttt{pre[i - 1]} 沒有告訴我們
那一段是不是剛好結束在第 $i - 1$ 天。 改用「以第 $i$ 天結尾」的狀態就沒有這
個問題。
\end{enumerate}

\chapter{動態規劃（二）：二維格子}
\label{ch:dp2d}

住處在一片棋盤狀街區的西北角路口，捷運站在東南角路口：往東要走過 $4$ 個街
口，往南要走過 $3$ 個街口。 不想繞遠路，所以每一步只往東或往南走一個街口。
有幾條不同的走法？ 如果其中兩個路口在施工、不能通過呢？ 往下讀之前先想一
想：不要一條一條列，走到某一個路口的走法數，和它西邊、北邊那兩個路口的走法
數有什麼關係？

第~9~章的狀態是一個數字：金額 $x$、第 $i$ 天、第 $i$ 盞燈。 這一章的狀態是
兩個數字，列和行，表格也從一排變成一整張。 第~9~章整理的四件事（狀態、轉
移、初始值、計算順序）、記憶化和填表的關係、把狀態畫成圖，在這一章都照用，
不再重講。 這一章只處理二維多出來的問題：邊界怎麼處理、走不過去的格子、只
留一列以節省記憶體、一格要記不只一個數字，以及正方形和矩形這兩種常見的題
型。

三題例題分別練最小成本、一格兩種狀態、最大正方形。 計數和障礙格的基本寫法
在第一節用章首的街區示範。 用到的語法只有二維的 \texttt{vector} 和
\texttt{vector<string>}（銜接書第~3~章 3.6 節）。 沒讀過銜接書的讀者，只要知
道 \texttt{vector<vector<long long>> f(r, vector<long long>(c, 0))} 會建出一
張 $r$ 列 $c$ 行、全部是 $0$ 的表格，\texttt{f[i][j]} 是第 $i$ 列第 $j$ 行，
就能讀這一章。

\section{格子上的路徑}
\label{sec:dp2-paths}

\subsection{一格只看兩個鄰居}
\label{sec:dp2-two}

把路口排成 $4$ 列 $5$ 行的格子，第 $i$ 列第 $j$ 行的路口記為 $(i, j)$。 住
處是 $(1, 1)$，捷運站是 $(4, 5)$，往東是 $j$ 加 $1$，往南是 $i$ 加 $1$。 用
$f(i, j)$ 表示從住處走到 $(i, j)$ 有幾條路線。

回答章首的問題。 走到 $(i, j)$ 的路線，最後一步只有兩種可能：從西邊的
$(i, j - 1)$ 往東走過來，或從北邊的 $(i - 1, j)$ 往南走下來。 這兩組路線最後
一步不同，不會重複；任何一條路線都屬於其中一組，也不會漏掉。 所以
\[
f(i, j) = f(i - 1, j) + f(i, j - 1).
\]
照第~9~章的四件事寫下來：

\begin{itemize}
\item \textbf{狀態}：$f(i, j)$ 是從 $(1, 1)$ 走到 $(i, j)$ 的路線數。
\item \textbf{轉移}：上面的式子。 第一列沒有北邊、第一行沒有西邊，不存在的那
一邊當成 $0$ 條。
\item \textbf{初始值}：$f(1, 1) = 1$，原地不動算一條。
\item \textbf{計算順序}：算 $(i, j)$ 時，北邊和西邊必須已經算好。 一列一列由
上往下，每一列由左往右，就滿足這個條件。
\end{itemize}

計算順序也可以從狀態圖看出來。 把每一格畫成一個點，箭頭從每一格指向它東邊
和南邊的格子（下圖）。 箭頭只往右、往下，沿著箭頭走永遠回不到走過的點，所以
這是一張有向無環圖，第~9~章說過的事都成立。 只要每一格排在它的北邊和西邊之
後，什麼順序都可以：一列一列填可以，一行一行填也可以，沿著斜線
（$i + j$ 相同的格子）一條一條填也可以。 最常用的是一列一列填，因為輸入也是
一列一列給的。

\begin{center}
\begin{tikzpicture}[x=1.15cm, y=0.95cm,
    st/.style={circle, draw=black!60, thick, minimum size=0.62cm, inner sep=0pt,
               font=\ttfamily\small},
    hst/.style={st, draw=memStroke, fill=memFill},
    ar/.style={-{Stealth[length=1.6mm]}, draw=black!40},
    har/.style={-{Stealth[length=2mm]}, very thick, draw=memStroke}]
\foreach \i in {1,...,3} {
  \foreach \j in {1,...,4} {
    \node[st] (p\i\j) at (\j,-\i) {};
  }
}
\node[hst] (q) at (3,-2) {};
\foreach \i in {1,...,3} {
  \foreach \j [evaluate=\j as \k using int(\j+1)] in {1,...,3} {
    \draw[ar] (p\i\j) -- (p\i\k);
  }
}
\foreach \i [evaluate=\i as \k using int(\i+1)] in {1,2} {
  \foreach \j in {1,...,4} {
    \draw[ar] (p\i\j) -- (p\k\j);
  }
}
\draw[har] (p22) -- (q);
\draw[har] (p13) -- (q);
\node[font=\small, right] at (4.4,-1) {$(1, 1)$ 在左上角};
\node[font=\small, right] at (4.4,-2) {每個點只有兩個箭頭進來};
\end{tikzpicture}
\diagcaption{$3$ 列 $4$ 行格子的狀態圖。 橘色的點是 $(2, 3)$，只有從 $(2, 2)$
和 $(1, 3)$ 進來的兩個箭頭，所以它的答案只跟這兩格有關。}
\label{fig:dp2-dag}
\end{center}

下圖左邊是 $4 \times 5$ 的路口全部都能走時填好的表。 第一列和第一行都是
$1$，因為只有一種走法：一路往東或一路往南。 右下角是 $35$。

\begin{center}
\begin{tikzpicture}[x=0.62cm, y=0.62cm,
    g/.style={draw=black!55, minimum size=0.62cm, inner sep=0pt, font=\small\ttfamily},
    blk/.style={g, fill=black!28},
    end/.style={g, draw=memStroke, thick, fill=memFill}]
\begin{scope}
\node[font=\small] at (2,1) {全部能走};
\foreach \row [count=\r from 0] in {{1,1,1,1,1},{1,2,3,4,5},{1,3,6,10,15},{1,4,10,20,35}} {
  \foreach \v [count=\c from 0] in \row {
    \node[g] at (\c,-\r) {\v};
  }
}
\node[end] at (4,-3) {35};
\end{scope}
\begin{scope}[xshift=4.6cm]
\node[font=\small] at (2,1) {$(1, 3)$、$(3, 4)$ 施工};
\foreach \row [count=\r from 0] in {{1,1,0,0,0},{1,2,2,2,2},{1,3,5,0,2},{1,4,9,9,11}} {
  \foreach \v [count=\c from 0] in \row {
    \node[g] at (\c,-\r) {\v};
  }
}
\node[blk] at (2,0) {x};
\node[blk] at (3,-2) {x};
\node[end] at (4,-3) {11};
\end{scope}
\end{tikzpicture}
\diagcaption{兩張填好的表，格子裡是 $f(i, j)$。 灰色的 \texttt{x} 是施工中的
路口，$f = 0$。}
\label{fig:dp2-count}
\end{center}

全部能走的時候，其實有公式：$7$ 步裡挑哪 $3$ 步往南，$\binom{7}{3} = 35$。
可是只要有一個路口不能走，公式就不能直接用了，填表的做法卻完全不受影響。

最直接的暴力解是把路線一條一條列出來。 列出來的條數就是答案本身，$500 \times
500$ 的格子沒有障礙時有 $\binom{998}{499}$ 條，大約是 $10^{299}$，一條一條列
是不可能的。 照轉移式直接寫成遞迴、不加記憶化，呼叫次數也和路線數同一個數量
級（本章練習第~4~題）。 填表只要每一格做一次加法，$O(rc)$。

\subsection{走不過去的格子}
\label{sec:dp2-blocked}

章首的第二個問題：$(1, 3)$ 和 $(3, 4)$ 兩個路口在施工。 施工的路口走不進去，
所以 $f = 0$。 它的東邊和南邊在加總時，從它那裡拿到的是 $0$，等於那個方向沒
有路，不用另外處理。 圖~\ref{fig:dp2-count} 右邊是填好的表，答案是 $11$。

注意第一列：$(1, 3)$ 擋住之後，$(1, 4)$ 和 $(1, 5)$ 都是 $0$。 第一列只能從
西邊過來，西邊被擋，後面整段都到不了。 第一行也一樣。

\begin{example}[label=ex:dp2-streets]{cpp}{(streets.cpp) 避開施工路口的路線數}
#include <iostream>
#include <string>
#include <vector>
using namespace std;

int main() {
    int r, c;
    cin >> r >> c;
    vector<string> g(r);
    for (string &row : g) {
        cin >> row;
    }
    // f[i][j]：走到第 i 列第 j 行的路口有幾條路線
    // 列、行從 1 編號；第 0 列、第 0 行是多開的一圈，永遠是 0
    vector<vector<long long>> f(r + 1, vector<long long>(c + 1, 0));
    for (int i = 1; i <= r; i++) {
        for (int j = 1; j <= c; j++) {
            if (g[i - 1][j - 1] == 'x') {   // 施工中：到不了
                f[i][j] = 0;
            } else if (i == 1 && j == 1) {  // 起點：不動也算一條
                f[i][j] = 1;
            } else {
                f[i][j] = f[i - 1][j] + f[i][j - 1];
            }
        }
    }
    cout << f[r][c] << '\n';
    return 0;
}
\end{example}

\iosample{
4 5 \\
..x.. \\
..... \\
...x. \\
..... \\
}{
11 \\
}

\begin{itemize}
\item 地圖用 \texttt{vector<string>} 存，每一列讀成一個字串。 \texttt{g} 的
列、行從 $0$ 編號，表格 \texttt{f} 從 $1$ 編號，所以第 $i$ 列第 $j$ 行的路口
是 \texttt{g[i - 1][j - 1]}。
\item 先判斷施工，再判斷起點。 起點如果也在施工，答案就是 $0$。
\item 這個範例假設 $r, c \le 30$。 路線數最多是 $\binom{58}{29} \approx 3
\times 10^{16}$，\texttt{long long} 放得下。 格子再大就要像第~9~章那樣取餘數。
\end{itemize}

\subsection{邊界：多開一圈，或是先判斷}
\label{sec:dp2-border}

第一列的格子沒有北邊，第一行沒有西邊。 處理的方法有兩種。

\textbf{多開一圈。} 表格開成 $(r + 1) \times (c + 1)$，第 $0$ 列和第 $0$ 行
當作圍牆，填上「不可能」對應的值：計數題是 $0$，求最少的題目是一個很大的數，
求最多的題目是一個很小的數。 轉移式不用任何 \texttt{if}，上面的
\texttt{streets.cpp} 就是這樣寫。 只有起點要特別處理，否則
$f(1, 1) = f(0, 1) + f(1, 0) = 0$。

\textbf{先判斷再用。} 表格開成 $r \times c$，從 $0$ 編號，轉移時寫
\texttt{if (i > 0)}、\texttt{if (j > 0)}，鄰居存在才拿來用。 轉移除了鄰居的
答案、還要用到鄰居的其他資料（例如下一節要用鄰居的高度）時，這種寫法比較直
接，不用替圍牆編造資料。

兩種都可以，同一支程式裡選一種用到底。 最常見的錯誤是混在一起：表格多開了
一圈，讀地圖時卻忘了減一，或者反過來。

\subsection{常見錯誤}

\begin{itemize}
\item \textbf{第一列、第一行直接設成 $1$}：沒有障礙的時候，第一列和第一行都
是 $1$，於是有人先把它們全部設好，再從 $(2, 2)$ 開始填。 範例的輸入會印出
\texttt{15}：$(1, 4)$、$(1, 5)$ 明明被 $(1, 3)$ 擋住，卻被設成 $1$。 讓第一
列、第一行也照轉移式算，障礙就會自動往後傳。
\item \textbf{只用正方形的格子測試}：列和行寫反了，$r = c$ 的輸入完全看不出
來。 自己測試時，一定要有一組列數和行數不同的輸入，例如範例的 $4 \times 5$。
\end{itemize}

\section{最小成本}
\label{sec:dp2-cost}

問「有幾種」用加法，問「最少」就把加法換成取最小值，表格的形狀和計算順序都
不變。 這一節的例題還多了一點變化：成本不在格子上，而在兩格之間的移動上。

\subsection{例題：上坡才耗電的搬運車}
\label{sec:dp2-ex-uphill}

\textbf{題意}（\akiref{only-uphill-drains-the-battery}{上坡才耗電的搬運車}）

\begin{itemize}
\item 倉庫地面分成 $m$ 列 $n$ 行，第 $i$ 列第 $j$ 行的高度是 $h_{i,j}$。
\item 搬運車從左上角開到右下角，每一步只能往右或往下一格。
\item 從高度 $a$ 開到高度 $b$ 時，$b > a$ 要耗電 $b - a$；平地或下坡不耗電。
\item 求最少的總耗電量。 $m, n \le 500$，$0 \le h_{i,j} \le 10^9$。
\end{itemize}

例如高度是
\[
\begin{matrix}
1 & 5 & 2 & 2 \\
3 & 3 & 4 & 1 \\
6 & 2 & 2 & 0
\end{matrix}
\]
走「下、右、下、右、右」，高度依序是 $1, 3, 3, 2, 2, 0$，只有第一步爬坡 $2$，
答案是 $2$。

\textbf{想法}

最直接的做法是列出所有路線，每條算一次耗電量。 $m \times n$ 的格子共有
$\binom{m + n - 2}{m - 1}$ 條路線，$m = n = 500$ 時大約 $10^{299}$ 條，不可能。

再試一個貪心：每一步往耗電比較少的方向走，一樣少的時候往比較低的格子走，留
在低處，之後好像比較省。 反例是高度
\[
\begin{matrix}
5 & 1 \\
4 & 9 \\
8 & 6
\end{matrix}
\]
第一步往右或往下都是下坡，貪心選比較低的 $1$；之後只能往下爬到 $9$，耗電
$8$，再到 $6$，總共 $8$。 往下走 $5 \to 4 \to 8 \to 6$ 只要 $4$。 先下到很低
的地方，後面就得從低處爬上來。 一步的好壞要看後面的路才知道，這正是第~9~章
說的「要用動態規劃」的訊號。

\textbf{狀態}：\texttt{best[i][j]} 是從起點開到 $(i, j)$ 最少的耗電量。 這
一題的轉移要用到鄰居的高度，所以表格從 $0$ 編號，用「先判斷」的寫法處理邊界。

\textbf{轉移}：最後一步從上面的 $(i - 1, j)$ 開下來，或從左邊的 $(i, j - 1)$
開過來。 從上面來的話，前半段一定是開到 $(i - 1, j)$ 的最省走法（否則換成更
省的前半段，整條就更省），總耗電是
\texttt{best[i - 1][j]} 加上這一步的
$\max(0,\ h_{i,j} - h_{i-1,j})$。 從左邊來也一樣，兩者取小的：
\[
\mathrm{best}(i, j) = \min\Bigl(\mathrm{best}(i - 1, j) + c(h_{i-1,j}, h_{i,j}),\
\mathrm{best}(i, j - 1) + c(h_{i,j-1}, h_{i,j})\Bigr),
\]
其中 $c(a, b) = \max(0, b - a)$ 是從高度 $a$ 開到高度 $b$ 的耗電量。 和第~\ref{sec:dp2-two}~
小節的計數只差兩處：加法換成取最小值，以及每個來源要先加上那一步自己的成本。
這一步的成本由「從哪一格來」決定，所以從上面來和從左邊來，要用不同的高度算。

\textbf{初始值}：\texttt{best[0][0]} $= 0$，起點本身不耗電。

\textbf{計算順序}：一列一列由上往下，每一列由左往右。

範例的表格（右邊的橘色是最省的那條路線）：

\begin{center}
\begin{tikzpicture}[x=0.7cm, y=0.7cm,
    g/.style={draw=black!55, minimum size=0.7cm, inner sep=0pt, font=\small\ttfamily},
    on/.style={g, draw=memStroke, thick, fill=memFill}]
\begin{scope}
\node[font=\small] at (1.5,1) {高度 $h$};
\foreach \row [count=\r from 0] in {{1,5,2,2},{3,3,4,1},{6,2,2,0}} {
  \foreach \v [count=\c from 0] in \row {
    \node[g] at (\c,-\r) {\v};
  }
}
\end{scope}
\begin{scope}[xshift=4.2cm]
\node[font=\small] at (1.5,1) {\texttt{best}};
\foreach \row [count=\r from 0] in {{0,4,4,4},{2,2,3,3},{5,2,2,2}} {
  \foreach \v [count=\c from 0] in \row {
    \node[g] at (\c,-\r) {\v};
  }
}
\node[on] at (0,0) {0};
\node[on] at (0,-1) {2};
\node[on] at (1,-1) {2};
\node[on] at (1,-2) {2};
\node[on] at (2,-2) {2};
\node[on] at (3,-2) {2};
\end{scope}
\end{tikzpicture}
\diagcaption{範例的高度和填好的 \texttt{best}。 例如 \texttt{best[1][2]}：從上面
來是 $4 + \max(0, 4 - 2) = 6$，從左邊來是 $2 + \max(0, 4 - 3) = 3$，取 $3$。}
\label{fig:dp2-uphill}
\end{center}

同樣的模型也出現在別的地方：規劃自行車路線時只在乎總爬升；物流車在單行道組
成的路網上，每一段路的過路費不同。 只要移動的方向不會繞回頭，成本都可以這樣
一格一格累積。

\textbf{程式}

\begin{example}[label=ex:dp2-uphill]{cpp}{(uphill.cpp) 上坡才耗電的搬運車}
#include <algorithm>
#include <climits>
#include <iostream>
#include <vector>
using namespace std;

// 從高度 a 開到高度 b 的耗電量：只有上坡要耗電
long long climb(long long a, long long b) {
    return max(0LL, b - a);
}

int main() {
    ios::sync_with_stdio(false);
    cin.tie(nullptr);
    int m, n;
    cin >> m >> n;
    vector<vector<long long>> h(m, vector<long long>(n));
    for (auto &row : h) {
        for (long long &x : row) {
            cin >> x;
        }
    }
    // best[i][j]：從 (0, 0) 開到 (i, j) 最少耗多少電
    const long long INF = LLONG_MAX;     // 還沒找到任何走法
    vector<vector<long long>> best(m, vector<long long>(n, INF));
    best[0][0] = 0;
    for (int i = 0; i < m; i++) {
        for (int j = 0; j < n; j++) {
            if (i > 0) {                    // 從上面那一格開下來
                long long c = climb(h[i - 1][j], h[i][j]);
                best[i][j] = min(best[i][j], best[i - 1][j] + c);
            }
            if (j > 0) {                    // 從左邊那一格開過來
                long long c = climb(h[i][j - 1], h[i][j]);
                best[i][j] = min(best[i][j], best[i][j - 1] + c);
            }
        }
    }
    cout << best[m - 1][n - 1] << '\n';
    return 0;
}
\end{example}

\iosample{
3 4 \\
1 5 2 2 \\
3 3 4 1 \\
6 2 2 0 \\
}{
2 \\
}

\begin{itemize}
\item \texttt{best} 一開始全部是 \texttt{INF}（\texttt{LLONG\_MAX}），意思是「還沒有找到任何走
法」，兩個方向各拿來比一次。 \texttt{INF} 只當比較的對象，從來不拿去加：
被加的 \texttt{best[i - 1][j]}、\texttt{best[i][j - 1]} 都已經算好，每一格都
到得了，一定是真正的耗電量。 如果題目有走不進去的格子，到不了的格子會一直是
\texttt{INF}，拿它去加就會溢位，要先檢查（本章練習第~1~題）。
\item 起點在迴圈裡兩個 \texttt{if} 都不成立，保持 $0$。
\item 每一步最多耗電 $10^9$，一條路線有 $m + n - 2 \le 998$ 步，總和接近
$10^{12}$，超過 \texttt{int}。 高度、\texttt{climb} 的回傳值和表格都用
\texttt{long long}。
\end{itemize}

\textbf{複雜度}

每一格看兩個鄰居，$O(mn)$，$m = n = 500$ 時是 $25$ 萬格。 讀入 $25$ 萬個數字
比計算還花時間，最大的測試資料在本書測試的環境下約 $0.05$ 秒。

\subsection{常見錯誤}
\label{sec:dp2-uphill-err}

\begin{itemize}
\item \textbf{表格用 \texttt{int}}：輸入
\begin{code}{text}{總耗電超過 int 的輸入}
1 7
0 1000000000 0 1000000000 0 1000000000 0
\end{code}
要爬三次 $10^9$，正確答案是 \texttt{3000000000}。 把高度、\texttt{climb} 和
表格都改成 \texttt{int}（初始值用 \texttt{INT\_MAX}），在本書測試的環境下印出
\texttt{-1294967296}。
\item \textbf{下坡也算耗電}：\texttt{climb} 寫成兩個高度差的絕對值，範例印
出 \texttt{5}，上面貪心的反例印出 \texttt{7}。
\item \textbf{用貪心}：每一步選耗電少的、或比較低的格子，上面的反例印出
\texttt{8}。
\item \textbf{兩個方向用同一個高度差}：從左邊來時寫成
\texttt{climb(h[i - 1][j], h[i][j])}，抄上一行時忘了改。 從哪一格來，就用
哪一格的高度。
\end{itemize}

\section{滾動陣列}
\label{sec:dp2-rolling}

範例~\ref{ex:dp2-uphill} 用了兩張 $500 \times 500$ 的表，各 $2$ MB，沒有問題。
如果格子是 $5000 \times 5000$，計算只要 $2500$ 萬次，很快；可是一張
\texttt{long long} 的表就要 $2500$ 萬 $\times 8$ 位元組 $= 200$ MB，很多題目的
記憶體限制放不下。

看轉移式：算第 $i$ 列只用到第 $i - 1$ 列和第 $i$ 列自己。 第 $i - 2$ 列以前
的值，以後再也用不到。 像抄筆記時只看前一頁，再往前翻的那些頁可以先收起來。
所以只要留一列。 這種寫法叫\textbf{滾動陣列}（rolling array）：同一個一維陣
列，算完一列就「滾」到下一列。 第~9~章 9.4 節的招牌燈程式之後提過「只記上一格和這一格」，
就是一維的版本。

一列的陣列 \texttt{best[0..n-1]} 怎麼兼顧「上一列」和「這一列」？ 由左往右
填第 $i$ 列，填到第 $j$ 格時：

\begin{itemize}
\item \texttt{best[j]} 還沒被覆寫，存的是上一列的值，就是「上面那一格」。
\item \texttt{best[j - 1]} 剛剛填好，是這一列的值，就是「左邊那一格」。
\end{itemize}

兩個需要的值剛好都在，填好之後把 \texttt{best[j]} 換成新的值。

\begin{center}
\begin{tikzpicture}[x=1.1cm, y=1cm,
    g/.style={draw=black!55, minimum width=1.1cm, minimum height=0.7cm, inner sep=0pt, font=\small\ttfamily},
    new/.style={g, draw=memStroke, thick, fill=memFill},
    old/.style={g, fill=black!8},
    cur/.style={g, draw=optStroke, very thick}]
\node[new] (a0) at (0,0) {2};
\node[new] (a1) at (1,0) {2};
\node[cur] (a2) at (2,0) {4};
\node[old] (a3) at (3,0) {4};
\foreach \j in {0,...,3} { \node[idxlbl] at (\j,-0.6) {\texttt{best[\j]}}; }
\draw[decorate, decoration={brace, amplitude=5pt, mirror}] (-0.5,-0.9) -- (1.5,-0.9)
    node[midway, below=5pt, font=\small] {這一列（第 $1$ 列）};
\draw[decorate, decoration={brace, amplitude=5pt, mirror}] (1.5,-0.9) -- (3.5,-0.9)
    node[midway, below=5pt, font=\small] {上一列（第 $0$ 列）};
\draw[diagarrow] (a1.north) to[bend left=40] node[above, font=\small] {左：$2 + 1$} (a2.north);
\draw[diagarrow] (a2.north east) to[out=60, in=120, looseness=6] node[above, font=\small] {上：$4 + 2$} (a2.north west);
\node[font=\small, right] at (4.0,0) {$\Rightarrow$ \texttt{best[2]} $= \min(3, 6) = 3$};
\end{tikzpicture}
\diagcaption{範例~\ref{ex:dp2-uphill} 的輸入，填第 $1$ 列第 $2$ 格之前的陣列。
橘色已經是第 $1$ 列的值，灰色還是第 $0$ 列的值；粉紅框的格子同時是「上面那一
格」的舊值，也是等一下要寫入新值的位置。}
\label{fig:dp2-roll}
\end{center}

下面的程式連高度也一列一列讀：只留上一列和這一列的高度，讀完的列就丟掉。 整
支程式的記憶體是 $O(n)$，和列數 $m$ 無關。

\begin{example}[label=ex:dp2-roll]{cpp}{(uphill\_roll.cpp) 只留一列的搬運車}
#include <algorithm>
#include <climits>
#include <iostream>
#include <vector>
using namespace std;

long long climb(long long a, long long b) {
    return max(0LL, b - a);
}

int main() {
    ios::sync_with_stdio(false);
    cin.tie(nullptr);
    int m, n;
    cin >> m >> n;
    vector<long long> above(n), here(n);   // 上一列和這一列的高度
    vector<long long> best(n);             // 只留一列的表格
    for (int i = 0; i < m; i++) {
        for (long long &x : here) {
            cin >> x;
        }
        for (int j = 0; j < n; j++) {
            // 此時 best[j] 是上一列的值，best[j - 1] 是這一列的值
            long long b = LLONG_MAX;
            if (i == 0 && j == 0) {
                b = 0;
            }
            if (i > 0) {
                b = min(b, best[j] + climb(above[j], here[j]));
            }
            if (j > 0) {
                long long c = climb(here[j - 1], here[j]);
                b = min(b, best[j - 1] + c);
            }
            best[j] = b;
        }
        swap(above, here);    // 這一列的高度，變成下一輪的「上一列」
    }
    cout << best[n - 1] << '\n';
    return 0;
}
\end{example}

\iosample{
3 2 \\
5 1 \\
4 9 \\
8 6 \\
}{
4 \\
}

\begin{itemize}
\item 輸入是上面貪心的反例，答案 $4$。
\item 新的值先算在 \texttt{b} 裡，兩個方向都比完才寫回 \texttt{best[j]}。 如
果先寫 \texttt{best[j] = ...} 再拿 \texttt{best[j]} 去比，就把上一列的值蓋掉
了。
\item \texttt{swap(above, here)} 交換兩個 \texttt{vector} 的內容，只交換內部
的指標，不會一個一個複製，是 $O(1)$（銜接書第~8~章）。 交換後，\texttt{here}
裡是更早的那一列，下一輪讀入時直接蓋掉。
\end{itemize}

這支程式也用完整的測試資料驗證過，和範例~\ref{ex:dp2-uphill} 的答案完全相同。

\textbf{填的方向不能反過來。} 把 \texttt{j} 改成由大到小，填第 $j$ 格時
\texttt{best[j - 1]} 還是上一列的值，「左邊那一格」拿錯了。 第 $0$ 列更糟：
\texttt{best[j - 1]} 是一開始的 $0$。 範例~\ref{ex:dp2-uphill} 的輸入會印出
\texttt{0}。 寫滾動陣列之前，先把「填第 $j$ 格時，陣列裡每一格是哪一列的值」
寫下來，再決定方向。

\textbf{什麼時候不能只留一列。}

\begin{itemize}
\item 要印出路線本身，不只是答案：一列的陣列沒有留下「每一格是從哪裡來的」。
可以另外開一張 \texttt{char} 的表記錄方向，每格只要 $1$ 位元組，比
\texttt{long long} 小八倍。
\item 轉移用到更前面的列，例如第 $i - 2$ 列：就留兩列或三列。
\item 轉移用到左上角 $(i - 1, j - 1)$：它的舊值在填 \texttt{best[j - 1]} 時已經被
蓋掉了，要先用一個變數存起來（本章練習第~6~題）。
\end{itemize}

不確定覆寫順序對不對時，改用兩個一維陣列 \texttt{prev} 和 \texttt{cur}：用
\texttt{prev} 算出整列 \texttt{cur}，算完再 \texttt{swap(prev, cur)}。 記憶體
多一列，但不會有「讀到一半被蓋掉」的問題。

\begin{apcsnote}
題目會寫出輸入的範圍和記憶體限制。 開二維表格之前先算一次：格子數乘上每格的
位元組數（\texttt{int} 是 $4$，\texttt{long long} 是 $8$），再乘上表格的張
數。 超過限制或接近限制，就改成滾動陣列；離限制很遠，用整張表比較好寫，也比
較好檢查。
\end{apcsnote}

\section{一格記兩個數字}
\label{sec:dp2-states}

第~9~章的大夜班，只記「前 $i$ 天最多多少」不夠，還要記第 $i$ 天在做什麼。
格子也一樣：只知道「走到這一格有幾條路線」，有時候不夠回答題目要的問題。 這
一節的例題要知道路線「走到這裡為止，已經做過某件事了沒有」。

\subsection{例題：至少經過一台飲水機的慢跑路線}
\label{sec:dp2-ex-fountain}

\textbf{題意}（\akiref{jogging-routes-past-a-fountain}{至少經過一台飲水機的慢跑路線}）

\begin{itemize}
\item 河濱步道區是 $H$ 列 $W$ 行的格子。 \texttt{.} 是步道，\texttt{*} 是有
飲水機的步道，\texttt{\#} 是花圃，不能踩。
\item 從左上角跑到右下角，每一步往右或往下。 只算\textbf{至少經過一台飲水機}
的路線，起點、終點是 \texttt{*} 也算經過。
\item 輸出路線數除以 $10^9 + 7$ 的餘數。 $H, W \le 500$。
\end{itemize}

例如
\begin{code}{text}{3 列 4 行的步道區}
....
.#*.
*...
\end{code}
繞過花圃的路線有 $4$ 條。 一路往右再往下的那條沒有經過飲水機，其他 $3$ 條各
經過 $(2, 3)$ 或 $(3, 1)$ 的飲水機，答案是 $3$。

\textbf{想法}

先看一個聽起來合理、但會錯的做法：對每一台飲水機，算出經過它的路線數，再全部
加起來。 問題在於一條路線可能經過好幾台飲水機。 例如
\begin{code}{text}{兩台飲水機在同一條路線上}
.**
...
\end{code}
從左上到右下有 $3$ 條路線：右右下經過兩台，右下右經過一台，下右右一台都沒經
過，答案是 $2$。 照上面的做法，經過 $(1, 2)$ 的有 $2$ 條、經過 $(1, 3)$ 的有
$1$ 條，加起來是 $3$，右右下被算了兩次。

\textbf{做法一：反過來數。} 「至少經過一台」的反面是「一台都沒經過」。 先算
全部的路線數，再把每台飲水機都當成花圃、算一次不經過飲水機的路線數，相減就
是答案。 兩次都是第~\ref{sec:dp2-blocked}~小節的計數，只是要注意相減時的餘
數（見本節的常見錯誤）。

\textbf{做法二：一格記兩個數。} 「走到 $(i, j)$ 的路線」分成兩組：一路上還沒
經過飲水機的，和已經經過的。 兩組之後的命運不同：前一組還要再找飲水機才算
數，後一組已經滿足條件了。 所以每一格記兩個數：

\begin{itemize}
\item \texttt{no[i][j]}：走到 $(i, j)$、包含這一格在內都還沒經過飲水機的路線數。
\item \texttt{yes[i][j]}：走到 $(i, j)$、包含這一格在內已經經過至少一台的路線數。
\end{itemize}

轉移時，先像平常一樣把北邊和西邊加起來，兩組分開加。 接著看這一格：

\begin{itemize}
\item 是 \texttt{.}：兩組都不變。
\item 是 \texttt{*}：走進這一格之後，原本「還沒經過」的路線也經過了。 所以
\texttt{yes} 是兩組加起來，\texttt{no} 變成 $0$。
\item 是 \texttt{\#}：走不進來，兩組都是 $0$。
\end{itemize}

答案是 \texttt{yes[H][W]}。 起點的「走進來之前」是：還沒經過的有 $1$ 條，已
經經過的有 $0$ 條；接著照起點自己的種類處理，起點是 \texttt{*} 就自動變成
\texttt{yes} $= 1$。

下表是範例每一格的（\texttt{no}, \texttt{yes}）。 $(2, 3)$ 是飲水機，北邊和
西邊進來的 $(1, 0)$、$(0, 0)$ 合起來全部變成 \texttt{yes}；$(3, 1)$ 也一樣。

\begin{center}
\begin{tabular}{c|cccc}
\toprule
 & 第 1 行 & 第 2 行 & 第 3 行 & 第 4 行 \\
\midrule
第 1 列 & $(1, 0)$ & $(1, 0)$ & $(1, 0)$ & $(1, 0)$ \\
第 2 列 & $(1, 0)$ & \texttt{\#} & $(0, 1)$ \texttt{*} & $(1, 1)$ \\
第 3 列 & $(0, 1)$ \texttt{*} & $(0, 1)$ & $(0, 2)$ & $(1, 3)$ \\
\bottomrule
\end{tabular}
\end{center}

右下角是 $(1, 3)$：沒經過飲水機的 $1$ 條，經過的 $3$ 條，答案是 $3$。 兩個數
加起來是 $4$，就是不管飲水機時的總路線數。

做法一只能回答「至少一台」這種剛好有反面可以數的問題。 做法二可以推廣，例如
「剛好經過一台」或「至少經過兩台」，只要多分幾組（本章練習第~7~題）。 下面的
程式用做法二。

同樣的模型也出現在別的地方：外送員的路線至少要經過一個取餐點、巡邏路線至少
要經過一個打卡點，都是「經過了沒有」這一個是非的資訊。

\textbf{程式}

\begin{example}[label=ex:dp2-fountain]{cpp}{(fountain.cpp) 至少經過一台飲水機的慢跑路線}
#include <iostream>
#include <string>
#include <vector>
using namespace std;

const long long MOD = 1000000007;

int main() {
    ios::sync_with_stdio(false);
    cin.tie(nullptr);
    int h, w;
    cin >> h >> w;
    vector<string> g(h);
    for (string &row : g) {
        cin >> row;
    }
    // no[i][j]：走到 (i, j)、一路上還沒經過飲水機的路線數
    // yes[i][j]：走到 (i, j)、已經經過至少一台飲水機的路線數
    // 列、行從 1 編號，第 0 列、第 0 行是多開的一圈，都是 0
    vector<vector<long long>> no(h + 1, vector<long long>(w + 1));
    vector<vector<long long>> yes = no;
    for (int i = 1; i <= h; i++) {
        for (int j = 1; j <= w; j++) {
            char cell = g[i - 1][j - 1];
            if (cell == '#') {
                continue;              // 花圃：兩種都是 0
            }
            long long inNo, inYes;     // 走進這一格之前
            if (i == 1 && j == 1) {
                inNo = 1;              // 起點：還沒經過任何一格
                inYes = 0;
            } else {
                inNo = (no[i - 1][j] + no[i][j - 1]) % MOD;
                inYes = (yes[i - 1][j] + yes[i][j - 1]) % MOD;
            }
            if (cell == '*') {         // 走進飲水機：都算經過了
                no[i][j] = 0;
                yes[i][j] = (inNo + inYes) % MOD;
            } else {
                no[i][j] = inNo;
                yes[i][j] = inYes;
            }
        }
    }
    cout << yes[h][w] << '\n';
    return 0;
}
\end{example}

\iosample{
3 4 \\
.... \\
.\#*. \\
*... \\
}{
3 \\
}

\begin{itemize}
\item \texttt{vector<vector<long long>> yes = no;} 複製一份同樣大小、全部是
$0$ 的表格。
\item 花圃直接 \texttt{continue}，兩張表的那一格保持一開始的 $0$。 這裡可以這
樣寫，是因為每一格只會被填一次；換成滾動陣列時就不行了，本章練習第~3~題會看到
原因。
\item 兩個小於 $10^9 + 7$ 的數相加，不超過 $2.1 \times 10^9$，\texttt{long
long} 放得下，每次加完就取餘數。
\item 起點是花圃時，\texttt{continue} 讓起點的兩個數都是 $0$，後面每一格也都
是 $0$，答案自然是 $0$。
\end{itemize}

\textbf{複雜度}

每一格做常數次運算，$O(HW)$，兩張表各 $25$ 萬格。 最大的測試資料在本書測試
的環境下約 $0.02$ 秒。

\subsection{常見錯誤}
\label{sec:dp2-fountain-err}

\begin{itemize}
\item \textbf{反過來數，但相減後沒有處理負數}：兩個數都先取過餘數，前面的可能
比後面的小。 例如 $18 \times 18$ 的格子全是 \texttt{.}，只有 $(1, 3)$ 是
\texttt{*}，全部的路線數取餘數後是 $333606206$，避開飲水機的是 $767883493$。
寫成 \texttt{(all - none) \% MOD}，在本書測試的環境下印出
\texttt{-434277287}；正確答案是 \texttt{565722720}。 C++ 的 \texttt{\%} 對
負數的結果是負的，要寫成 \texttt{(all - none + MOD) \% MOD}。
\item \textbf{把經過每台飲水機的路線數加起來}：前面的 $2 \times 3$ 例子會得到
$3$，正確答案是 $2$。
\item \textbf{起點的飲水機沒算到}：把起點寫成「\texttt{no} $= 1$ 之後直接
\texttt{continue}」，跳過了後面判斷 \texttt{*} 的那幾行。 輸入
\begin{code}{text}{起點就有飲水機的輸入}
3 3
*..
.#.
...
\end{code}
正確答案是 \texttt{2}，這樣寫會印出 \texttt{0}。 起點也要經過「看這一格是什
麼」那一步，只有「走進來之前」的值要特別給。
\end{itemize}

\subsection{格子加上一點資訊}

這一節的狀態是「位置加上一個是非」。 一般來說，二維格子的狀態可以寫成「位置
$(i, j)$，再加上一點資訊」，這點資訊要回答一個問題：從這裡開始，之後怎麼走才
合規定、能得幾分，除了位置，還取決於過去的什麼？ 例題取決於「經過飲水機了沒
有」，所以多記一個是非；如果取決於「已經經過幾台，最多算到 $K$ 台」，就多記
$0$ 到 $K$ 這 $K + 1$ 種。

估計時間的方法和第~9~章一樣：狀態數是 $H \times W \times$（資訊的種類數），
再乘上每個狀態要看幾個前一狀態。 $500 \times 500$ 的格子，每格記十幾種資
訊，仍然在幾百萬的數量級。 如果資訊多到要記「走過哪些格子」，種類數就變成指
數，表格開不下，這時要回到第~4~章的回溯，只能處理很小的格子。

\section{最大正方形}
\label{sec:dp2-square}

前面幾節的表格，一格存的是「走到這裡」的答案。 這一節換一種問法：每一格當作
某個圖形的右下角，存「以這一格為右下角，最大能做到多大」。 這種狀態在找正方
形時特別好用。

\subsection{例題：單一花種的取景框}
\label{sec:dp2-ex-frames}

\textbf{題意}（\akiref{single-flower-photo-frames}{單一花種的取景框}）

\begin{itemize}
\item 花圃是 $n$ 列 $m$ 行的格子，每格是一個小寫字母，代表花的種類。
\item 邊長 $k$ 的正方形取景框對齊格線、整個在花圃裡；框裡的 $k^2$ 格全是同一
種花才合格。
\item 有 $q$ 個問題，每個問題給一個 $k$，回答有幾個位置可以放下合格的邊長 $k$
取景框。
\item $n, m \le 500$，$q \le 2 \times 10^5$，$1 \le k \le 500$。
\end{itemize}

例如
\begin{code}{text}{4 列 5 行的花圃}
aabbb
aabbb
cabbb
ccaaa
\end{code}
邊長 $1$ 有 $20$ 個位置，邊長 $2$ 有 $5$ 個（左上角的 \texttt{a} 一個，右邊
\texttt{b} 的 $3 \times 3$ 裡四個），邊長 $3$ 有 $1$ 個，邊長 $4$ 沒有。

\textbf{想法}

最直接的做法：每個問題枚舉框的每個位置，再檢查框裡的 $k^2$ 格。 $n = m =
500$、$k = 250$ 時，位置有 $251^2$ 個，每個檢查 $250^2$ 格，一個問題就大約
$4 \times 10^9$ 次；$q$ 個問題更不用說。 就算把每個 $k$ 的答案事先算好存起來，$k$ 從 $1$ 到
$500$ 全部檢查一遍也來不及。

\textbf{換一個問法。} 對每一格 $(i, j)$，問「以它為右下角、全是同一種花的正
方形，邊長最大是多少」，記為 \texttt{side[i][j]}。 如果
\texttt{side[i][j]} $= s$，以 $(i, j)$ 為右下角、邊長 $k \le s$ 的正方形都包
在那個 $s \times s$ 的正方形裡，一定合格；邊長 $k > s$ 的一定不合格。 所以這
一格剛好替邊長 $1, 2, \ldots, s$ 各貢獻一個位置，邊長 $k$ 的答案就是
\[
\text{\texttt{side}} \ge k \text{ 的格子數}.
\]
先數出 \texttt{side} 剛好等於每個 $d$ 的格子數 \texttt{cnt[d]}，再從大到小
累加，\texttt{cnt[k]} 就變成「至少是 $k$ 的格子數」，每個問題查一次表。

剩下的問題是怎麼算 \texttt{side}。 先看一個會錯的想法：只看上面和左邊兩格。
如果這一格和上面、左邊的花都一樣，就取兩者的最小值加一。 反例：
\begin{code}{text}{只看兩個鄰居會錯的花圃}
ba
aa
\end{code}
右下角的 \texttt{a}，上面和左邊都是 \texttt{a}、\texttt{side} 都是 $1$，照這
個規則得到 $2$；可是 $2 \times 2$ 的框裡左上角是 \texttt{b}，不合格。 上面和
左邊兩個正方形合起來，剛好缺了左上角那一塊。

正確的規則要看三個鄰居：上面、左邊、左上。 如果這一格和三個鄰居的花都相同，
\[
\text{\texttt{side[i][j]}} = \min\bigl(\text{\texttt{side[i-1][j]}},\
\text{\texttt{side[i][j-1]}},\ \text{\texttt{side[i-1][j-1]}}\bigr) + 1,
\]
否則是 $1$（只有自己這一格）。 第一列、第一行沒有完整的三個鄰居，也是 $1$。

直覺是下圖：三個鄰居各自的正方形，加上這一格自己，剛好蓋住一個大一號的正方
形。 三個裡面只要有一個小，大正方形就會缺一塊，所以由最小的那一個決定。

\begin{center}
\begin{tikzpicture}[x=0.55cm, y=0.55cm]
\draw[black!25] (0,0) grid (5,5);
\fill[memFill] (4,0) rectangle (5,1);
\draw[very thick, coreStroke] (2,1) rectangle (5,4);
\draw[very thick, optStroke, dashed] (1,0) rectangle (4,3);
\draw[very thick, greenStroke, dotted] (1,1) rectangle (4,4);
\draw[ultra thick, memStroke] (1,0) rectangle (5,4);
\node[font=\small] at (4.5,0.5) {$\ast$};
\node[font=\small, right, align=left] at (5.6,3.5) {\textcolor{coreStroke}{實線}：上面那一格的正方形};
\node[font=\small, right, align=left] at (5.6,2.5) {\textcolor{optStroke}{虛線}：左邊那一格的正方形};
\node[font=\small, right, align=left] at (5.6,1.5) {\textcolor{greenStroke}{點線}：左上那一格的正方形};
\node[font=\small, right, align=left] at (5.6,0.5) {橘色粗框：加上 $\ast$ 之後的 $4 \times 4$};
\end{tikzpicture}
\diagcaption{三個鄰居的 \texttt{side} 都是 $3$ 時，三個 $3 \times 3$ 正方形加上
$\ast$ 這一格，蓋滿了以 $\ast$ 為右下角的 $4 \times 4$。}
\label{fig:dp2-square}
\end{center}

\begin{whybox}[三個取最小值為什麼剛好]
設 $(i, j)$ 和三個鄰居的花都是 $c$，三個鄰居的 \texttt{side} 最小值是 $t$。

\textbf{至少是 $t + 1$。} 以三個鄰居為右下角、邊長 $t$ 的正方形都合格（它們
的 \texttt{side} 都 $\ge t$），而且分別全是那一格的花，也就是 $c$。 以
$(i, j)$ 為右下角、邊長 $t + 1$ 的正方形裡，除了 $(i, j)$ 自己以外的每一格，
都落在這三個正方形的至少一個裡面：在最後一列的落在左邊那個裡，在最後一行的
落在上面那個裡，其他的落在左上那個裡。 所以整個正方形都是 $c$。

\textbf{不會超過 $t + 1$。} 假設以 $(i, j)$ 為右下角、邊長 $s$ 的正方形合格。
把它的最後一列去掉，得到以 $(i - 1, j)$ 為右下角、邊長 $s - 1$ 的正方形（少
了一列，寬度也只取右邊的 $s - 1$ 行），它是大正方形的一部分，也合格，所以
\texttt{side[i-1][j]} $\ge s - 1$。 左邊和左上同理。 因此三個都 $\ge s - 1$，
$t \ge s - 1$，也就是 $s \le t + 1$。

這一格和三個鄰居只要有一格的花不同，任何邊長 $2$ 以上的正方形都會同時框到兩
種花，所以只能是 $1$。
\end{whybox}

範例的 \texttt{side}：

\begin{center}
\begin{tikzpicture}[x=0.62cm, y=0.62cm,
    g/.style={draw=black!55, minimum size=0.62cm, inner sep=0pt, font=\small\ttfamily},
    hl/.style={g, draw=memStroke, thick, fill=memFill}]
\begin{scope}
\node[font=\small] at (2,1) {花圃};
\foreach \row [count=\r from 0] in {{a,a,b,b,b},{a,a,b,b,b},{c,a,b,b,b},{c,c,a,a,a}} {
  \foreach \v [count=\c from 0] in \row {
    \node[g] at (\c,-\r) {\v};
  }
}
\end{scope}
\begin{scope}[xshift=4.6cm]
\node[font=\small] at (2,1) {\texttt{side}};
\foreach \row [count=\r from 0] in {{1,1,1,1,1},{1,2,1,2,2},{1,1,1,2,3},{1,1,1,1,1}} {
  \foreach \v [count=\c from 0] in \row {
    \node[g] at (\c,-\r) {\v};
  }
}
\node[hl] at (4,-2) {3};
\end{scope}
\end{tikzpicture}
\diagcaption{範例的花圃和 \texttt{side}。 第 $2$ 列第 $1$ 行（從 $0$ 編號）的
\texttt{a} 和上面、左上都是 \texttt{a}，但左邊是 \texttt{c}，所以是 $1$。}
\label{fig:dp2-side}
\end{center}

\texttt{side} 等於 $1, 2, 3$ 的格子分別有 $15, 4, 1$ 格。 從大到小累加：至少
$3$ 的有 $1$ 格，至少 $2$ 的有 $5$ 格，至少 $1$ 的有 $20$ 格，就是範例的答案。

同樣的模型也出現在別的地方：在一塊有些地方不能用的空地上，找能蓋最大正方形
建物的位置；在影像裡找顏色一致的最大方塊。

\textbf{程式}

\begin{example}[label=ex:dp2-frames]{cpp}{(frames.cpp) 單一花種的取景框}
#include <algorithm>
#include <iostream>
#include <string>
#include <vector>
using namespace std;

int main() {
    ios::sync_with_stdio(false);
    cin.tie(nullptr);
    int n, m;
    cin >> n >> m;
    vector<string> g(n);
    for (string &row : g) {
        cin >> row;
    }
    // side[i][j]：右下角在 (i, j)、全是同一種花的正方形的最大邊長
    vector<vector<int>> side(n, vector<int>(m, 1));
    vector<long long> cnt(502, 0);  // cnt[d]：side 剛好是 d 的格子數
    for (int i = 0; i < n; i++) {
        for (int j = 0; j < m; j++) {
            char c = g[i][j];
            if (i > 0 && j > 0 && c == g[i - 1][j] &&
                c == g[i][j - 1] && c == g[i - 1][j - 1]) {
                side[i][j] = min({side[i - 1][j], side[i][j - 1],
                                  side[i - 1][j - 1]}) + 1;
            }
            cnt[side[i][j]]++;
        }
    }
    // 由大到小累加：cnt[k] 變成 side 至少是 k 的格子數
    for (int k = 500; k >= 1; k--) {
        cnt[k] += cnt[k + 1];
    }
    int q;
    cin >> q;
    while (q--) {
        int k;
        cin >> k;
        cout << cnt[k] << '\n';
    }
    return 0;
}
\end{example}

\iosample{
4 5 \\
aabbb \\
aabbb \\
cabbb \\
ccaaa \\
4 \\
1 2 3 4 \\
}{
20 \\
5 \\
1 \\
0 \\
}

\begin{itemize}
\item \texttt{side} 一開始全是 $1$，條件不成立的格子（包括第一列、第一行）就
保持 $1$。
\item \texttt{min(\{a, b, c\})} 用大括號一次比三個數，需要
\texttt{<algorithm>}。
\item \texttt{side} 最大是 $500$，累加時會讀到 \texttt{cnt[501]}，所以開
$502$ 格。 $k$ 比 $n$ 或 $m$ 大時，\texttt{cnt[k]} 自然是 $0$，不用特別判斷。
\item 輸出最多 $2 \times 10^5$ 行，用 \texttt{'\textbackslash n'} 而不用
\texttt{endl}，避免每一行都清一次輸出緩衝區。
\end{itemize}

\textbf{複雜度}

填表 $O(nm)$，累加 $O(500)$，每個問題 $O(1)$，總共 $O(nm + q)$。 最大的測試資
料在本書測試的環境下約 $0.06$ 秒。

\subsection{常見錯誤}

\begin{itemize}
\item \textbf{只數剛好等於 $k$ 的格子}：少了累加那個迴圈，範例印出
\texttt{15 4 1 0}（每個數一行）。 右邊那塊 \texttt{b} 的 \texttt{side} 是
$3$，它裡面以同一格為右下角的 $2 \times 2$ 也合格，卻沒被算到。
\item \textbf{只檢查上面和左邊的花}：條件裡少了 \texttt{g[i - 1][j - 1]}，
上面的 \texttt{ba}／\texttt{aa} 問邊長 $2$ 時印出 \texttt{1}，正確是 \texttt{0}。
\item \textbf{每個問題重新掃一遍}：每個問題都把 \texttt{side} 整張表掃一次，
數有幾格 $\ge k$，一個問題 $25$ 萬格，$2 \times 10^5$ 個問題就是 $5 \times
10^{10}$，來不及。 問題很多時，先把所有可能的答案算好。
\end{itemize}

\section{從正方形到矩形}
\label{sec:dp2-rect}

矩形比正方形多一個自由度：寬和高可以不同。 上一節「以這一格為右下角的最大正
方形」只要一個數字就描述得完，矩形不行，所以要換別的方法。 這一節整理三種常
見的工具，章末的實作練習會用到。

\subsection{二維前綴和}
\label{sec:dp2-prefix2d}

第~5~章的前綴和推廣到二維：$S(i, j)$ 是左上角 $i$ 列 $j$ 行那一塊的總和，
$S(0, \cdot) = S(\cdot, 0) = 0$。 它本身就是一張由上往下、由左往右填的表：
\[
S(i, j) = S(i - 1, j) + S(i, j - 1) - S(i - 1, j - 1) + a_{i,j}.
\]
上面那塊和左邊那塊重疊的部分（左上那塊）被加了兩次，所以減掉一次。 有了 $S$，
第 $r_1$ 到 $r_2$ 列、第 $c_1$ 到 $c_2$ 行的矩形總和用同樣的想法算：
\[
S(r_2, c_2) - S(r_1 - 1, c_2) - S(r_2, c_1 - 1) + S(r_1 - 1, c_1 - 1).
\]
\begin{code}{cpp}{二維前綴和}
// a 和 S 都從 1 編號，S 多開一圈 0
for (int i = 1; i <= R; i++) {
    for (int j = 1; j <= C; j++) {
        S[i][j] = S[i-1][j] + S[i][j-1] - S[i-1][j-1] + a[i][j];
    }
}
// 第 r1 到 r2 列、第 c1 到 c2 行的總和
long long rect = S[r2][c2] - S[r1-1][c2]
               - S[r2][c1-1] + S[r1-1][c1-1];
\end{code}
有了它，每個矩形的總和都是 $O(1)$。 枚舉所有矩形是 $O(R^2 C^2)$，$R = C = 50$
時大約 $6 \times 10^6$ 個，來得及；$R = C = 200$ 就有 $16$ 億個，太多了。

\subsection{總和最大的矩形}
\label{sec:dp2-maxrect}

要找總和最大的矩形，先固定它的上邊在第 $a$ 列、下邊在第 $b$ 列。 這時每一行
在第 $a$ 到第 $b$ 列的部分，可以壓成一個數：這一行在這幾列的總和。 矩形就變
成這排數字裡連續的一段，問題變成一維的最大區間和，用第~9~章 9.6 節的 \texttt{endAt}（或第~5~章的前綴和）一次掃過去。

例如
\[
\begin{matrix}
1 & -2 & -3 & 2 \\
-4 & 3 & 2 & -1 \\
2 & 1 & 4 & -5
\end{matrix}
\]
固定第 $2$ 到第 $3$ 列，四行壓成 $-2, 4, 6, -6$，最大的連續一段是 $4 + 6 =
10$，也就是第 $2$ 到第 $3$ 列、第 $2$ 到第 $3$ 行的矩形。 對每一組 $(a, b)$ 都
做一次，取最大的，這個例子的答案就是 $10$。

\begin{code}{cpp}{固定上下兩列，壓成一維}
long long best = LLONG_MIN;
for (int a = 1; a <= R; a++) {
    vector<long long> col(C + 1, 0);  // 第 j 行在第 a 到 b 列的總和
    for (int b = a; b <= R; b++) {
        long long endAt = 0;          // 以第 j 行結尾的最大總和
        for (int j = 1; j <= C; j++) {
            col[j] += A[b][j];        // b 往下一列，每行多加一格
            endAt = col[j] + max(endAt, 0LL);  // 前面是負的就不接
            best = max(best, endAt);
        }
    }
}
\end{code}
\texttt{b} 往下移一列時，\texttt{col} 不用重算，每一行加上新的一格就好。 時間
是 $O(R^2 C)$。 $R$ 和 $C$ 差很多時，讓比較小的那一邊當成被固定的「列」，例如
$R = 20$、$C = 10^5$ 是 $4 \times 10^7$，反過來就是 $2 \times 10^{11}$。

\subsection{全空的最大矩形}

另一種常見的題目：格子裡有些格子是壞的，找不含壞格子、面積最大的矩形。 逐列
處理，對第 $i$ 列的每一行，記下「從第 $i$ 列往上數，連續幾格是好的」，叫做
高度 \texttt{up[j]}。 \texttt{up} 本身是一個很簡單的動態規劃：好的格子是上一
列的 \texttt{up[j]} 加一，壞的格子是 $0$，可以用滾動陣列只留一列。

把第 $i$ 列的 \texttt{up} 畫成一排長條，以第 $i$ 列為底邊的矩形，就是這排長條
裡的矩形。 每一根長條往左、往右能延伸到第一根比它矮的長條為止，用第~6~章
6.3 節的單調堆疊，一列是 $O(C)$，整張表 $O(RC)$。 例如某一
列的高度是 $2, 3, 3, 1$，以高度 $3$ 為高的矩形可以橫跨第 $2$、$3$ 行，面積
$6$；以高度 $2$ 為高的可以橫跨前三行，面積也是 $6$。

\section{什麼時候用格子上的動態規劃}
\label{sec:dp2-when}

\begin{itemize}
\item 輸入是一張格子或表格，移動只朝固定的方向（往右和往下，或一列一列往下），
問「有幾種」「最少」「最多」。
\item 格子數在幾百萬以內，每一格只看常數個鄰居。 $500 \times 500$、
$1000 \times 1000$、$2000 \times 2000$ 都是常見的範圍。
\item 每一格的答案，只需要「怎麼走到這一格」的少量資訊，例如經過了沒有、用了幾
次；這些資訊放進狀態，表格多幾層。
\item 問正方形時，用「以這一格為右下角」當狀態；問矩形時，先想二維前綴和、固定
兩列壓成一維、逐列的高度三種工具。
\end{itemize}

看起來像、但不能直接用的情況：

\begin{itemize}
\item 可以往四個方向走，例如迷宮。 狀態圖有環（往右再往左就回到原地），找不到計
算順序，第~9~章說過的事在這裡一樣成立。 每一步的成本都一樣時，用第~12~章的
廣度優先搜尋；成本不同時，要用加權最短路徑（《C++ 演算法解題 2》第~8~章）。
\item 規定路線不能重複經過同一格，而且可以往回走：要記住走過哪些格子，狀態數
是指數，只有很小的格子可以用第~4~章的回溯。
\end{itemize}

\section{本章常見錯誤}
\label{sec:dp2-pitfalls}

各節的錯誤，歸納起來是五類。 寫完格子的動態規劃，交出去之前逐項檢查。

\begin{center}
\begin{tabular}{p{0.2\linewidth}p{0.46\linewidth}p{0.22\linewidth}}
\toprule
錯誤 & 例子 & 見 \\
\midrule
邊界與初始值 & 第一列、第一行直接設成 $1$，忽略障礙；起點的飲水機沒算；
多開一圈卻沒有減一 & 第~\ref{sec:dp2-blocked}、\ref{sec:dp2-border}、
\ref{sec:dp2-fountain-err}~節 \\
轉移漏了來源 & 正方形只看上、左兩格；兩個方向用同一個高度差 &
第~\ref{sec:dp2-ex-frames}、\ref{sec:dp2-uphill-err}~節 \\
狀態少了資訊 & 把經過每台飲水機的路線數加起來，重複計算；用貪心選下一步 &
第~\ref{sec:dp2-ex-uphill}、\ref{sec:dp2-ex-fountain}~節 \\
滾動陣列的覆寫順序 & 填的方向反了，用到還沒更新的值；先寫回再比較；障礙格
\texttt{continue} 留下上一列的值 & 第~\ref{sec:dp2-rolling}~節、練習第~3~題 \\
型別與餘數 & 耗電總和用 \texttt{int}；相減後沒有加回 $10^9 + 7$ &
第~\ref{sec:dp2-uphill-err}、\ref{sec:dp2-fountain-err}~節 \\
\bottomrule
\end{tabular}
\end{center}

另外一個和格子特別有關的習慣：自己測試時，一定要有一組列數和行數不同的輸入，
以及一組只有一列或只有一行的輸入。 列、行寫反和邊界沒處理好，用正方形的格子
都看不出來。

\section{本章練習}
\label{sec:dp2-exercises}

\begin{drills}[label=drills:dp2]
\item 每一格有一個過路費，起點和終點也要付，從左上角往右、往下走到右下角：
\[
\begin{matrix}
3 & 1 & 4 & 1 \\
2 & 7 & 1 & 8 \\
5 & 2 & 1 & 3 \\
1 & 6 & 2 & 1
\end{matrix}
\]
\begin{enumerate}
\item 寫出狀態、轉移和初始值，手動填出整張表，答案是多少？ 寫出一條最省的
路線。
\item 如果有些格子不能走，有人把不能走的格子設成 \texttt{LLONG\_MAX}，轉移照
樣寫 \texttt{min(best[i-1][j], best[i][j-1]) + a[i][j]}。 這樣會出什麼問題？
怎麼改？
\end{enumerate}

\item 用範例~\ref{ex:dp2-roll} 跑範例~\ref{ex:dp2-uphill} 的輸入（$3$ 列 $4$ 行
那一組）。
\begin{enumerate}
\item 第 $0$ 列算完時，\texttt{best} 的內容是什麼？
\item 第 $1$ 列剛算完 \texttt{j = 1} 那一格時，\texttt{best} 的內容是什麼？
哪幾格是第 $1$ 列的值，哪幾格還是第 $0$ 列的值？
\end{enumerate}

\item 有人把範例~\ref{ex:dp2-streets} 改成滾動陣列：
\begin{code}{cpp}{練習 3}
vector<long long> f(c + 1, 0);
for (int i = 1; i <= r; i++) {
    for (int j = 1; j <= c; j++) {
        if (g[i - 1][j - 1] == 'x') {
            continue;
        }
        if (i == 1 && j == 1) {
            f[j] = 1;
        } else {
            f[j] = f[j] + f[j - 1];
        }
    }
}
cout << f[c] << '\n';
\end{code}
範例~\ref{ex:dp2-streets} 的輸入會印出什麼？ 錯在哪裡？ 怎麼改？

\item 直接照轉移式寫成遞迴，不加記憶化：
\begin{code}{cpp}{練習 4}
long long paths(int i, int j) {
    if (i == 1 || j == 1) {
        return 1;
    }
    return paths(i - 1, j) + paths(i, j - 1);
}
\end{code}
\begin{enumerate}
\item 呼叫 \texttt{paths(4, 5)}，回傳值是多少？ \texttt{paths} 一共被呼叫幾
次（包括第一次）？
\item 說明一般來說，呼叫次數 $= 2 \times$ 回傳值 $- 1$。（提示：每一次呼叫不
是直接回傳 $1$，就是剛好再呼叫兩次。）
\item 照第~9~章 \texttt{stamps\_memo.cpp} 的方式加上記憶化（查表的呼叫也算一
次），\texttt{paths(4, 5)} 一共被呼叫幾次？
\end{enumerate}

\item 照範例~\ref{ex:dp2-frames} 的規則，填出下面花圃的 \texttt{side}，並回答
邊長 $1$ 到 $4$ 各有幾個合格的位置。
\begin{code}{text}{練習 5 的花圃}
aabba
aabbb
bbbbb
abbbb
aabbb
\end{code}

\item 第~\ref{sec:dp2-paths}~節的路線計數，如果除了往東、往南，還可以往東南斜
走一步（從 $(i, j)$ 直接到 $(i + 1, j + 1)$）：
\begin{enumerate}
\item 寫出新的轉移式。
\item 改成只留一列的滾動陣列，由左往右填。 為什麼需要多一個變數？ 寫出迴圈。
\end{enumerate}

\item 把例題改成「剛好經過一台飲水機」的路線數。
\begin{enumerate}
\item 每一格要記幾個數？ 寫出它們的意義和走進 \texttt{*} 時的轉移。
\item 第~\ref{sec:dp2-ex-fountain}~小節 $2 \times 3$ 的例子（\texttt{.**}／
\texttt{...}）答案是多少？
\end{enumerate}
\end{drills}

\section{實作練習}
\label{sec:dp2-practice}

下面是 Aki-Judge 上和本章相關的題目，依小節分組。 難度分四階：暖身是直接套用
本章的寫法；基本要自己定出狀態和轉移；實戰的題敘比較長，要先把題目轉換成本章
的形狀；挑戰會和其他章的技巧合在一起，或需要更多維的狀態。

\begin{center}
\begin{tabular}{llc}
\toprule
小節 & 題目 & 難度 \\
\midrule
路線計數與障礙 & \akiref{night-market-snack-plans}{夜市小吃組合數} & 暖身 \\
 & \akiref{routes-through-one-way-corridors}{單向走廊的路線數} & 暖身 \\
 & \akiref{seepage-from-the-top-row}{從頂層滲到底層的水路} & 基本 \\
 & \akiref{stepping-over-puddles}{跨過水坑的上學路線} & 基本 \\
 & \akiref{clear-just-one-shelf}{只移走一座貨架} & 基本 \\
\midrule
最小成本與最大分數 & \akiref{best-route-through-the-stamp-booth}{必經集章點的最高分路線} & 暖身 \\
 & \akiref{tolls-to-every-drop-off}{每個送貨點的最低過路費} & 基本 \\
 & \akiref{drilling-down-the-ore-pyramid}{礦層金字塔的鑽探路線} & 基本 \\
 & \akiref{best-score-at-every-exit-slot}{彈珠斜板各出口的最高分} & 基本 \\
 & \akiref{the-route-with-the-strongest-weakest-signal}{訊號最弱處最強的路線} & 基本 \\
 & \akiref{best-route-when-one-square-is-closed}{單格封閉後的最佳路線} & 挑戰 \\
 & \akiref{capped-battery-least-starting-charge}{有容量上限的電池：最少起始電量} & 挑戰 \\
\midrule
一格記多個數字 & \akiref{cheapest-routes-and-one-coin-more}{最省和只多一枚的路線} & 基本 \\
 & \akiref{one-take-drone-footage}{只錄一段的空拍路線} & 基本 \\
 & \akiref{the-cart-that-hates-turning}{轉彎次數有限的推車} & 實戰 \\
 & \akiref{waivers-from-the-supply-stations}{補給站的免罰券} & 實戰 \\
 & \akiref{a-workout-plan-without-back-to-back-muscle-groups}{不連練同一部位的課表} & 挑戰 \\
 & \akiref{toggle-tiles-on-a-round-trip}{來回踩燈磚} & 挑戰 \\
 & \akiref{two-boats-dragging-a-net}{雙船拖網} & 挑戰 \\
\midrule
正方形與矩形 & \akiref{best-crop-that-keeps-the-subject}{保留主角的最佳裁切} & 基本 \\
 & \akiref{booth-area-with-a-minimum-size}{尺寸有下限的最佳攤位區} & 基本 \\
 & \akiref{smallest-rectangle-reaching-the-target}{發電量達標的最小屋頂區塊} & 基本 \\
 & \akiref{the-most-fertile-stone-free-plot}{避開石頭的最肥沃菜園} & 實戰 \\
 & \akiref{sticker-sizes-on-a-cracked-wall}{裂磚牆上貼得下的壁貼尺寸} & 挑戰 \\
\bottomrule
\end{tabular}
\end{center}

本章的三題例題也在 Aki-Judge 上，讀完各節之後可以先自己寫一次再對照。

\section{讀完你應該會}
\label{sec:dp2-checklist}

\begin{itemize}
\item 看到只能往右、往下走的格子題，能用「最後一步從哪裡來」寫出計數、最少、最
多三種轉移，並正確處理第一列、第一行和走不進去的格子。
\item 能在「多開一圈」和「先判斷再用」兩種邊界寫法中選一種，寫完後用列數、行數
不同的輸入測試。
\item 能把二維表格改成只留一列的滾動陣列，說出填第 $j$ 格時陣列裡每一格是哪一
列的值，並判斷覆寫的方向。
\item 轉移時發現只知道「走到這裡的答案」不夠，能把缺的資訊（例如經過了沒有）加
進狀態，並估計表格的大小。
\item 能用「以這一格為右下角的最大邊長」解正方形的題目，並說明為什麼要取三個鄰
居的最小值；看到矩形的題目，能想到二維前綴和、固定兩列壓成一維、逐列的高度三
種工具。
\end{itemize}

\textbf{延伸閱讀}（書末會整理完整的參考資料）

\begin{itemize}
\item Antti Laaksonen，\textit{Competitive Programmer's Handbook}，7.3 節：格
子上的路徑，一頁就講完，可以拿來快速複習。
\item 劉汝佳，《算法競賽入門經典》第 2 版，9.1 節：從數字三角形講格子型的動態
規劃，比較遞迴、記憶化與填表。
\item Steven Halim 等，\textit{Competitive Programming 4}，Book 1，3.5.2 節：
總和最大的矩形，從暴力一路改進到固定兩列的做法。
\item Thomas H. Cormen 等，\textit{Introduction to Algorithms}，第 4 版，
14.4 節：最長共同子序列，同樣是一格看上、左、左上三格的二維表格，《C++ 演算法解題 2》第~2~章會
再遇到。
\end{itemize}

\addcontentsline{toc}{section}{練習解答}
\section*{練習解答}

\begin{enumerate}
\item
\begin{enumerate}
\item \texttt{best[i][j]}：從左上角走到 $(i, j)$、包含兩端的最少過路費。 轉移
是 $\min(\text{上}, \text{左}) + a_{i,j}$，第一列只能從左、第一行只能從上；
初始值是起點自己的 $3$。 填好的表：
\[
\begin{matrix}
3 & 4 & 8 & 9 \\
5 & 11 & 9 & 17 \\
10 & 12 & 10 & 13 \\
11 & 17 & 12 & 13
\end{matrix}
\]
答案是 $13$。 從右下角往回找「是從哪一格來的」：$13 = 12 + 1$ 來自左邊，
$12 = 10 + 2$ 來自上面，依此類推，路線是
$3 \to 1 \to 4 \to 1 \to 1 \to 2 \to 1$，也就是右、右、下、下、下、右。
\item 不能走的格子是 \texttt{LLONG\_MAX}，它南邊或東邊的格子取最小值時，如果
另一邊也是 \texttt{LLONG\_MAX}（兩邊都到不了），就會算
\texttt{LLONG\_MAX + a[i][j]}，溢位成負數，變成「最省」。 改法：不能走的格子
直接跳過、保持 \texttt{LLONG\_MAX}；能走的格子，兩個來源都是
\texttt{LLONG\_MAX} 時保持 \texttt{LLONG\_MAX}，否則才拿比較小的那個加上
\texttt{a[i][j]}。 也可以用一個不會溢位的「很大的數」，例如 $4 \times
10^{18}$ 或 \texttt{LLONG\_MAX / 2}，但要確認「很大的數加上一條路線的總和」仍
然不會超過 \texttt{long long}，而且最後要判斷答案是不是還那麼大（到不了）。
\end{enumerate}

\item
\begin{enumerate}
\item \texttt{0 4 4 4}。
\item \texttt{2 2 4 4}。 前兩格是第 $1$ 列的值（\texttt{best[1][0]} $= 2$、
\texttt{best[1][1]} $= 2$），後兩格還是第 $0$ 列的值。 這正是
圖~\ref{fig:dp2-roll} 的狀態。
\end{enumerate}

\item 印出 \texttt{15}，正確是 $11$。 一整張表時，施工的格子一開始就是 $0$，
\texttt{continue} 不會出錯；只留一列時，\texttt{f[j]} 裡是上一列同一行的值，
\texttt{continue} 讓它原封不動留下來，等於施工的路口「到得了」，而且路線數和
它北邊的路口一樣。 範例裡 $(3, 4)$ 施工，\texttt{f[4]} 留著第 $2$ 列的 $2$，
最後算出 $15$。 改法是把 \texttt{continue} 換成 \texttt{f[j] = 0;}（再
\texttt{continue}）。 這也是滾動陣列的一般原則：每一格每一輪都要寫入，不能
因為「不能走」就跳過。

\item
\begin{enumerate}
\item 回傳 $35$，一共被呼叫 $69$ 次。
\item 回傳 $1$ 的呼叫是遞迴樹的葉子，答案是所有葉子的 $1$ 加起來，所以葉子有
「回傳值」個。 其他每一次呼叫剛好有兩個子呼叫，這樣的樹，內部的點數比葉子少
一，總數是 $2 \times$ 回傳值 $- 1$。 回傳值本身就是路線數，所以不加記憶化的
遞迴，呼叫次數和路線數同一個數量級，$500 \times 500$ 的格子完全不可能。
\item $25$ 次。 需要算的是 $i \ge 2$、$j \ge 2$ 的 $3 \times 4 = 12$ 個狀態，每
個第一次被算時呼叫兩次，加上最外面那一次，$1 + 2 \times 12 = 25$。
\end{enumerate}

\item \texttt{side} 是
\[
\begin{matrix}
1 & 1 & 1 & 1 & 1 \\
1 & 2 & 1 & 2 & 1 \\
1 & 1 & 1 & 2 & 2 \\
1 & 1 & 2 & 2 & 3 \\
1 & 1 & 1 & 2 & 3
\end{matrix}
\]
邊長 $1$ 到 $4$ 各有 $25, 9, 2, 0$ 個位置。 例如第 $4$ 列第 $4$ 行（從 $0$ 編
號）：四格都是 \texttt{b}，上面是 $3$，左邊和左上都是 $2$，取最小值加一是
$3$。

\item
\begin{enumerate}
\item $f(i, j) = f(i - 1, j) + f(i, j - 1) + f(i - 1, j - 1)$，三個不存在的鄰居
都當成 $0$。
\item 填第 $j$ 格時，需要上一列的 \texttt{f[j]}（上）和 \texttt{f[j - 1]}
（左上），以及這一列的 \texttt{f[j - 1]}（左）。 可是 \texttt{f[j - 1]} 只有一
格，填第 $j - 1$ 格時已經把上一列的值蓋掉了。 所以在蓋掉之前，先把舊的值存在
變數 \texttt{diag} 裡：
\begin{code}{cpp}{練習 6 的解答}
for (int i = 1; i <= r; i++) {
    long long diag = 0;                   // 上一列的 f[j - 1]
    for (int j = 1; j <= c; j++) {
        long long up = f[j];     // 上一列的 f[j]，等一下會被蓋掉
        if (g[i - 1][j - 1] == 'x') {
            f[j] = 0;
        } else if (i == 1 && j == 1) {
            f[j] = 1;
        } else {
            f[j] = up + f[j - 1] + diag;
        }
        diag = up;               // 下一格的左上，就是這一格的上
    }
}
\end{code}
\end{enumerate}

\item
\begin{enumerate}
\item 三個數：走到這裡經過了 $0$ 台、剛好 $1$ 台、$2$ 台以上的路線數。 經過 $2$
台以上的不會再變回合格，但仍然要記下來，否則下一台飲水機會把它們誤算成「剛好
$1$ 台」。（只記 $0$ 台和 $1$ 台、遇到 \texttt{*} 時把 $1$ 台的直接丟掉，也
可以，效果一樣。） 走進 \texttt{*} 時，$0$ 台的變成 $1$ 台，$1$ 台和 $2$ 台
以上的都變成 $2$ 台以上；走進 \texttt{.} 時三組都不變。 答案是右下角的「剛好
$1$ 台」。
\item $1$。 只有右下右這條剛好經過一台。
\end{enumerate}
\end{enumerate}

% 佔位：本章尚未開始。book.tex 裡先註解掉。
\chapter{圖的表示與走訪}

本章尚未開始。

% 佔位：本章尚未開始。book.tex 裡先註解掉。
\chapter{格子圖與狀態 BFS}

本章尚未開始。

\chapter{樹}
\label{ch:trees}

% 本章的圖共用的樣式（點、強調的點、第二焦點、邊）
\tikzset{
    tv/.style={circle, draw=black!60, thick, minimum size=0.72cm, inner sep=0pt, font=\small},
    tvhl/.style={tv, draw=memStroke, fill=memFill},
    tvalt/.style={tv, draw=optStroke, fill=optFill},
    te/.style={thick, draw=black!60},
    tehl/.style={very thick, draw=memStroke},
    tlbl/.style={font=\scriptsize, text=black!60},
}

過年回老家，長輩拿出一本族譜。 最上面是開基祖，每個人底下連著自己的子女，一
代接一代往下畫。 你想知道三件事：自己是第幾代、整本族譜一共傳了幾代、某位叔公
這一房（他和他所有的子孫）一共有幾個人。 族譜如果只寫了「每個人的父親是誰」，
你會怎麼算？ 往下讀之前，先想一想第一件事和第三件事，哪一件要從上往下算，哪
一件要從下往上算。

自己是第幾代，只要知道父親是第幾代再加一，所以從開基祖往下一代一代算。 一房
有幾個人，要先知道每個子女那一房有幾個人，再加上自己，所以要從最下面往上算。
這兩個方向是這一章的骨架：族譜、公司的組織、電腦的資料夾、城市的自來水管網，
畫出來都是同一種形狀，叫\textbf{樹}（tree）。

這一章先講樹的名詞和兩種存法，再講輸入只給「誰和誰相連」、沒有說誰是上層時，
怎麼自己定出上下關係。 接著是由上往下算深度、由下往上算高度和子樹的大小，並
且整理出一個不用遞迴的由下往上順序。 後半講兩個在樹上很常用的結論：一條邊被
多少對點的路徑經過，以及樹上最遠的兩點（樹直徑）怎麼用兩次走訪找到。 最後簡
單看一下二元樹的三種走訪順序，並預告第~20~章的樹形動態規劃。

這一章會一直用到前面三章的東西。 第~6~章的例題「層層彙整的報告」已經用佇列
由下往上算過一棵樹；第~8~章量過遞迴可以多深；第~11~章講過 BFS、DFS，以及遞迴
太深時怎麼改成迴圈。 這些不再重講，用到時會指回去。 語法上只需要
\texttt{vector<vector<int>>} 和 \texttt{pair}（銜接書第~8~章）。 沒讀過銜接書的讀
者，只要會用 \texttt{vector} 存鄰接串列、會寫 BFS，就能讀這一章。

\section{有根樹的名詞與存法}
\label{sec:tree-terms}

第~11~章的圖可以有環，兩點之間也可能有好幾條路。 \textbf{樹}是沒有環、而且連
通的無向圖。 $n$ 個點的樹有下面幾個性質，它們其實互相等價，知道其中兩三個就
能推出其他的：

\begin{itemize}
\item 剛好有 $n - 1$ 條邊。 一個連通的圖至少要 $n - 1$ 條邊；剛好 $n - 1$ 條時
不會有環，多一條就一定有環。
\item 任兩點之間剛好只有一條路徑（不重複經過點）。 有兩條的話，兩條路合起來就
繞成一個環。
\item 拿掉任何一條邊，樹就斷成兩塊。 這一點第~\ref{sec:tree-edge}~節會用到。
\end{itemize}

題目常常寫「$n$ 個點、$n - 1$ 條邊，任兩點都能到達」，或「任兩點之間的路線只
有一條」，這兩種說法都是在說輸入是一棵樹。

\subsection{有根樹}

族譜、組織圖都有一個明確的頂端。 在樹上指定一個點當作\textbf{根}（root），
其他點都掛在它下面，就是\textbf{有根樹}（rooted tree）。 下圖是章首的族譜，
九個人編號 $1$ 到 $9$，開基祖是 $4$ 號：

\begin{center}
\begin{tikzpicture}[x=1.15cm, y=1.05cm]
\node[tvhl] (4) at (3,3) {4};
\node[tv] (2) at (1.4,2) {2};
\node[tv] (7) at (3,2) {7};
\node[tv] (9) at (4.6,2) {9};
\node[tv] (1) at (0.6,1) {1};
\node[tv] (5) at (2.2,1) {5};
\node[tv] (3) at (3,1) {3};
\node[tv] (6) at (1.6,0) {6};
\node[tv] (8) at (2.8,0) {8};
\foreach \a/\b in {4/2, 4/7, 4/9, 2/1, 2/5, 7/3, 5/6, 5/8} {\draw[te] (\a) -- (\b);}
\foreach \y/\d in {3/0, 2/1, 1/2, 0/3} {\node[tlbl, anchor=east] at (-0.6,\y) {深度 \d};}
\end{tikzpicture}
\diagcaption{章首的族譜。 根是 $4$；左邊標出每一層的深度。}
\label{fig:tree-family}
\end{center}

用這張圖說明名詞：

\begin{itemize}
\item 一條邊連著的兩個點，上面那個是下面那個的\textbf{父節點}（parent），下面
的是上面的\textbf{子節點}（child）。 $2$ 的父節點是 $4$，子節點是 $1$、$5$。
根沒有父節點，其他每個點都剛好有一個。
\item 沒有子節點的點叫\textbf{葉}（leaf），圖上是 $1$、$3$、$6$、$8$、$9$。
同一個父節點底下的點互為\textbf{兄弟}（sibling）。
\item 從某個點一直往上走到根，路上經過的點都是它的\textbf{祖先}（ancestor）；
反過來，它是那些點的\textbf{後代}（descendant）。 $6$ 的祖先是 $5$、$2$、$4$。
\item 一個點和它所有的後代，合起來叫做以它為根的\textbf{子樹}（subtree）。 $2$
的子樹是 $\{2, 1, 5, 6, 8\}$，共 $5$ 個點。 章首問的「叔公這一房」，就是叔公的子
樹。
\item 根到一個點的路徑有幾條邊，叫這個點的\textbf{深度}（depth）。 根的深度是
$0$，$6$ 的深度是 $3$。
\item 從一個點往下走到最遠的葉要幾條邊，叫這個點的\textbf{高度}（height）。 葉
的高度是 $0$，$2$ 的高度是 $2$（$2 \to 5 \to 6$）。 根的高度就是整棵樹的高度，
這裡是 $3$。
\end{itemize}

深度和高度有些題目從 $1$ 開始數（根是「第 $1$ 代」），有些數的是點而不是邊。
本書一律從 $0$ 開始數邊，讀題時先確認題目的定義，差一就是答案錯。

把每條邊想成從父節點指向子節點，有根樹就是第~9~章的有向無環圖，由上往下或由
下往上的計算都是照先後順序算。

\subsection{兩種存法}

有根樹最常見的兩種存法：

\begin{itemize}
\item \textbf{父節點陣列}：\texttt{parent[v]} 是 $v$ 的父節點，根填 $0$。 一個
陣列就存下整棵樹，很多題目的輸入直接就是這個陣列。 從一個點往上走到根很容易，
一直做 \texttt{v = parent[v]} 就好；但想知道一個點的子節點有哪些，得把整個陣列
翻一遍。
\item \textbf{子節點串列}：\texttt{children[v]} 是一個 \texttt{vector}，放 $v$
的所有子節點。 寫法和第~11~章的鄰接串列一樣，只是每條邊只存「父指向子」一個方
向。 由上往下走訪要用它。
\end{itemize}

圖~\ref{fig:tree-family} 的父節點陣列是

\begin{center}
\begin{tabular}{c|ccccccccc}
\toprule
$v$ & 1 & 2 & 3 & 4 & 5 & 6 & 7 & 8 & 9 \\
\midrule
\texttt{parent[v]} & 2 & 4 & 7 & 0 & 2 & 5 & 4 & 5 & 4 \\
\bottomrule
\end{tabular}
\end{center}

從父節點陣列建子節點串列，順便找出根：

\begin{code}{cpp}{由父節點陣列建子節點串列}
int root = 0;
for (int v = 1; v <= n; v++) {
    if (parent[v] == 0) {
        root = v;                       // 沒有父節點的就是根
    } else {
        children[parent[v]].push_back(v);
    }
}
\end{code}

依 $v = 1, 2, \ldots$ 的順序加入，每個串列裡的子節點是由小到大排好的，例如
\texttt{children[4]} 是 $2, 7, 9$。

注意這張表的根是 $4$，不是 $1$；$1$ 號的父節點 $2$ 反而比它大。 題目沒有明說
的話，根不一定是 $1$，父節點的編號也不一定比子節點小，第~6~章「層層彙整的報
告」的常見錯誤就是吃了這個虧。

有些題目反過來，給每個點的子節點名單（先給個數，再列出編號），這時可以直接存成
\texttt{children}；要找根，就找哪一個點沒有出現在任何人的名單裡。 做法是開一
個陣列記錄「當過子節點」，最後沒被記錄的那一個就是根。

\begin{apcsnote}
樹的題目在找根和編號上常有陷阱：根不一定是 $1$、父節點的編號不一定比較小、
$n$ 可能等於 $1$（沒有任何邊）。 寫完之後，至少用三組自己造的輸入測試：只有一
個點、一條很長的鏈、根的編號在最後面。
\end{apcsnote}

\section{給無向邊時怎麼定根}
\label{sec:tree-rooting}

更常見的輸入是 $n - 1$ 行「$u$ $v$」，只說這兩點之間有一條邊，沒說誰是父節點。
這時先照第~11~章存成無向的鄰接串列，再自己挑一個根（題目指定的，或任選一個，
例如 $1$），從根出發走一遍，走的時候把每條邊定成「先走到的那一端是父節點」。

在樹上走訪，不需要第~11~章的 \texttt{seen} 陣列。 從父節點走到 $u$ 之後，$u$
的鄰居只有兩種：它的父節點，和它的子節點。 如果還有別的鄰居先前已經走過，就表
示有兩條路通到那個點，樹上不可能。 所以只要「不走回父節點」，就不會重複走到同
一個點：

\begin{code}{cpp}{從根出發做 BFS，定出父節點和走訪順序}
vector<int> parent(n + 1, 0);
vector<int> order = {root};             // 同時是佇列和走訪順序
for (int i = 0; i < (int)order.size(); i++) {
    int u = order[i];
    for (int v : adj[u]) {
        if (v != parent[u]) {           // 不走回父節點
            parent[v] = u;
            order.push_back(v);
        }
    }
}
\end{code}

這段程式用一個 \texttt{vector} 加上索引 \texttt{i} 當佇列，代替
\texttt{queue}：\texttt{i} 指著目前要處理的點，新的點接在尾端。 每個點只會被放進去一
次，迴圈結束時 \texttt{order} 剛好裝著 $n$ 個點，順序就是 BFS 拜訪的順序。 東
西不從前面拿掉，是因為下一節要再用一次這個順序。

根的 \texttt{parent} 是 $0$，而點從 $1$ 編號，不會有鄰居剛好是 $0$，所以根的
每個鄰居都會被當成子節點。 如果點從 $0$ 編號，根的父節點要改成 $-1$。

走訪的同時就能由上往下算深度：在 \texttt{parent[v] = u;} 後面加一行
\texttt{depth[v] = depth[u] + 1;}。 處理 $u$ 時它的深度已經算好，因為 $u$ 是從
它的父節點那裡放進來的，父節點比它早處理。 同樣的寫法可以算「從根沿路累加」的
任何量，例如根到每個點的路線總長。

\subsection{遞迴的寫法}

用遞迴寫，就是多傳一個參數記住「從哪裡來」：

\begin{code}{cpp}{遞迴走訪，不走回來的地方}
void dfs(int u, int from) {
    parent[u] = from;
    for (int v : adj[u]) {
        if (v != from) {
            depth[v] = depth[u] + 1;
            dfs(v, u);
        }
    }
}
// 呼叫：dfs(root, 0);
\end{code}

它比 BFS 短，但深度和樹的高度一樣。 一條 $2 \times 10^5$ 個點的鏈，從一端出發
就要疊 $2 \times 10^5$ 層，照第~8~章量過的上限會當掉（第~\ref{sec:ex-roads}~小
節有實測）。 這一章的例題 $n$ 都到 $2 \times 10^5$，一律用上面的 BFS 順序；需
要遞迴的「離開」時機時，用第~11~章的改法二。

\subsection{例題：聯絡網找發起人}
\label{sec:ex-phonetree}

\textbf{題意}（\akiref{who-started-the-phone-tree}{聯絡網找發起人}）

\begin{itemize}
\item 班上有 $n$ 位同學，緊急聯絡網是一棵有根樹：發起人先打電話給自己負責的
「下線」，每個人接到電話後再打給自己的下線。 除了發起人，每個人剛好有一位
「上線」。
\item 原始名單弄丟了，只剩 $n - 1$ 筆通話紀錄「$u$ $v$」（兩人之間直接通過話，
不知道誰打給誰），以及每位同學回報的下線人數 $c_i$。
\item 保證剛好有一種發起人的選法，讓每個人的下線數都等於 $c_i$。 輸出發起人，
以及每位同學的上線（發起人輸出 $0$）。
\item $n \le 2 \times 10^5$。
\end{itemize}

例如下圖左邊是 $8$ 位同學的通話紀錄，每個點旁邊標著回報的 $c_i$。 發起人是
$6$，定好方向後是右邊的樣子：$6$ 只打給 $3$，$3$ 再打給 $1$、$5$、$8$。

\begin{center}
\begin{tikzpicture}[x=1.1cm, y=1.0cm]
\begin{scope}
\node[font=\small] at (2,3.2) {通話紀錄（沒有方向）};
\node[tv] (3) at (2,1.5) {3};
\node[tv] (1) at (0.8,2.4) {1};
\node[tv] (6) at (3.2,2.4) {6};
\node[tv] (5) at (1,0.7) {5};
\node[tv] (8) at (3,0.7) {8};
\node[tv] (2) at (0.2,-0.3) {2};
\node[tv] (7) at (1.6,-0.3) {7};
\node[tv] (4) at (3.4,-0.3) {4};
\foreach \a/\b in {3/1, 5/3, 6/3, 2/5, 5/7, 8/3, 4/8} {\draw[te] (\a) -- (\b);}
\foreach \v/\c in {1/0, 5/2} {\node[tlbl, left] at (\v.west) {$c = \c$};}
\foreach \v/\c in {3/3, 6/1, 8/1} {\node[tlbl, right] at (\v.east) {$c = \c$};}
\foreach \v/\c in {2/0, 7/0, 4/0} {\node[tlbl, below] at (\v.south) {$c = \c$};}
\end{scope}
\begin{scope}[xshift=6cm]
\node[font=\small] at (2,3.2) {以 6 為根};
\node[tvhl] (6) at (2,2.4) {6};
\node[tv] (3) at (2,1.5) {3};
\node[tv] (1) at (0.8,0.6) {1};
\node[tv] (5) at (2,0.6) {5};
\node[tv] (8) at (3.2,0.6) {8};
\node[tv] (2) at (1.5,-0.3) {2};
\node[tv] (7) at (2.5,-0.3) {7};
\node[tv] (4) at (3.2,-0.3) {4};
\foreach \a/\b in {6/3, 3/1, 3/5, 3/8, 5/2, 5/7, 8/4} {\draw[te] (\a) -- (\b);}
\end{scope}
\end{tikzpicture}
\diagcaption{左：通話紀錄和每人回報的下線數。 右：以發起人 $6$ 為根定出上下關係。}
\label{fig:tree-phone}
\end{center}

\textbf{想法}

如果知道發起人，就是上一段的 BFS：從發起人出發，不走回上線，每個人的上線就是
\texttt{parent}。 問題只剩下：發起人是誰？

最直接的做法是每個人都試一次：假設他是發起人，做一次 BFS 定出方向，再檢查每個
人的下線數是不是 $c_i$。 一次 $O(n)$，試 $n$ 個人是 $O(n^2)$，$n = 2 \times
10^5$ 時是 $4 \times 10^{10}$，來不及。 不過它很適合拿來對拍。

一個看起來合理的猜法是「下線最多的人是發起人」，畢竟他在最上面。 圖
~\ref{fig:tree-phone} 就是反例：$3$ 要打三通電話，是最忙的人，但發起人是只打
一通的 $6$。 下線多少和在不在最上面沒有關係。

換個方向想：每個人在通話紀錄裡出現幾次（度數），和下線數有什麼關係？ 一個不是
發起人的人，通過電話的對象是他的上線加上所有下線，所以度數 $= c_i + 1$。 發起
人沒有上線，度數 $= c_i$。 所以只要找度數剛好等於 $c_i$ 的那個人。 題目保證答
案存在，真正的發起人一定滿足這個條件，其他人的度數都比 $c_i$ 多一，不會被誤認。

圖~\ref{fig:tree-phone} 的度數和 $c_i$：

\begin{center}
\begin{tabular}{c|cccccccc}
\toprule
同學 & 1 & 2 & 3 & 4 & 5 & 6 & 7 & 8 \\
\midrule
度數 & 1 & 1 & 4 & 1 & 3 & 1 & 1 & 2 \\
$c_i$ & 0 & 0 & 3 & 0 & 2 & 1 & 0 & 1 \\
\bottomrule
\end{tabular}
\end{center}

只有 $6$ 的兩個數相等。

同樣的模型也出現在別的地方：網路設備只留下「哪兩台之間有線」的紀錄，要從某台
主機回推每台設備的上游；水管網只記錄管子接在哪兩個閥門之間，要從水塔定出每段
管子的水流方向。

\textbf{程式}

\begin{example}[label=ex:phonetree]{cpp}{(phonetree.cpp) 聯絡網找發起人}
#include <iostream>
#include <vector>
using namespace std;

int main() {
    ios::sync_with_stdio(false);
    cin.tie(nullptr);
    int n;
    cin >> n;
    vector<int> c(n + 1);
    for (int v = 1; v <= n; v++) {
        cin >> c[v];
    }
    vector<vector<int>> adj(n + 1);
    for (int i = 0; i < n - 1; i++) {
        int u, v;
        cin >> u >> v;
        adj[u].push_back(v);
        adj[v].push_back(u);
    }

    int root = 0;
    for (int v = 1; v <= n; v++) {
        if ((int)adj[v].size() == c[v]) {   // 沒有上線
            root = v;
        }
    }

    // 從根做 BFS，把每條邊定成「上線 → 下線」
    vector<int> parent(n + 1, 0);
    vector<int> order = {root};
    for (int i = 0; i < (int)order.size(); i++) {
        int u = order[i];
        for (int v : adj[u]) {
            if (v != parent[u]) {           // 不走回上線
                parent[v] = u;
                order.push_back(v);
            }
        }
    }

    cout << root << '\n';
    for (int v = 1; v <= n; v++) {
        cout << parent[v] << (v == n ? '\n' : ' ');
    }
    return 0;
}
\end{example}

\iosample{
8 \\
0 0 3 0 2 1 0 1 \\
3 1 \\
5 3 \\
6 3 \\
2 5 \\
5 7 \\
8 3 \\
4 8 \\
}{
6 \\
3 5 6 8 3 0 5 3 \\
}

BFS 從 $6$ 出發，每一輪處理 \texttt{order[i]}：

\begin{center}
\begin{tabular}{cclll}
\toprule
\texttt{i} & 處理 & 鄰居 & 新的子節點 & \texttt{order} \\
\midrule
 & & & & 6 \\
0 & 6 & 3 & 3 & 6 3 \\
1 & 3 & 1、5、6、8 & 1、5、8（6 是上線，跳過）& 6 3 1 5 8 \\
2 & 1 & 3 & 無 & 6 3 1 5 8 \\
3 & 5 & 3、2、7 & 2、7 & 6 3 1 5 8 2 7 \\
4 & 8 & 3、4 & 4 & 6 3 1 5 8 2 7 4 \\
5～7 & 2、7、4 & 只有上線 & 無 & \\
\bottomrule
\end{tabular}
\end{center}

\begin{itemize}
\item \texttt{adj[v].size()} 就是度數。 它的型別是 \texttt{size\_t}（不帶正負
號），和 \texttt{int} 直接比較時 \texttt{-Wall -Wextra} 會警告，所以先轉成
\texttt{int}。
\item $n = 1$ 時沒有任何通話紀錄，$1$ 號的度數 $0$ 等於 $c_1 = 0$，發起人是
$1$，BFS 只處理他自己，輸出 \texttt{1} 和 \texttt{0}，不用特別處理。
\end{itemize}

\textbf{複雜度}

算度數、找發起人 $O(n)$，BFS 每個點處理一次、每條邊從兩端各看一次，也是
$O(n)$。 最大的測試資料在本書測試的環境下，CPU 時間不到 $0.1$ 秒。

\subsection{常見錯誤}

\begin{itemize}
\item \textbf{下線最多的當發起人}：把找根的迴圈改成找 $c_i$ 最大的人。 範例的
輸入會印出 \texttt{3} 和 \texttt{3 5 0 8 3 3 5 3}。 如果範例剛好是發起人打最
多電話，這個錯誤就看不出來。
\item \textbf{用 \texttt{parent[v] == 0} 判斷還沒走過}：把
\texttt{if (v != parent[u])} 換成 \texttt{if (parent[v] == 0)}，想用「還沒有父
節點」代替 \texttt{seen}。 問題在根：它的 \texttt{parent} 本來就是 $0$，處理
$3$ 的時候，鄰居 $6$ 看起來「還沒走過」，於是被設成 $3$ 的子節點。 範例的輸入
第二行會印出 \texttt{3 5 6 8 3 3 5 3}，發起人的上線變成 $3$。 要用這種寫法，
\texttt{parent} 的初始值得用 $-1$，根設成 $0$。
\item \textbf{直接照輸入的方向當上下線}：以為每筆紀錄的第一個人是上線。 範例裡
\texttt{3 1} 剛好是對的，\texttt{5 3} 就反了。 紀錄本來就沒有方向，一定要從根
走一遍。
\end{itemize}

\section{由下往上：高度、子樹大小與子樹和}
\label{sec:tree-bottom-up}

深度由上往下算，父節點先算好，子節點加一。 下面三個量剛好相反，要先知道所有子
節點的值，才算得出自己的：

\begin{itemize}
\item \textbf{高度}：葉是 $0$；其他點是「子節點高度的最大值加一」。
\item \textbf{子樹大小}：$v$ 的子樹有幾個點。 $\mathit{sz}(v) = 1 +$ 所有子節
點的 $\mathit{sz}$ 加起來。
\item \textbf{子樹和}：每個點有一個值 $a_v$，$v$ 的子樹裡所有點的值加起來。
$\mathit{sum}(v) = a_v +$ 所有子節點的 $\mathit{sum}$ 加起來。 子樹大小就是每
個 $a_v$ 都是 $1$ 的子樹和。
\end{itemize}

第~8~章章首算資料夾的總大小，就是子樹和：每個資料夾的大小是自己的檔案加上每
個子資料夾的總大小。 那裡用遞迴，函式回傳的那一刻，子資料夾的答案已經算完。
在樹上，遞迴版本是在 \texttt{dfs(v, u)} 回來之後加一行
\texttt{sz[u] += sz[v];}。 問題又是深度：鏈狀的樹會讓遞迴疊到 $n$ 層。

\subsection{把 BFS 順序倒過來}

不用遞迴也能得到「子節點一定先算完」的順序。 看上一節的 \texttt{order}：一個點
是在它的父節點被處理時才放進去的，所以\textbf{在 \texttt{order} 裡，每個點都排
在它所有子節點的前面}。 倒過來從最後一個看到第一個，每個點就都排在它所有子節
點的後面。 輪到 $v$ 的時候，它的子節點都已經看過了，而每個子節點在被看到時已經
把自己的值加給 $v$，所以 $v$ 的值已經是最終的值，可以再把它加給自己的父節點：

\begin{code}{cpp}{倒著看 BFS 順序，由下往上}
for (int i = n - 1; i >= 1; i--) {      // 跳過根 order[0]
    int v = order[i], p = parent[v];
    height[p] = max(height[p], height[v] + 1);
    sz[p] += sz[v];
}
\end{code}

每個點只把自己的值交給父節點一次，不需要子節點串列，也不需要遞迴。

第~6~章的「層層彙整的報告」是另一種得到這個順序的方法：記下每個人還有幾個部屬
沒交，一開始把沒有部屬的人放進佇列，誰的部屬全交齊了才輪到誰。 兩種方法要保證
的是同一件事：處理一個點之前，它的子節點都處理完了。 比較一下：

\begin{center}
\begin{tabular}{p{0.24\linewidth}p{0.33\linewidth}p{0.33\linewidth}}
\toprule
 & 第~6~章：數還差幾個 & 本章：倒著看 BFS 順序 \\
\midrule
需要的資料 & 只要父節點陣列 & 子節點串列或鄰接串列 \\
順帶得到 & 只有由下往上的順序 & 由上往下的順序（深度）也有了 \\
推廣 & 第~18~章的拓撲排序，不限於樹 & 只適用於樹 \\
\bottomrule
\end{tabular}
\end{center}

輸入是無向邊時，本來就要做一次 BFS 定根，倒著看同一個 \texttt{order} 最省事。
輸入只有父節點陣列、又只需要由下往上的量時，第~6~章的做法連子節點串列都不用
建。

\subsection{把前面的東西放在一起}

下面的程式讀入族譜的父節點陣列和每個人的值，算出每個人的深度、高度、子樹大小
和子樹和。 值 $a_v$ 可以想成每個人這一戶要訂幾個便當，子樹和就是「這一房一共要
訂幾個」。

\begin{example}[label=ex:tree-basic]{cpp}{(tree.cpp) 深度、高度、子樹大小與子樹和}
#include <algorithm>
#include <iostream>
#include <vector>
using namespace std;

int main() {
    int n;
    cin >> n;
    vector<int> parent(n + 1);
    vector<vector<int>> children(n + 1);
    int root = 0;
    for (int v = 1; v <= n; v++) {
        cin >> parent[v];
        if (parent[v] == 0) {
            root = v;                       // 沒有父節點的就是根
        } else {
            children[parent[v]].push_back(v);
        }
    }
    vector<long long> sum(n + 1);
    for (int v = 1; v <= n; v++) {
        cin >> sum[v];                      // 先放自己的值
    }

    // 由上往下：order 是佇列也是走訪順序
    vector<int> order = {root};
    vector<int> depth(n + 1, 0);
    for (int i = 0; i < (int)order.size(); i++) {
        int u = order[i];
        for (int c : children[u]) {
            depth[c] = depth[u] + 1;        // 父節點已算好
            order.push_back(c);
        }
    }

    // 由下往上：倒著看 order
    vector<int> height(n + 1, 0), sz(n + 1, 1);
    for (int i = n - 1; i >= 1; i--) {      // 跳過根 order[0]
        int v = order[i], p = parent[v];
        height[p] = max(height[p], height[v] + 1);
        sz[p] += sz[v];
        sum[p] += sum[v];
    }

    cout << "root " << root << "\norder";
    for (int v : order) {
        cout << ' ' << v;
    }
    cout << '\n';
    for (int v = 1; v <= n; v++) {
        cout << v << ": " << depth[v] << ' ' << height[v] << ' '
             << sz[v] << ' ' << sum[v] << '\n';
    }
    return 0;
}
\end{example}

\iosample{
9 \\
2 4 7 0 2 5 4 5 4 \\
3 2 4 6 1 2 5 3 2 \\
}{
root 4 \\
order 4 2 7 9 1 5 3 6 8 \\
1: 2 0 1 3 \\
2: 1 2 5 11 \\
3: 2 0 1 4 \\
4: 0 3 9 28 \\
5: 2 1 3 6 \\
6: 3 0 1 2 \\
7: 1 1 2 9 \\
8: 3 0 1 3 \\
9: 1 0 1 2 \\
}

輸出的每一行是「點：深度 高度 子樹大小 子樹和」。 由下往上那個迴圈，每一步之
後父節點的三個值：

\begin{center}
\begin{tabular}{cccccc}
\toprule
\texttt{i} & $v$ & $p$ & \texttt{height[p]} & \texttt{sz[p]} & \texttt{sum[p]} \\
\midrule
8 & 8 & 5 & 1 & 2 & 4 \\
7 & 6 & 5 & 1 & 3 & 6 \\
6 & 3 & 7 & 1 & 2 & 9 \\
5 & 5 & 2 & 2 & 4 & 8 \\
4 & 1 & 2 & 2 & 5 & 11 \\
3 & 9 & 4 & 1 & 2 & 8 \\
2 & 7 & 4 & 2 & 4 & 17 \\
1 & 2 & 4 & 3 & 9 & 28 \\
\bottomrule
\end{tabular}
\end{center}

$5$ 在 \texttt{i} $= 5$ 時把值交給 $2$，那時它的兩個子節點 $8$、$6$（\texttt{i}
$= 8, 7$）已經交過了。 根 $4$ 最後才收齊三個子節點，\texttt{sz[4]} $= 9$ 正好
是全部的人數，可以拿來檢查。

\begin{itemize}
\item \texttt{sum} 用 \texttt{long long}。 題目的值常到 $10^9$，$n$ 個加起來
就超過 \texttt{int}。
\item \texttt{sz} 初始值是 $1$（自己），\texttt{height} 是 $0$（先當作葉），由下往上只會把它們變大。
\item \texttt{order} 的長度不是 $n$，表示輸入不是一棵樹（例如有環、不連通）。 題目
保證是樹時不用檢查。
\end{itemize}

\subsection{什麼時候用、什麼時候不能用}

\begin{itemize}
\item 「每個點的答案只由子節點的答案決定」：子樹大小、子樹和、子樹裡的最大值、
到最遠的葉有多遠，都倒著看 BFS 順序。 「每個點的答案只由父節點決定」：深度、
從根累加的距離，順著看。
\item 一個點的答案同時要用到上面和下面的資訊時，例如「離每個點最遠的點有多遠」，
一次由下往上不夠，要再配合一次由上往下，或用第~\ref{sec:tree-diameter}~節的直
徑性質。
\item 不是樹的圖（有環，或一個點有兩個「上層」）沒有這種順序可以倒過來；第~18~
章會處理有向無環圖的情況。
\end{itemize}

\section{每條邊的貢獻}
\label{sec:tree-edge}

樹的每條邊都有一個特別的性質：拿掉它，樹剛好斷成兩塊。 以任一點為根時，這兩塊
很好描述：邊的下端是子節點 $c$，一塊是 $c$ 的子樹，有 $\mathit{sz}(c)$ 個點；另
一塊是其他所有點，有 $n - \mathit{sz}(c)$ 個。 所以算出子樹大小，每條邊兩側各
有幾個點就都知道了。

\subsection{例題：每條道路的通訊量}
\label{sec:ex-roads}

\textbf{題意}（\akiref{messages-through-every-road}{每條道路的通訊量}）

\begin{itemize}
\item 山區有 $n$ 個村莊，由 $n - 1$ 條雙向道路連成一棵樹。
\item 郵局要在每一對村莊之間各寄一封信，信沿著兩村之間唯一的路線送。
\item 依輸入順序，對每條道路求出有多少對村莊的路線會經過它。
\item $2 \le n \le 2 \times 10^5$。
\end{itemize}

例如下圖的 $7$ 個村莊。 道路 $4$--$2$ 拿掉之後，一邊是 $\{4, 5, 6\}$，另一邊是
$\{1, 2, 3, 7\}$。 一對村莊的路線會經過這條路，剛好是兩個村莊分在兩邊的時候，
例如 $5$ 和 $7$ 的路線是 $5 \to 4 \to 2 \to 1 \to 7$。 所以有 $3 \times 4 =
12$ 對。

\begin{center}
\begin{tikzpicture}[x=1.15cm, y=1.0cm]
\node[tv] (1) at (2,3) {1};
\node[tv] (2) at (1,2) {2};
\node[tv] (7) at (3,2) {7};
\node[tv] (3) at (0.2,1) {3};
\node[tvhl] (4) at (1.8,1) {4};
\node[tvhl] (5) at (1.2,0) {5};
\node[tvhl] (6) at (2.4,0) {6};
\foreach \a/\b in {1/2, 1/7, 2/3, 4/5, 4/6} {\draw[te] (\a) -- (\b);}
\draw[tehl, dashed] (2) -- (4);
\node[tlbl, right] at (4.east) {$\mathit{sz} = 3$};
\node[tlbl, left] at (2.west) {$\mathit{sz} = 5$};
\node[tlbl, right] at (1.east) {$\mathit{sz} = 7$};
\end{tikzpicture}
\diagcaption{以村莊 $1$ 為根。 虛線的道路 $4$--$2$ 下端是 $4$，拿掉後分成 $4$
的子樹（橘色，$3$ 個點）和其他 $4$ 個點。}
\label{fig:tree-roads}
\end{center}

\textbf{想法}

最直接的做法：對每條路，把它拿掉，從一端做一次 BFS 數出那一側有幾個點 $k$，答案
是 $k(n - k)$。 一次 BFS $O(n)$，$n - 1$ 條路就是 $O(n^2)$，$n = 2 \times 10^5$
時是 $4 \times 10^{10}$ 步，來不及。 它可以拿來對拍。

每次都重新數，浪費在同一件事做了 $n - 1$ 次。 照這一節開頭的說法，以村莊 $1$ 為
根，一次 BFS 加上倒著累加，就有每個點的子樹大小。 道路 $(u, v)$ 的兩端，一個是
另一個的父節點；下端的子節點 $c$ 那一側有 $\mathit{sz}(c)$ 個村莊，答案是
$\mathit{sz}(c) \times (n - \mathit{sz}(c))$。

剩下的是：輸入的 $u$、$v$ 哪一個是下端？ 看 \texttt{parent}：如果
\texttt{parent[v]} 是 $u$，下端是 $v$；否則下端是 $u$。 輸入的順序不能信，圖
~\ref{fig:tree-roads} 的範例裡，道路寫成 \texttt{4 2} 的時候，下端是寫在前面的
$4$。

這個「拿掉一條邊，看兩側各有幾個」的想法用途很廣。 每條邊被 $\mathit{sz}(c)
\times (n - \mathit{sz}(c))$ 對點的路線經過，把所有邊的這個數加起來，等於每一對
點的距離加總（本章練習第~\ref{drill:tree-pairsum}~題）。 一條邊如果有長度
$w$，乘上 $w$ 再加，就是帶長度的距離總和。

同樣的模型也出現在別的地方：電力公司評估每段線路壞掉時會切斷多少對用戶之間的
供電；物流公司估計每條路線被多少對倉庫之間的運送用到，決定先拓寬哪一條。

\textbf{程式}

\begin{example}[label=ex:roads]{cpp}{(roads.cpp) 每條道路的通訊量}
#include <iostream>
#include <vector>
using namespace std;

int main() {
    ios::sync_with_stdio(false);
    cin.tie(nullptr);
    int n;
    cin >> n;
    vector<int> eu(n - 1), ev(n - 1);
    vector<vector<int>> adj(n + 1);
    for (int i = 0; i < n - 1; i++) {
        cin >> eu[i] >> ev[i];
        adj[eu[i]].push_back(ev[i]);
        adj[ev[i]].push_back(eu[i]);
    }

    // 以村莊 1 為根：BFS 順序與父節點
    vector<int> parent(n + 1, 0);
    vector<int> order = {1};
    for (int i = 0; i < (int)order.size(); i++) {
        int u = order[i];
        for (int v : adj[u]) {
            if (v != parent[u]) {
                parent[v] = u;
                order.push_back(v);
            }
        }
    }

    // 倒著累加子樹大小
    vector<long long> sz(n + 1, 1);
    for (int i = n - 1; i >= 1; i--) {
        sz[parent[order[i]]] += sz[order[i]];
    }

    for (int i = 0; i < n - 1; i++) {
        // 兩端裡，父節點是另一端的那個是子節點
        int c = (parent[ev[i]] == eu[i]) ? ev[i] : eu[i];
        cout << sz[c] * (n - sz[c]) << (i == n - 2 ? '\n' : ' ');
    }
    return 0;
}
\end{example}

\iosample{
7 \\
2 1 \\
2 3 \\
4 2 \\
4 5 \\
6 4 \\
1 7 \\
}{
10 6 12 6 6 6 \\
}

\begin{itemize}
\item 道路要依輸入順序輸出，所以把兩端存在 \texttt{eu}、\texttt{ev}，等子樹大小
算完再回頭看。
\item 判斷下端只看 \texttt{parent[ev[i]] == eu[i]}，不用反過來再檢查一次：兩端
一定是父子關係，不是 $\mathit{ev}$ 的父節點是 $\mathit{eu}$，就是反過來。
\item \texttt{sz} 直接開成 \texttt{long long}，乘法就用 \texttt{long long} 算。
答案最大在一條鏈的正中間：$10^5 \times 10^5 = 10^{10}$，超過 \texttt{int}。
\end{itemize}

\textbf{複雜度}

一次 BFS、一次倒著累加、一次輸出，都是 $O(n)$。 最大的測試資料在本書測試的環境
下，CPU 時間不到 $0.1$ 秒。

\subsection{常見錯誤}

\begin{itemize}
\item \textbf{把輸入的第二個點當成子節點}：寫成 \texttt{int c = ev[i];}。 範例
的輸入會印出 \texttt{0 6 10 6 12 6}：第一條路寫成 \texttt{2 1}，把根 $1$ 當成
下端，算出 $7 \times 0$。
\item \textbf{乘法用 \texttt{int}}：\texttt{sz} 開成 \texttt{vector<int>}。 輸
入 $n = 2 \times 10^5$ 的一條鏈 $1 - 2 - \cdots - n$，第 $10^5$ 條路的答案是
\texttt{10000000000}，在本書測試的環境下印出 \texttt{1410065408}。 加了
\texttt{-Wall -Wextra} 也沒有警告。
\item \textbf{遞迴算子樹大小}：用上一節的 \texttt{dfs(u, from)} 在回來時累加，
範例和隨機資料都對。 同樣那條 $2 \times 10^5$ 個點的鏈，在本書測試的環境下
\texttt{Segmentation fault}；實測鏈長 $10^5$ 還能跑完，$1.5 \times 10^5$ 就當掉。
這題的測試資料裡就有這樣的鏈。
\end{itemize}

\section{樹直徑}
\label{sec:tree-diameter}

樹上距離最遠的兩個點，它們之間的距離叫\textbf{樹直徑}（tree diameter）；那條路
徑本身也常叫直徑。 在圖~\ref{fig:tree-family} 的族譜上把邊當成無向的，$6$ 到
$3$ 是 $6 - 5 - 2 - 4 - 7 - 3$，$5$ 條邊，沒有更遠的兩點，直徑是 $5$。 直徑不
一定只有一條，$8$ 到 $3$ 也是 $5$。

問法常常不直接說「直徑」：網路上兩台設備之間最多隔幾段線路、步道上最長的單程路
線、要準備多長的管子才拉得到任兩個點。 這些都是在問直徑。

\subsection{兩次走訪}

最直接的做法：每個點各做一次 BFS，找出離它最遠的點，取最大的。 $n$ 次 BFS 是
$O(n^2)$，$n = 2 \times 10^5$ 時來不及。

只要兩次就夠了：

\begin{enumerate}
\item 從任意一點 $s$ 出發走一遍，找出離 $s$ 最遠的點 $u$。
\item 從 $u$ 出發再走一遍，離 $u$ 最遠的點 $v$ 和 $u$ 的距離就是直徑，$u$、$v$
是一條直徑的兩端。
\end{enumerate}

在族譜上試一次，從 $s = 9$ 出發：最遠的是 $6$ 和 $8$，距離都是 $4$（$9 - 4 - 2
- 5 - 6$），取 $u = 6$。 再從 $6$ 出發，最遠的是 $3$，距離 $5$，就是直徑。 如
果第一次取的是 $8$，第二次一樣找到 $3$，距離也是 $5$。

為什麼第一次找到的 $u$ 一定是某條直徑的一端？ 直覺是：$s$ 不管在哪裡，要走到
最遠，總得沿著樹上「最長的那條幹道」往某一頭走到底；走到底的地方就是幹道的一
端。 下面的框把它說清楚。

\begin{whybox}[兩次走訪]
記 $d(p, q)$ 為 $p$、$q$ 的距離。 設 $a$、$b$ 是一條直徑的兩端，$d(a, b) = D$。
從 $s$ 出發，$u$ 是最遠的點。 要說明：有一個點和 $u$ 的距離是 $D$。 這樣第二
次從 $u$ 出發找到的最遠距離就是 $D$。

會用到兩個事實：

\begin{enumerate}
\item[(1)] 兩條路徑 $p \to q$、$q \to r$ 如果只在 $q$ 碰到，接起來就是 $p$ 到
$r$ 的那條唯一的路徑，所以 $d(p, r) = d(p, q) + d(q, r)$。 這裡用到「樹上兩點
的路徑只有一條」。
\item[(2)] 一定有 $d(p, r) \le d(p, q) + d(q, r)$：先到 $q$ 再到 $r$ 是一種走
法，真正的路徑是把其中走過去又走回來的部分刪掉。 這裡用到「邊長不是負的」，刪
掉不會變長。
\end{enumerate}

\begin{center}
\begin{tikzpicture}[x=1cm, y=0.9cm]
\node[tv] (a) at (0,1) {$a$};
\node[tvhl] (x) at (3,1) {$x$};
\node[tv] (b) at (6,1) {$b$};
\node[tv] (s) at (2.4,-0.6) {$s$};
\node[tvalt] (u) at (4.2,2.5) {$u$};
\draw[te] (a) -- (x) -- (b);
\draw[te, dashed] (x) -- (s);
\draw[te, dashed] (x) -- (u);
\end{tikzpicture}
\diagcaption{直徑 $a$--$b$ 畫成水平線。 $x$ 是直徑上離 $s$ 最近的點，$u$ 是離
$s$ 最遠的點。 虛線表示中間可能還有很多點。}
\label{fig:tree-proof}
\end{center}

令 $x$ 是 $a$--$b$ 這條路徑上離 $s$ 最近的點（$s$ 本來就在上面時 $x = s$）。
$s$ 到 $x$ 的路徑只在 $x$ 碰到直徑，由 (1)：
\[
d(s, a) = d(s, x) + d(x, a), \qquad d(s, b) = d(s, x) + d(x, b).
\]
$u$ 是離 $s$ 最遠的，所以 $d(s, u) \ge d(s, b)$；由 (2)，$d(s, u) \le d(s, x) +
d(x, u)$。 兩式合起來得到 $d(x, u) \ge d(x, b)$，同樣也有 $d(x, u) \ge d(x,
a)$。

看 $x$ 到 $u$ 的路徑。 從 $x$ 往 $a$ 和往 $b$ 的路徑，第一條邊不同（它們是直徑
的兩半）；從同一點出發、第一條邊不同的兩條路徑，在樹上只會在出發點碰到，否則就
有環。 所以 $x \to u$ 最多和其中一邊共用第一條邊，假設它和 $x \to a$ 不共用（另
一種情況把 $a$、$b$ 對調）。 那麼 $x \to a$ 和 $x \to u$ 只在 $x$ 碰到，由
(1)：
\[
d(a, u) = d(a, x) + d(x, u) \ge d(a, x) + d(x, b) = D.
\]
$D$ 是所有距離裡最大的，所以 $d(a, u) = D$。 （$x = a$ 或 $u = x$ 時，某條路徑
只有一個點，上面的說法一樣成立。）
\end{whybox}

證明裡用到的兩件事，正好是這個方法的適用範圍：

\begin{itemize}
\item \textbf{一定要是樹。} 下圖六個點圍成一圈，再在 $4$ 掛上 $7$。 從 $1$ 出
發，最遠的是 $6$（距離 $3$）；從 $6$ 出發，最遠的距離還是 $3$。 可是 $7$ 到
$5$ 是 $7 - 4 - 1 - 3 - 5$，距離 $4$。 有環的圖，最遠的兩點要老實從每個點各做
一次 BFS。
\item \textbf{邊長不能是負的。} 負的邊長讓事實 (2) 不成立，本章練習第
~\ref{drill:tree-negative}~題請你造一個反例。 邊長是 $0$ 沒有關係。
\end{itemize}

\begin{center}
\begin{tikzpicture}[x=1.1cm, y=1.0cm]
\node[tvhl] (1) at (0,1) {1};
\node[tv] (3) at (1,2) {3};
\node[tv] (5) at (2.2,2) {5};
\node[tvalt] (6) at (3.2,1) {6};
\node[tv] (2) at (2.2,0) {2};
\node[tv] (4) at (1,0) {4};
\node[tv] (7) at (1,-1.1) {7};
\foreach \a/\b in {1/3, 3/5, 5/6, 6/2, 2/4, 4/1, 4/7} {\draw[te] (\a) -- (\b);}
\end{tikzpicture}
\diagcaption{有環的圖上，兩次走訪會失敗：從 $1$ 找到 $6$，從 $6$ 最遠只有 $3$，
真正最遠的是 $7$ 和 $5$，距離 $4$。}
\label{fig:tree-cycle-cex}
\end{center}

\subsection{邊有長度的樹}

第~11~章說過，邊有不同長度的圖要用續集的 Dijkstra 演算法求最短距離。 樹不用：
兩點之間只有一條路徑，沿著它把長度加起來就是距離。 從起點做一次 BFS，在
\texttt{parent[v] = u;} 的地方寫 \texttt{dist[v] = dist[u] + w;}，每個點的距離
就對了。 這裡的 BFS 只是「一種保證父節點先處理的走訪順序」，不再代表「一層一層
由近到遠」，換成第~11~章的改法一也可以。 所以兩次走訪在帶長度的樹上照樣能用。

\subsection{例題：校園地下管線的最長管路}
\label{sec:ex-pipes}

\textbf{題意}（\akiref{longest-underground-pipe-run-on-campus}{校園地下管線的最長管路}）

\begin{itemize}
\item 校園地下有 $n$ 個閥門，由 $n - 1$ 段水管連成一棵樹，每段水管有長度 $w$。
\item 兩個閥門的距離是路線上水管長度的總和，路線經過幾段水管叫段數。
\item 輸出所有配對裡最大的距離 $D$，以及距離等於 $D$ 的配對裡，段數最多是多少
（記作 $K$）。
\item $n \le 2 \times 10^5$，$1 \le w \le 10^9$。
\end{itemize}

例如下圖 $8$ 個閥門。 從 $3$ 往外有四條管線，末端是 $1$、$8$、$4$、$5$，長度
分別是 $6$、$6$、$6$、$5$，段數是 $2$、$1$、$3$、$1$。 最遠的配對要在長度 $6$
的三個末端裡挑兩個，距離都是 $12$；其中 $1$ 和 $4$ 的段數 $2 + 3 = 5$ 最多，
答案是 \texttt{12 5}。

\begin{center}
\begin{tikzpicture}[x=1.15cm, y=0.95cm]
\node[tv] (3) at (0,0) {3};
\node[tv] (6) at (-1.6,0.9) {6};
\node[tvalt] (1) at (-3.2,1.8) {1};
\node[tv] (8) at (1.8,1.4) {8};
\node[tv] (2) at (-1.2,-1.0) {2};
\node[tv] (7) at (-2.4,-2.0) {7};
\node[tvalt] (4) at (-3.6,-3.0) {4};
\node[tv] (5) at (1.8,-1.2) {5};
\draw[te] (1) -- node[tlbl, above] {2} (6);
\draw[te] (6) -- node[tlbl, above] {4} (3);
\draw[te] (3) -- node[tlbl, above left] {6} (8);
\draw[te] (3) -- node[tlbl, above right] {1} (2);
\draw[te] (2) -- node[tlbl, above right] {2} (7);
\draw[te] (7) -- node[tlbl, above right] {3} (4);
\draw[te] (3) -- node[tlbl, below left] {5} (5);
\end{tikzpicture}
\diagcaption{$8$ 個閥門，邊上的數字是水管長度。 粉色的 $1$、$4$ 是距離 $12$、段
數 $5$ 的配對。}
\label{fig:tree-pipes}
\end{center}

\textbf{想法}

只求 $D$ 的話，就是帶長度的樹直徑，兩次走訪。 麻煩在 $K$：距離 $12$ 的配對有三
組，兩次走訪只找到其中一組，不一定是段數最多的那組。

把「距離」換成一對數（總長，段數），比較時先比總長，總長相同再比段數。 這樣
「最大的距離」就是：總長最大，總長相同的裡面段數最多，正好是題目要的 $(D, K)$。
每段水管的「長度」變成 $(w, 1)$，沿路相加時兩個數各自相加。

換了比較方式，兩次走訪還對嗎？ 回去看證明，只用到兩件事：路徑唯一，以及「多走
一段不會變短」。 前者是樹本身的性質；後者因為每段水管是 $(w, 1)$，$w \ge 1$，
加上去之後總長變大，一定比原來「大」。 比較和加法也配合得起來：兩對數各加上同
一對數，原本誰大，加完還是誰大。 所以證明一字不改都成立，兩次走訪都用數對比較
就好。

\texttt{pair} 本身的比較就是先比 \texttt{first}、相同再比 \texttt{second}，直接
拿 \texttt{pair<long long, int>} 當距離。

同樣的模型也出現在別的地方：要訂一條從任一機房都能拉到任一機房的光纖，長度要
夠，接頭也要備齊；物流中心之間最遠的配送路線，路程相同時看要經過幾個轉運站。

\textbf{程式}

\begin{example}[label=ex:pipes]{cpp}{(pipes.cpp) 校園地下管線的最長管路}
#include <iostream>
#include <utility>
#include <vector>
using namespace std;

using Dist = pair<long long, int>;          // （總長, 段數）

int n;
vector<vector<pair<int, int>>> adj;         // （鄰居, 水管長度）

// 從 src 走遍整棵樹，算出到每個閥門的距離，回傳最遠的閥門
int farthest(int src, vector<Dist> &dist) {
    vector<int> parent(n + 1, 0), order = {src};
    dist.assign(n + 1, Dist(0, 0));
    int best = src;
    for (int i = 0; i < (int)order.size(); i++) {
        int u = order[i];
        if (dist[u] > dist[best]) {         // pair 的比較
            best = u;
        }
        for (auto [v, w] : adj[u]) {
            if (v != parent[u]) {
                parent[v] = u;
                dist[v] = Dist(dist[u].first + w, dist[u].second + 1);
                order.push_back(v);
            }
        }
    }
    return best;
}

int main() {
    ios::sync_with_stdio(false);
    cin.tie(nullptr);
    cin >> n;
    adj.assign(n + 1, {});
    for (int i = 0; i < n - 1; i++) {
        int u, v, w;
        cin >> u >> v >> w;
        adj[u].push_back({v, w});
        adj[v].push_back({u, w});
    }
    vector<Dist> dist;
    int a = farthest(1, dist);              // 第一次：從 1 出發
    int b = farthest(a, dist);              // 第二次：從 a 出發
    cout << dist[b].first << ' ' << dist[b].second << '\n';
    return 0;
}
\end{example}

\iosample{
8 \\
1 6 2 \\
6 3 4 \\
3 8 6 \\
2 3 1 \\
7 2 2 \\
4 7 3 \\
3 5 5 \\
}{
12 5 \\
}

兩次走訪各算出的距離：

\begin{center}
\begin{tabular}{c|cccccccc|c}
\toprule
出發點 & 1 & 2 & 3 & 4 & 5 & 6 & 7 & 8 & 最遠 \\
\midrule
1 & $(0,0)$ & $(7,3)$ & $(6,2)$ & $(12,5)$ & $(11,3)$ & $(2,1)$ & $(9,4)$ & $(12,3)$ & 4 \\
4 & $(12,5)$ & $(5,2)$ & $(6,3)$ & $(0,0)$ & $(11,4)$ & $(10,4)$ & $(3,1)$ & $(12,4)$ & 1 \\
\bottomrule
\end{tabular}
\end{center}

第一次從 $1$ 出發，$4$ 和 $8$ 的總長都是 $12$，段數 $5$ 比較多，選 $4$。 第二次
從 $4$ 出發，最遠的是 $(12, 5)$ 的 $1$。

\begin{itemize}
\item \texttt{farthest} 一邊走訪一邊更新 \texttt{best}。 每個點在被處理時，它的
距離已經在父節點那裡算好了，所以可以直接比。
\item \texttt{for (auto [v, w] : adj[u])} 是結構化綁定（銜接書第~8~章），
\texttt{v}、\texttt{w} 是串列裡每個 \texttt{pair} 的兩個成員。
\item 第二次呼叫時 \texttt{dist} 會被 \texttt{assign} 整個重設，裡面留下的是從
$a$ 出發的距離。
\item 總長最大約 $2 \times 10^5 \times 10^9 = 2 \times 10^{14}$，要用
\texttt{long long}。 段數最多 $n - 1$，\texttt{int} 就夠。
\end{itemize}

\textbf{複雜度}

兩次走訪各 $O(n)$。 最大的測試資料在本書測試的環境下，CPU 時間不到 $0.1$ 秒。

\subsection{常見錯誤}

\begin{itemize}
\item \textbf{只比總長，最後印出那一組的段數}：把 \texttt{dist[u] > dist[best]}
改成只比 \texttt{.first}。 範例的輸入會印出 \texttt{12 3}：第一次從 $1$ 出發，
$8$ 比 $4$ 先被處理，兩者總長相同就留下 $8$；第二次從 $8$ 出發，又先遇到總長
$12$、只有 $3$ 段的 $1$。 走訪順序剛好對的時候會過，所以範例過了也不代表對，
這題的測試資料有一半以上會錯。
\item \textbf{只在第二次比段數}：第一次只比總長，找到的端點可能不在段數最多的
那一組裡（範例就是 $8$），第二次怎麼比都救不回來，範例的輸入會印出
\texttt{12 4}。 兩次都要用同一種比較。
\item \textbf{總長用 \texttt{int}}：把 \texttt{Dist} 改成
\texttt{pair<int, int>}，輸入一條 $2 \times 10^5$ 個閥門、每段 $10^9$ 的鏈，正確
答案是 \texttt{199999000000000 199999}，在本書測試的環境下印出
\texttt{2147459072 174992}，連段數都錯了：溢位之後比較大小的結果也不可靠。
\end{itemize}

\begin{apcsnote}
「兩次走訪」的證明不需要記，但要記得它的條件：樹、邊長不是負的。 題目在這兩點
上有變化時（多了一條邊、某些點不能當端點、距離的定義不同），先用「每個點各做一
次 BFS」的暴力解在小資料上對拍，確定兩次走訪還成立再交。
\end{apcsnote}

\section{二元樹的三種走訪順序}
\label{sec:tree-binary}

每個點最多兩個子節點，而且分得出左、右的有根樹，叫\textbf{二元樹}（binary
tree）。 只有一個子節點時，它是左子節點還是右子節點也算在形狀裡。 存法用兩個
陣列：\texttt{L[v]}、\texttt{R[v]} 是 $v$ 的左、右子節點，沒有就填 $0$。 根是
沒有出現在任何 \texttt{L}、\texttt{R} 裡的那個點。

遞迴走一棵二元樹時，「處理自己」可以放在三個時間點，得到三種順序：

\begin{itemize}
\item \textbf{前序}（preorder）：自己、左子樹、右子樹。
\item \textbf{中序}（inorder）：左子樹、自己、右子樹。
\item \textbf{後序}（postorder）：左子樹、右子樹、自己。
\end{itemize}

第~8~章的呼叫樹印的就是中序：先走左邊，回來印自己，再走右邊。 下面的程式只走
一次，在三個時間點各記錄一次，同時得到三種順序：

\begin{example}[label=ex:traversal]{cpp}{(traversal.cpp) 二元樹的前序、中序、後序}
#include <iostream>
#include <vector>
using namespace std;

vector<int> L, R, pre, in, post;

void walk(int v) {
    if (v == 0) {
        return;
    }
    pre.push_back(v);       // 第一次到 v
    walk(L[v]);
    in.push_back(v);        // 左子樹走完，回到 v
    walk(R[v]);
    post.push_back(v);      // 兩邊都走完，要離開 v
}

void print(const char *name, const vector<int> &a) {
    cout << name;
    for (int v : a) {
        cout << ' ' << v;
    }
    cout << '\n';
}

int main() {
    int n;
    cin >> n;
    L.assign(n + 1, 0);
    R.assign(n + 1, 0);
    vector<bool> isChild(n + 1, false);
    for (int v = 1; v <= n; v++) {
        cin >> L[v] >> R[v];
        isChild[L[v]] = isChild[R[v]] = true;   // 0 也會被標，不影響
    }
    int root = 1;
    while (isChild[root]) {
        root++;
    }
    walk(root);
    print("pre:", pre);
    print("in:", in);
    print("post:", post);
    return 0;
}
\end{example}

\iosample{
7 \\
0 0 \\
7 0 \\
0 0 \\
0 0 \\
1 3 \\
2 5 \\
0 4 \\
}{
pre: 6 2 7 4 5 1 3 \\
in: 7 4 2 6 1 5 3 \\
post: 4 7 2 1 3 5 6 \\
}

\begin{center}
\begin{tikzpicture}[x=1.1cm, y=1.0cm]
\node[tvhl] (6) at (2,3) {6};
\node[tv] (2) at (0.8,2) {2};
\node[tv] (5) at (3.2,2) {5};
\node[tv] (7) at (0.2,1) {7};
\node[tv] (1) at (2.6,1) {1};
\node[tv] (3) at (3.8,1) {3};
\node[tv] (4) at (0.8,0) {4};
\foreach \a/\b in {6/2, 6/5, 2/7, 5/1, 5/3, 7/4} {\draw[te] (\a) -- (\b);}
\end{tikzpicture}
\diagcaption{範例的二元樹。 $2$ 只有左子節點 $7$，$7$ 只有右子節點 $4$。}
\label{fig:tree-binary}
\end{center}

三種順序的規律：

\begin{itemize}
\item 前序的第一個、後序的最後一個都是根。
\item 中序裡，根把序列切成兩半：左邊是左子樹的點，右邊是右子樹的點。 範例的中
序 \texttt{7 4 2 | 6 | 1 5 3} 就是這樣。
\item 每一棵子樹在三種順序裡都是連在一起的一段。
\end{itemize}

這幾條合起來，給前序和中序就能一層一層切出整棵樹的形狀（本章練習第
~\ref{drill:tree-binary}~題）；給後序和中序也可以。 只給前序和後序就不一定了，
同一組序列可能對應好幾棵樹。

這支程式用遞迴，深度等於樹的高度。 二元樹也可能是一條很長的鏈，$n$ 到 $10^5$
以上時，要照第~11~章的改法二改寫：堆疊裡記住每個點走到哪一步（還沒走左邊、左
邊走完、兩邊都走完）。

\section{預告：樹形動態規劃}
\label{sec:tree-dp-preview}

這一章由下往上算的量，高度、子樹大小、子樹和，都是同一個形狀：每個點的答案，用
它所有子節點的答案組出來。 把「每個點的答案」當成狀態、「怎麼用子節點的答案組
出來」當成轉移，這就是第~9~章的動態規劃搬到樹上，叫\textbf{樹形動態規劃}（tree
DP）。 順序也已經有了：倒著看 BFS 順序。

到目前為止，每個點只存一個數。 第~20~章的題目，每個點常常要存好幾個數，例如：

\begin{itemize}
\item 在樹上挑一些點，相鄰的兩點不能都挑，要讓挑到的值總和最大。 每個點要存兩
個數：挑自己時子樹裡的最大總和、不挑自己時的最大總和。
\item 直徑也可以用樹形 DP 算：最長的路徑一定有一個「最高點」，在那個點把往下最
長的兩條鏈接起來。 每個點存「往下最長的鏈」，順便更新答案。
\end{itemize}

這些留到第~20~章，用的仍然是這一章的迴圈，只是每個點多存幾個數。

\section{本章常見錯誤}
\label{sec:tree-pitfalls}

\begin{center}
\begin{tabular}{p{0.22\linewidth}p{0.46\linewidth}p{0.2\linewidth}}
\toprule
錯誤 & 例子 & 見 \\
\midrule
假設根是 $1$，或父節點的編號比較小 & 照編號順序由下往上算；直接用 $1$ 當族譜的根 &
第~\ref{sec:tree-terms}~節、第~6~章 \\
走回父節點 & 無向的樹走訪時沒記「從哪裡來」；用 \texttt{parent[v] == 0} 判斷沒走
過，把根當成子節點 & 第~\ref{sec:ex-phonetree}~節 \\
照輸入的方向當父子 & 把邊的第二個點當成子節點 & 第~\ref{sec:ex-phonetree}、
\ref{sec:ex-roads}~節 \\
遞迴太深 & 一條 $2 \times 10^5$ 個點的鏈，遞迴算子樹大小當掉 &
第~\ref{sec:ex-roads}~節 \\
\texttt{int} 溢位 & $\mathit{sz} \times (n - \mathit{sz})$、子樹和、帶長度的距
離 & 第~\ref{sec:ex-roads}、\ref{sec:ex-pipes}~節 \\
兩次走訪用錯地方 & 有環的圖、負的邊長；比較的標準兩次不一致 &
第~\ref{sec:tree-diameter}~節 \\
\bottomrule
\end{tabular}
\end{center}

另外兩個邊界：只有一個點的樹（沒有邊，\texttt{n - 1} 是 $0$，很多迴圈一次都不
跑），以及深度、高度從 $0$ 還是 $1$ 數。 交出去之前各用一組小輸入確認。

\section{本章練習}
\label{sec:ch13-exercises}

\begin{drills}[label=drills:ch13]
\item 一棵樹有 $8$ 個點，父節點陣列是 \texttt{parent} $= [3, 0, 2, 3, 2, 5, 5,
4]$（依序是 $1$ 到 $8$ 號，$0$ 表示根）。
\begin{enumerate}
\item 找出根，畫出這棵樹，寫出每個點的深度、高度和子樹大小。
\item 子節點串列照編號由小到大，寫出從根出發的 BFS 順序。
\item 下面三種順序，哪些倒過來看（從最後一個往前）可以拿來由下往上算子樹大小？
說明理由。 (i) $1, 2, \ldots, 8$；(ii) $2, 5, 7, 6, 3, 4, 8, 1$；(iii) 依深度由
小到大，深度相同時編號小的在前。
\end{enumerate}

\item \label{drill:tree-binary}
\begin{enumerate}
\item 一棵二元樹：$1$ 的左、右子節點是 $2$、$3$；$2$ 只有左子節點 $4$；$4$ 只
有右子節點 $5$；$3$ 的左、右子節點是 $6$、$7$。 寫出前序、中序、後序。
\item 一棵二元樹的前序是 $5, 2, 1, 4, 3, 6$，中序是 $1, 2, 3, 4, 5, 6$。 畫出這
棵樹。
\item 前序是 $1, 2, 3$、後序是 $3, 2, 1$ 的二元樹有幾棵？ 全部畫出來。
\end{enumerate}

\item \label{drill:tree-pairsum} 證明：把樹上每條邊的 $\mathit{sz}(c) \times (n
- \mathit{sz}(c))$ 加起來，等於所有點對 $\{p, q\}$ 的距離總和。 用範例
~\ref{ex:roads} 的輸出驗算一次。

\item 有人把範例~\ref{ex:roads} 改成「每條路拿掉後從一端 BFS 數一側的點數」。
\begin{enumerate}
\item 時間複雜度是多少？ $n = 2 \times 10^5$ 時大約要走幾步？
\item 範例~\ref{ex:tree-basic} 由下往上那個迴圈，\texttt{sz[p] += sz[v];} 一共
執行幾次？ 為什麼不是 $O(n^2)$，雖然一個點可能有很多子節點？
\end{enumerate}

\item \label{drill:tree-ecc} 設 $a$、$b$ 是一條直徑的兩端。 證明：對樹上任意一
點 $s$，離 $s$ 最遠的距離等於 $\max(d(s, a), d(s, b))$。 （提示：沿用證明框的
$x$，說明 $d(x, u) \le \max(d(x, a), d(x, b))$。） 這個性質讓「每個點離最遠的
點多遠」只要三次走訪：先兩次找到 $a$、$b$，再從 $b$ 走一次。

\item \label{drill:tree-negative} 造一棵有負邊長的樹（四個點以內），讓兩次走訪
得到的答案比真正最遠的距離小。 說明證明框裡哪一步不成立。

\item 第~6~章的「層層彙整的報告」，改成用本章的方法：先從總經理做 BFS，再倒著
看 BFS 順序算每個人完成報告的時刻。 輸入只有每個人的直屬主管，程式要先多建什
麼？ 寫出倒著看的那個迴圈。
\end{drills}

\section{實作練習}
\label{sec:ch13-practice}

下面是 Aki-Judge 上和本章相關的題目，依小節分組。 難度分四階：暖身是直接套用
本章的寫法；基本要自己決定根、決定由上往下還是由下往上；實戰的題敘比較長，要
先把題目轉換成本章的形狀；挑戰會和其他章的技巧合在一起。

\begin{center}
\begin{tabular}{llc}
\toprule
小節 & 題目 & 難度 \\
\midrule
表示與定根 & \akiref{full-path-of-a-library-category}{圖書分類的完整路徑} & 暖身 \\
 & \akiref{which-way-the-water-flows}{灌溉渠道的水流方向} & 暖身 \\
 & \akiref{checking-reported-generations}{社團傳承的自報代數} & 暖身 \\
 & \akiref{the-club-mentor-register}{學長姐傳承名冊} & 基本 \\
 & \akiref{indented-meeting-notes}{會議紀錄的縮排大綱} & 實戰 \\
\midrule
高度與子樹 & \akiref{days-until-the-memo-reaches-the-bottom}{公告傳到最底層要幾天} & 暖身 \\
 & \akiref{overloaded-delivery-roads}{超過載量的運輸路段} & 暖身 \\
 & \akiref{who-sits-on-the-longest-relay-chains}{最長轉達鏈上的人} & 基本 \\
 & \akiref{top-score-in-each-group}{帶領範圍內的最高分} & 基本 \\
 & \akiref{interchangeable-teams}{可以互換的小組} & 基本 \\
 & \akiref{zone-meters}{分區電表的裝設位置} & 基本 \\
 & \akiref{total-reporting-steps}{回報路徑總步數} & 實戰 \\
 & \akiref{removing-one-manager}{拿掉一層主管} & 實戰 \\
\midrule
邊的貢獻 & \akiref{stranded-rooms}{封走廊後受困的房間數} & 暖身 \\
 & \akiref{the-old-system-only-stores-binary-trees}{只能存二元樹的舊系統} & 挑戰 \\
\midrule
樹直徑 & \akiref{the-two-farthest-bike-stations}{相距最遠的兩個租借站} & 基本 \\
 & \akiref{longest-line-after-the-outage}{停電後各區最長的線路} & 基本 \\
 & \akiref{tour-route-with-the-most-junction-stops}{解說站最多的導覽路線} & 基本 \\
 & \akiref{the-longest-one-way-hike}{最長單程健行路線的起訖站} & 實戰 \\
 & \akiref{night-inspection-starting-points}{巡檢員的出發地點} & 實戰 \\
 & \akiref{temporary-link-proposals-between-two-campuses}{兩校區網路的臨時連線提案} & 挑戰 \\
\midrule
二元樹 & \akiref{when-to-sign}{簽核的時機} & 暖身 \\
 & \akiref{a-family-tree-in-numbered-slots}{書架格號族譜} & 暖身 \\
 & \akiref{the-inspection-robots-move-log}{巡檢機器人的移動紀錄} & 基本 \\
\bottomrule
\end{tabular}
\end{center}

本章的三題例題也在 Aki-Judge 上，讀完各節之後可以先自己寫一次再對照。 直徑那
一組有幾題會用到練習第~\ref{drill:tree-ecc}~題的性質。

\section{讀完你應該會}
\label{sec:ch13-checklist}

\begin{itemize}
\item 拿到父節點陣列、子節點名單或無向邊三種輸入，都能存成程式要用的形式，找出
根，並且不假設根是 $1$、父節點編號比較小。
\item 能從根做一次 BFS，用「不走回父節點」定出每條邊的方向，順著 BFS 順序算深
度，倒著算高度、子樹大小和子樹和，不用遞迴。
\item 看到「每條邊兩側」「經過這條邊的配對」，能用 $\mathit{sz}(c) \times (n -
\mathit{sz}(c))$ 一次算出全部的邊，並分得出輸入的兩端哪一個是子節點。
\item 會用兩次走訪求樹直徑（含帶長度的樹），說得出它需要「是樹」和「邊長不是負
的」，並能在比較標準改變時判斷證明是否還成立。
\item 能寫出一棵二元樹的前序、中序、後序，給前序和中序能畫出原來的樹。
\end{itemize}

\textbf{延伸閱讀}（書末會整理完整的參考資料）

\begin{itemize}
\item Antti Laaksonen，\textit{Competitive Programmer's Handbook}，第 14 章：
篇幅很短，有走訪、子樹大小、直徑的兩種算法，以及「每個點最遠多遠」的做法，可以
當第~20~章的預習。
\item 吳邦一，《AP325：從 APCS 實作題檢測三級到五級》，第 8 章：樹的各種儲存方
式和由下往上的走訪，題目很多，適合要準備 APCS 的讀者多練。
\item Thomas H. Cormen 等，\textit{Introduction to Algorithms}，第 4 版，10.3
節：有根樹的幾種表示法；12.1 節：二元搜尋樹與中序走訪。
\end{itemize}

\addcontentsline{toc}{section}{練習解答}
\section*{練習解答}

\begin{enumerate}
\item
\begin{enumerate}
\item 根是 $2$。 子節點：$2$ 有 $3$、$5$；$3$ 有 $1$、$4$；$4$ 有 $8$；$5$ 有
$6$、$7$。
\begin{center}
\begin{tabular}{c|cccccccc}
\toprule
點 & 1 & 2 & 3 & 4 & 5 & 6 & 7 & 8 \\
\midrule
深度 & 2 & 0 & 1 & 2 & 1 & 2 & 2 & 3 \\
高度 & 0 & 3 & 2 & 1 & 1 & 0 & 0 & 0 \\
子樹大小 & 1 & 8 & 4 & 2 & 3 & 1 & 1 & 1 \\
\bottomrule
\end{tabular}
\end{center}
\item $2, 3, 5, 1, 4, 6, 7, 8$。
\item 要求的是：倒過來看時，每個點都在它所有子節點之後，也就是原本的順序裡，每
個點都在它的子節點之前。 (i) 不行：$3$ 是 $1$ 的父節點，卻排在 $1$ 後面；倒過來
看時先處理 $3$，把還沒加上 $1$ 的 \texttt{sz[3]} 交給 $2$，\texttt{sz[2]} 會少
$1$。 (ii) 可以，逐一檢查父節點都在前面（它是第~11~章改法一那種堆疊走訪的順
序）。 (iii) 可以：父節點的深度比子節點小一，一定排在前面。 BFS 順序就是 (iii)
的一種。
\end{enumerate}

\item
\begin{enumerate}
\item 前序 $1, 2, 4, 5, 3, 6, 7$；中序 $4, 5, 2, 1, 6, 3, 7$；後序 $5, 4, 2, 6,
7, 3, 1$。 把這棵樹寫成範例~\ref{ex:traversal} 的輸入跑一次可以核對。
\item 前序的第一個 $5$ 是根；中序裡 $5$ 的左邊 $1, 2, 3, 4$ 是左子樹，右邊 $6$
是右子樹。 左子樹的前序是接下來的四個 $2, 1, 4, 3$，根是 $2$；中序 $1 | 2 | 3,
4$，$1$ 是 $2$ 的左子節點，右子樹是 $\{3, 4\}$，前序 $4, 3$，根是 $4$，中序 $3,
4$ 表示 $3$ 是 $4$ 的左子節點。 所以：$5$ 的左 $2$、右 $6$；$2$ 的左 $1$、右
$4$；$4$ 的左 $3$。
\item $4$ 棵。 前序和後序都只說明 $1$ 在最上面、$2$ 在中間、$3$ 在最下面，排成
一條鏈；但 $2$ 可以是 $1$ 的左或右子節點，$3$ 可以是 $2$ 的左或右子節點，前序、
後序都不變，$2 \times 2 = 4$ 種。 中序就分得出來：四棵的中序各不相同。
\end{enumerate}

\item 一對點 $\{p, q\}$ 的距離，是它們之間的路徑有幾條邊。 換個方式數：對每條
邊，問它在幾對點的路徑上，再全部加起來，數到的是同一個總數（每一對點，路徑上
的每條邊各被數到一次）。 一條邊在 $\{p, q\}$ 的路徑上，剛好是拿掉這條邊後 $p$、
$q$ 分在兩邊的時候，這樣的配對有 $\mathit{sz}(c) \times (n - \mathit{sz}(c))$
對。 範例的輸出加起來是 $10 + 6 + 12 + 6 + 6 + 6 = 46$。 直接算：以 $1$ 為根
時各點深度是 $1{:}0$、$2{:}1$、$7{:}1$、$3{:}2$、$4{:}2$、$5{:}3$、$6{:}3$，
$21$ 對點的距離逐一加起來也是 $46$。

\item
\begin{enumerate}
\item 每條路一次 BFS，$O(n)$，共 $n - 1$ 條，$O(n^2)$。 $n = 2 \times 10^5$ 時
約 $4 \times 10^{10}$ 步，遠超過時間限制。
\item 迴圈 \texttt{i} 從 $n - 1$ 到 $1$，剛好 $n - 1$ 次。 每次是一個點把值交
給父節點，每個非根的點只交一次；子節點多的點只是「收到」很多次，收的次數加起來
也是 $n - 1$。 這和「每條邊只看一次」是同一件事。
\end{enumerate}

\item 沿用證明框的記號，$u$ 是離 $s$ 最遠的點。 $x \to u$ 最多和 $x \to a$、$x
\to b$ 其中一邊共用第一條邊，設它不和 $x \to a$ 共用。 如果 $d(x, u) > d(x, b)$，
由事實 (1)，$d(a, u) = d(a, x) + d(x, u) > d(a, x) + d(x, b) = D$，和 $D$ 最大
矛盾，所以 $d(x, u) \le d(x, b)$。 由事實 (2)，$d(s, u) \le d(s, x) + d(x, u)
\le d(s, x) + d(x, b) = d(s, b)$。 另一方面 $d(s, u)$ 本來就是最大的，不小於
$d(s, a)$、$d(s, b)$。 所以離 $s$ 最遠的距離等於 $\max(d(s, a), d(s, b))$。

\item 一條鏈 $1 - 2 - 3 - 4$，邊長依序是 $3$、$-10$、$8$。 從 $2$ 出發：到 $1$
是 $3$、到 $3$ 是 $-10$、到 $4$ 是 $-2$，最遠的是 $1$。 從 $1$ 出發：到 $2$ 是
$3$、到 $3$ 是 $-7$、到 $4$ 是 $1$，最遠只有 $3$。 但 $3$ 到 $4$ 的距離是 $8$。
不成立的是事實 (2)：從 $1$ 到 $4$ 的真正路徑經過長度 $-10$ 的邊，「刪掉走回頭
的部分不會變長」在這裡變成會變長，證明裡 $d(s, u) \le d(s, x) + d(x, u)$ 那一步
就斷了。

\item 輸入只有每個人的主管 \texttt{boss}，要從總經理往下 BFS，得先建出每個人的
部屬串列 \texttt{children}（和第~\ref{sec:tree-terms}~節由父節點陣列建子節點串
列一樣），並找出 \texttt{boss} 是 $0$ 的總經理。 BFS 得到 \texttt{order} 之後：
\begin{code}{cpp}{練習 7}
// ready[v]：部屬裡最晚的完成時刻，一開始都是 0
for (int i = n - 1; i >= 0; i--) {
    int v = order[i];
    finish[v] = ready[v] + t[v];
    if (boss[v] != 0) {
        ready[boss[v]] = max(ready[boss[v]], finish[v]);
    }
}
\end{code}
這次迴圈要跑到 \texttt{i = 0}，因為總經理自己的 \texttt{finish} 也要算。 輪到
$v$ 時，它的部屬都在 \texttt{order} 的後面，已經處理過了，\texttt{ready[v]} 是
最終的值。
\end{enumerate}

% 佔位：本章尚未開始。book.tex 裡先註解掉。
\chapter{綜合練習}

本章尚未開始。


% 原本的第 15–20 章（背包 DP、子序列 DP、區間 DP、DAG 與拓樸排序、並查集、樹形 DP）
% 移到《C++ 演算法解題 2：動態規劃與圖論》（repo cpp-algorithms-book-2）。

% 書末（不編章號）
% 書末：參考資料與延伸閱讀（不編章號）
\chapter*{參考資料與延伸閱讀}
\addcontentsline{toc}{chapter}{參考資料與延伸閱讀}
\markboth{參考資料與延伸閱讀}{參考資料與延伸閱讀}
\label{ch:references}

各章最後的「延伸閱讀」提到的書，完整的書目資料都列在這裡，每本附一兩句它適合
做什麼，以及和本書哪幾章搭配。 括號裡的章節號碼是那本書的章節，以這裡列出的
版本為準；換了版本，章節號碼可能不同。

本書的講解、例子和程式都是自己寫的，參考這些書的地方是題型的分類、講解的先後
和經典的結論。 想看另一種講法、更嚴謹的證明，或想多做一些題目時，可以從下面
挑一本。

\section*{篇幅短、適合快速複習}

\begin{itemize}
\item Antti Laaksonen，\textit{Competitive Programmer's Handbook}，2018 年 7 月
3 日草稿版，作者公開的 PDF。

每個主題只用幾頁、一兩個小例子講完，適合讀完本書某一章後，花十幾分鐘從另一個
角度再看一次，也適合想快速看過全貌的人。 本書第~1~到第~13~章的延伸閱讀幾乎
都引用了它：第~1~章（第 1、2 章），第~2~章（第 3 章），第~3~章（3.3 節），
第~4~章（第 5 章），第~5~章（2.4、8.1、9.1 節），第~6~章（4.5、8.2、8.3 節），
第~7~章（第 6 章），第~8~章（3.1 節），第~9~章（第 7 章），第~10~章（7.3 節），
第~11~章（第 11、12 章），第~12~章（第 12 章），第~13~章（第 14 章）。
\end{itemize}

\section*{題型分類細、題目多}

\begin{itemize}
\item 吳邦一，《AP325：從 APCS 實作題檢測三級到五級》，v1.5，2024 年 9 月，作者
公開的講義。

以 APCS 實作題的程度為目標，每個主題附很多例題和習題，還有測試資料，適合要準
備 APCS 的讀者多練。 它的第 0 章（預備知識）到第 8 章（樹上演算法）和本書第~1~%
到第~13~章的範圍大致平行：遞迴、排序與二分搜、佇列與堆疊、貪心與掃描線、分治、
動態規劃、圖、樹都有。 各章引用的是第~13~章（第 8 章）和第~14~章（第 0 章）。

\item Steven Halim、Felix Halim、Suhendry Effendy，\textit{Competitive
Programming 4}，Book 1 與 Book 2，2020（作者自行出版）。

把題目依技巧分得很細，每一節後面附一長串分好類的練習題，適合打算參加程式競
賽、想大量練習的讀者。 第 1 章教怎麼從題目的限制判斷題型、估計複雜度，和本書
第~1、14~章的想法相同。 和本書搭配的地方：第~1~章（Book 1 的 1.3 節），
第~3~章（Book 1 的 3.3.1 節），第~5~章（Book 1 的 3.5.2 節），第~10~章（Book 1
的 3.5.2 節），第~12~章（Book 1 的 4.2.5 節，Book 2 的 8.2.2 節），第~14~章
（Book 1 的 1.3 節）。

\item 劉汝佳，《算法競賽入門經典》，第 2 版，清華大學出版社，2014。

簡體中文。 例題多，程式寫得精簡，適合想用中文多看幾種題型的讀者。 和本書搭配
的地方：第~4~章（第 7 章），第~6~章（6.1 節），第~9~章（9.1、9.2 節），第~10~%
章（9.1 節），第~12~章（6.4、7.5 節）。
\end{itemize}

\section*{嚴謹的證明與演算法設計}

\begin{itemize}
\item Thomas H. Cormen、Charles E. Leiserson、Ronald L. Rivest、Clifford Stein，
\textit{Introduction to Algorithms}，第 4 版，The MIT Press，2022。

大學演算法課常用的教科書，每個演算法都有完整的正確性證明和分析。 想知道本書
省略的證明細節，或要修演算法課的讀者，可以對照著讀。 和本書搭配的地方：第~2~章
（2.1、2.3 節，第 8 章），第~5~章（16.1 節），第~6~章（第 6 章，16.1 節），第~7~%
章（15.1、15.2 節），第~8~章（2.3、4.4 節），第~9~章（14.1、14.3 節），第~10~%
章（14.4 節），第~11~章（20.1 到 20.3 節），第~13~章（10.3、12.1 節）。

\item Steven S. Skiena，\textit{The Algorithm Design Manual}，第 3 版，Springer，
2020。

重點放在怎麼想出演算法、怎麼替一條看似合理的規則找反例，以及怎麼把實際的問題
對應到已知的問題，後半是一份常見問題的目錄。 適合想理解「做法是怎麼想出來的」
的大學生和自學的人。 和本書搭配的地方：第~1~章（第 2 章），第~2~章（4.1 節），
第~3~章（5.1 節），第~7~章（1.2、1.3 節），第~8~章（1.4 節，第 5 章），第~11~章
（第 7 章），第~14~章（第 13 章）。
\end{itemize}

\section*{本系列的其他書}

\begin{itemize}
\item 《C++ 程式設計：給學過 C 的人》：各章提到的「銜接書」，本書用到的 C++ 語法
和 STL 都在那裡。
\item 《C++ 演算法解題 2：動態規劃與圖論》：本書的下一本。 背包、子序列、區間、
位元、數位等動態規劃，拓樸排序、並查集、加權最短路徑、最小生成樹和樹形 DP 都
在那裡。
\end{itemize}


\end{document}
